% PM 트랙 Day 1 슬라이드 본문 — 손으로 쓴 정본(v2, 한 프레임 = 한 쪽).
% 래퍼: Day1_슬라이드.tex(슬라이드판) / Day1_슬라이드_노트.tex(노트판, 우측에 \note).
% 화면에는 결론만, 설명·근거·예외는 각 프레임의 \note{}에. Day1_슬라이드.md 는 이 파일의 텍스트 미러다.
% 이전 자동 생성본(172쪽)은 Day1_슬라이드_body_v1_172p.tex 에 보관.

\newcolumntype{L}[1]{>{\raggedright\arraybackslash}p{#1}}
\newcommand{\ra}{$\rightarrow$}
\newcommand{\til}{\textasciitilde}
\newcommand{\keymsg}[1]{\par\vfill{\color{mDarkTeal}\bfseries #1}\par}
\newcommand{\figcap}[1]{\par\smallskip{\footnotesize\color{black!60}#1}\par}
\newcommand{\fig}[2][0.66]{\centering\includegraphics[height=#1\textheight,width=\textwidth,keepaspectratio]{#2}\par}
\newcommand{\subhead}[1]{{\color{mDarkTeal}\bfseries #1}\par}
\newenvironment{tbl}[2][\small]{\begin{center}\usebeamercolor[fg]{normal text}\setlength{\tabcolsep}{4pt}\renewcommand{\arraystretch}{1.15}#1\begin{tabular}{#2}\toprule}{\bottomrule\end{tabular}\end{center}}

% =====================================================================
% 도입
% =====================================================================
\begin{frame}{PM 트랙 Day 1 — 오늘의 사례와 세 질문}
\begin{itemize}
\item 프로젝트 매니지먼트 4일 특강 · 1일차
\item 사례 회사: ㈜온담F\&B홀딩스 ONE ONDAM 프로그램 \textbf{P1} 「차세대 주문·재고 통합 플랫폼 구축」
\item 승인 예산 \textbf{168억 원} · 2026-01-05 \til{} 2027-05-29 · 하이브리드(스프린트 36회 + 게이트 G0\til{}G5)
\end{itemize}
\keymsg{오늘의 세 질문 — 이것은 무엇인가 / 무엇을 따를 것인가 / 어떻게 시작할 것인가}
\note{표지. 사례 회사 문서(교재/공통/가상회사\_온담F앤B.pdf)는 §1.1·§3.4·§4.2·§6만 사전 읽기로 안내했다(약 25분). 읽어 왔는지 묻되, 안 읽어 온 학생도 0교시 사례 브리핑 20분으로 따라올 수 있게 해 확인한다. 읽지 않은 수강생은 워크북 §0.3 한 페이지 요약을 4교시 전까지 읽도록 안내한다. 오늘 하루 모든 예시는 이 P1 하나로 통일되며, 세 질문(정의 / 표준 / 착수)이 오전·오후의 순서다. 교재 서문 참조. 예상 소요 3분.}
\end{frame}

\begin{frame}{오늘의 흐름}
\begin{columns}[T]
\begin{column}{0.46\textwidth}
\subhead{오전 — 이론 (제1\til{}2장)}
\begin{itemize}
\item \textbf{0교시} 20분 · 사례 브리핑 — 온담F\&B와 P1, 책받침 한 장
\item \textbf{1교시} 80분 · 용어 정리 · 프로젝트란 무엇인가 · 왜 실패하는가 — 제1장
\item \textbf{2교시} 90분 · 표준의 지형 — PMBOK 6\ra{}7\ra{}8, PRINCE2, ISO, 테일러링 — 제2장
\end{itemize}
\end{column}
\begin{column}{0.52\textwidth}
\subhead{오후 — 착수 이론과 실습이 번갈아}
\begin{itemize}
\item \textbf{3교시} 90분 · 착수 이론 — BC · 편익 소유 · 헌장 · 이해관계자 — 제3장 전반
\item \textbf{4교시} 90분 · \textbf{실습 ①} BC 요약 + 헌장 — 워크북 §1\til{}2
\item \textbf{5교시} 75분 · 거버넌스 · 권한 없는 PM · 프리모템 · 분야별 착수 — 제3장 후반
\item \textbf{6교시} 75분 · \textbf{실습 ②} 등록부 · RACI · 프리모템 — 워크북 §3\til{}5
\item \textbf{마무리} 20분 · 학술 배경 · 읽을 자료 · Day2 예고 — 제4\til{}5장
\end{itemize}
\end{column}
\end{columns}
\keymsg{오늘 만든 다섯 산출물이 Day2 계획 산출물의 입력이 된다}
\note{오전은 이론(제1\til{}2장), 오후는 착수 이론과 실습이 번갈아 온다. 실습 ①(4교시)은 비즈니스 케이스 요약과 헌장, 실습 ②(6교시)는 등록부·RACI·프리모템이다. 점심은 2교시 뒤. 오늘 만든 다섯 산출물이 Day2 계획 산출물(WBS·RTM·일정·원가 기준선·리스크 정량화)의 입력이 된다는 점을 미리 말해 둔다. 예상 소요 3분.}
\end{frame}

\begin{frame}{학습 목표 — 오늘이 끝나면}
\begin{columns}[T]
\begin{column}{0.5\textwidth}
\subhead{오전 이론}
\begin{enumerate}
\item 프로젝트·프로그램·운영을 \textbf{돈과 권한의 언어}로 구분한다
\item CHAOS·Flyvbjerg 수치를 \textbf{왜 그대로 인용하면 안 되는지} 말한다
\item PMBOK 6\ra{}7\ra{}8 계보 — 사내 표준을 \textbf{재매핑 대장} 네 칸으로
\item 테일러링 기준 = \textbf{되돌림 비용} — P1과 P2에 다른 처방
\end{enumerate}
\end{column}
\begin{column}{0.48\textwidth}
\subhead{오후 실습}
\begin{enumerate}\setcounter{enumi}{4}
\item BC 결함을 목록화하고 \textbf{편익 오너가 PM이 아닌 이유}를 설명한다
\item 헌장 15항목 · 이해관계자 3속성 7유형 · RACI · 프리모템을 온담 P1으로 \textbf{직접 작성}한다
\end{enumerate}
\end{column}
\end{columns}
\keymsg{자격시험 수업이 아니다 — 기초는 다루지 않고, 판단을 다룬다}
\note{여섯 목표 중 앞의 넷은 오전 이론, 뒤의 둘은 오후 실습에 대응한다. "기초는 다루지 않는다"는 교재 전제(주니어 PM과 3\til{}7년차 혼합 기수)를 명시하고, 이 과정은 자격시험 수업이 아니라는 점을 말한다 — 표준이 "이렇게 하라"고만 말하는 지점에서 "왜 그런가, 언제 틀리는가"를 다루는 과정이다. 목표 1은 온담 ONE ONDAM이 왜 프로그램인지 설명할 수 있다는 뜻이고, 목표 6은 워크북 다섯 산출물이다. 교재 서문·제4장 4.1 방향 5(성찰적 실무자) 참조. 예상 소요 3분.}
\end{frame}

% =====================================================================
\section{0교시 (20분) — 사례 브리핑: 온담F\&B홀딩스와 P1}
% =====================================================================
\begin{frame}{사례를 어떻게 쓰는가 — 통독하지 않는다}
\begin{itemize}
\item 사례 문서(25쪽)는 \textbf{사전}이다 — 소설처럼 읽고 외우는 문서가 아니다
\item 교재의 온담 단락마다 \textbf{■ 사례 사실} 줄이 그 단락이 쓰는 사실과 출처(§·숫자 표 ID)를 먼저 준다
\item 워크북의 산출물마다 \textbf{◆ 실습 입력 카드}가 필요한 정본 §·앞 산출물·추가 조건을 모아 준다
\item 책상 위에는 \textbf{사례 한 장}(책받침) — 타임라인·인물 13명·숫자 12
\item 오늘 20분 동안 그 한 장을 같이 읽는다 — 회사 / 왜 P1인가 / 날짜 / 사람
\end{itemize}
\keymsg{사례를 다 읽어 왔는지 묻지 않는다 — 다음 네 장을 보면 오늘 수업을 따라올 수 있다}
\note{학습안내(공통/00\_학습안내.pdf) §5의 요지. 사전에 사례 §1.1·§3.4·§4.2·§6만 읽어 오라고 했지만, 안 읽어 온 학생도 이 교시로 따라올 수 있어야 한다. 책받침(공통/01\_사례\_한장.pdf)을 나눠 주고 1면 왼쪽 위부터 같이 본다. 예상 소요 2분.}
\end{frame}

\begin{frame}{온담F\&B홀딩스 — 어떤 회사인가}
\begin{columns}[T]
\begin{column}{0.5\textwidth}
\subhead{규모와 돈}
\begin{itemize}
\item 카페·베이커리·델리 3브랜드 \textbf{615개 매장}(직영 227 / 가맹 388) + B2B 원료 유통
\item 2025 매출 \textbf{8,740억}, 영업이익 312억(\textbf{3.6\%}) — 이익이 얇다
\item 부채비율 \textbf{187\%}, 사모펀드 350억 → \textbf{2027년 IPO} 조항
\end{itemize}
\end{column}
\begin{column}{0.5\textwidth}
\subhead{왜 P1인가}
\begin{itemize}
\item 매장 시스템 \textbf{OD-POS 2.1}(2011 자체개발) 수명 종료, 원 개발자 1명 재직
\item 인터페이스 47본 중 \textbf{13본 설계문서 없음}, 재고를 실시간으로 못 본다
\item 두 번의 실패 — \textbf{2019 스마트오더 34억 중단}, \textbf{2022 WMS 12억 축소}
\end{itemize}
\end{column}
\end{columns}
\keymsg{돈이 얇고 시계가 빠른 회사가 세 번째 시도를 한다 — 이사회는 "세 번째 실패는 없다"를 문서로 걸었다}
\note{책받침 1면 왼쪽. 사례 §1.1·§1.5·§2.1·§6.5. 두 번의 실패가 남긴 감사 지적 세 가지(편익 근거 부재 · 손실 확정 절차 부재 · 형상관리 미비)가 이사회 승인 조건 5항이고, 오늘 쓰는 헌장·비즈니스 케이스가 그 답이라는 것을 여기서 한 번 말해 둔다. 재고 정확도 91.3\%가 순환실사 SKU 69\%만의 숫자라는 점은 Day4까지 따라오는 복선 — 지금은 "분모가 있다"고만. 예상 소요 4분.}
\end{frame}

\begin{frame}{P1 — 무엇을, 얼마에, 언제까지}
\begin{columns}[T]
\begin{column}{0.52\textwidth}
\subhead{범위 — 다섯 산출물}
\begin{itemize}
\item ① 클라우드 POS·매장 운영 ② 주문 허브 ③ 통합 재고 ④ 가맹 발주 포털 ⑤ 통합 계층·데이터 전환(1,842 테이블·7.4TB)
\item \textbf{제외}: SAP ECC 교체 · 물류 설비(P2) · 교육(P3) · 규제 등록(P4) · 운영 조직(P5)
\end{itemize}
\end{column}
\begin{column}{0.46\textwidth}
\subhead{돈과 시간}
\begin{itemize}
\item 예산 \textbf{3층} — 승인 168.0 / 집행 한도 156.0 / 계획 원가 145.8
\item SI 계약 68.2억 = \textbf{FP 12,400 × 55만 원} 고정가, 수행사 한빛디지털
\item 2026-01-05 → 2027-05-29, 스프린트 36회 + 게이트 G0\til{}G5 — \textbf{하이브리드}
\end{itemize}
\end{column}
\end{columns}
\keymsg{P1은 프로그램 ONE ONDAM(428억, P1\til{}P5)의 척추다 — 편익 214억은 P1 혼자 만들지 않는다}
\note{책받침 1면 오른쪽 위와 왼쪽 아래. 사례 §3.1\til{}3.3, §2.5. 예산 3층은 오늘 헌장 항목 8과 Day2 원가 기준선의 핵심이므로 세 숫자를 칠판에 적어 둔다. 제외 범위의 각 항목에 담당 프로젝트가 붙어 있다는 것 — "제외에는 주인이 있다" — 이 헌장 항목 5의 요점. 예상 소요 4분.}
\end{frame}

\begin{frame}{P1 타임라인 — 움직일 수 없는 날짜 셋}
\fig[0.62]{fig_0_1_case_timeline.pdf}
\keymsg{2026-09-11 · 2027-02-26\til{}28 · 2027-03-31 — 이 셋은 협상 대상이 아니다. 컷오버 창은 IPO 예비심사에서 거꾸로 센 날짜다}
\note{책받침 1면 타임라인. 사례 §3.4·§3.5·§5.4. 붉은 세 날짜를 칠판에 적고 "이 중 미룰 수 있는 게 있나"를 묻는다 — 없다. 그것이 1교시 NTCP의 Pace = time-critical 판정의 근거이고, 성수기 잠금(7\til{}8월·12월·명절 전 2주) 때문에 컷오버 창이 1\til{}2월 아니면 3\til{}5월뿐이라는 것이 헌장 §9 제약이다. 초록색 P2(물류센터 설비, 194억)는 P1 밖의 의존 — 리드타임 34→41주 통지가 이미 와 있다. 예상 소요 5분.}
\end{frame}

\begin{frame}{사람 13명과 세 개의 시계}
\begin{columns}[T]
\begin{column}{0.54\textwidth}
\subhead{오늘 이름을 알아야 할 여덟}
\begin{itemize}\small
\item \textbf{김선호} 대표 — 이사회 의장, "일정이 먼저다" · \textbf{정미란} CFO — 스폰서·CCB 위원장, "148억으로"
\item \textbf{오정근} 영업본부장 — Senior User, "성수기 절대 금지" · \textbf{박세연} IT본부장 — Senior Supplier, OD-POS 원 설계자
\item \textbf{한동석} 물류본부장 — 3PL 계약 소관 · \textbf{박준영} 서비스운영팀장 — SAC 판정권
\item \textbf{서정필} 가맹점주협의회장 — 388점 서면 동의 · \textbf{이재현} 수행사 PM — "대기·추가는 변경대가"
\end{itemize}
\end{column}
\begin{column}{0.44\textwidth}
\subhead{세 개의 시계 (사례 §6.3)}
\begin{itemize}\small
\item \textbf{대표} — IPO: 2027-3Q 예비심사 전 안정 가동 → "더 앞당기라"
\item \textbf{CFO} — 부채비율 187 → 160\%: 1차연도 176억을 \textbf{148억}으로
\item \textbf{수행사·감리} — FP 12,400 안에서 손익률 6.5\%: 당기면 변경대가, 줄이면 재협상
\end{itemize}
\end{column}
\end{columns}
\keymsg{세 시계는 맞지 않는다 — PM이 맞추는 게 아니라 헌장의 우선순위 칸에 적어 스폰서가 결정하게 한다}
\note{책받침 2면. 사례 §4.2·§6.3. 나머지 다섯(나영선·최영주·강윤지·윤태호·김미르)은 3교시 이해관계자 절에서 등장하므로 여기서는 "책받침에 있다"고만. 세 개의 시계는 1교시 NTCP Pace 축, 3교시 헌장 항목 13(우선순위), 6교시 프리모템까지 세 번 되돌아오는 장치 — 처음 말할 때 "이 충돌은 오늘 안에 해결되지 않고 문서에 적힌다"고 예고한다. 예상 소요 5분.}
\end{frame}

% =====================================================================
\section{1교시 (80분) — 프로젝트란 무엇인가, 왜 실패하는가}
% =====================================================================

\begin{frame}{이 교시에서 배우는 것}
\begin{itemize}
\item 이 과정에서 "PM"이 무엇을 뜻하는가 — 프로젝트 매니저 / PO / 프로덕트 매니저의 3중 충돌
\item 프로젝트·프로그램·포트폴리오·운영 — 세 표준의 정의를 나란히 놓으면 보이는 것 [§1.1]
\item 프로젝트는 왜 실패하는가 — 통계를 그대로 믿으면 안 되는 이유 [§1.2]
\item 프로젝트를 임시조직으로 본다는 것 [§1.3]
\item 유형은 하나가 아니다 — NTCP 다이아몬드 · 분야별 되돌림 비용 [§1.4\til{}1.5]
\end{itemize}
\keymsg{§1.2(실패 통계)가 가장 길다 — 경력자가 가장 많이 틀리는 지점}
\note{제1장 전체를 90분에 다룬다. 시간 배분은 용어 정리 5분, §1.1 15분, §1.2 25분, §1.3 15분, §1.4 15분, §1.5 10분, 정리 5분. §1.2(실패 통계)가 가장 길다 — 경력자가 가장 많이 틀리는 지점이기 때문이다. §1.4 NTCP는 워크북 §1과 직결되고, §1.5 되돌림 비용은 2교시 테일러링의 토대가 된다. 예상 소요 2분.}
\end{frame}

\begin{frame}{용어 정리 — 이 과정에서 "PM"은 프로젝트 매니저다}
\begin{tbl}{L{0.17\textwidth}L{0.2\textwidth}L{0.34\textwidth}L{0.19\textwidth}}
\textbf{약칭} & \textbf{원어} & \textbf{책임지는 것} & \textbf{이 과정에서} \\ \midrule
\textbf{PM} & Project Manager & 정해진 사양·일정·예산 안에서 \textbf{산출물(output)} 인도 & \textbf{PM 트랙 4일 (오늘)} \\
PO & Product Owner (Scrum) & 백로그 우선순위, 스프린트 단위 가치 판단 & Product 트랙 \\
프로덕트 매니저 & Product Manager & 제품의 존재 이유 — 시장·사용자·수명 전체의 \textbf{성과(outcome)} & Product 트랙 4일 \\
\end{tbl}
\begin{itemize}
\item 명함이 아니라 \textbf{무엇에 책임지는가}로 구분한다
\item 오늘 이후 "PM" = 온담 P1의 \textbf{발주사} 프로젝트 매니저(여러분)
\end{itemize}
\keymsg{output은 PM의 책임, outcome은 사업 오너의 책임 — 오늘 하루 반복되는 구분}
\note{이 슬라이드는 혼합 기수를 위한 오리엔테이션이며 교재 본문 밖이다. 한국 기업에서 세 역할이 "PM"이라는 같은 글자로 불리면서 충돌한다 — 명함이 아니라 무엇에 책임지는가로 구분한다. 깊이 들어가지 말고, "output은 PM의 책임, outcome은 사업 오너의 책임"이라는 3교시 §3.1 ITO 모델의 예고편으로만 쓴다. 산출물/성과의 구분이 오늘 반복해서 나온다는 것만 심어 둔다. 예상 소요 5분.}
\end{frame}

\begin{frame}{프로젝트·프로그램·포트폴리오 — 세 표준의 정의}
\begin{tbl}[\footnotesize]{L{0.15\textwidth}L{0.27\textwidth}L{0.27\textwidth}L{0.22\textwidth}}
\textbf{표준} & \textbf{프로젝트} & \textbf{프로그램} & \textbf{포트폴리오} \\ \midrule
\textbf{ISO 21502:2020} & 합의된 목표를 위한 임시적 노력. 시작·끝이 정의됨 & 개별 관리로는 얻을 수 없는 \textbf{편익}을 위해 함께 관리 & 전략 목표를 위한 프로젝트·프로그램·\textbf{운영}의 집합 \\
\textbf{PMBOK 8판} (2025) & 고유한 산출물을 위한 임시적 노력 + \textbf{가치} 인도 & 편익과 통제를 위해 조정 관리 & 전략 목표를 위한 그룹 관리 \\
\textbf{MSP 5판} (2020) & 산출물(outputs)을 만드는 것 & outputs \ra{} capabilities \ra{} outcomes \ra{} \textbf{benefits} & (MoP가 다룸) \\
\end{tbl}
\begin{itemize}
\item 공통 단어 셋: \textbf{임시적}(끝이 정의되어 있다) / \textbf{편익} / \textbf{전략·운영}
\end{itemize}
\keymsg{정의는 개념 논쟁이 아니다 — 프로젝트라 부르는 순간 예산·책임자·해산이 생긴다}
\note{§1.1 개념. "누가 돈을 대고 누가 책임지는가"의 논쟁임을 먼저 말한다. "임시적(temporary)"은 짧다는 뜻이 아니라 끝이 정의되어 있다는 뜻이다. 프로젝트라 부르는 순간 예산·책임자·해산이 생기고, 운영이라 부르면 연간 예산 안에서 반복된다. outputs\ra{}capabilities\ra{}outcomes\ra{}benefits 네 단어 사슬은 오늘 하루 종일 반복되므로 칠판에 남겨 둔다. PRINCE2 7의 정의("비즈니스 케이스에 따라 … 만들어진 임시조직")도 언급 — "임시조직"과 "비즈니스 케이스"가 §1.3과 §3.1의 예고다. 예상 소요 8분.}
\end{frame}

\begin{frame}{온담으로 읽기 — ONE ONDAM은 프로그램이다}
\fig[0.7]{fig_1_1_hierarchy}
\figcap{그림 1-1. 포트폴리오·프로그램·프로젝트·운영의 관계와 MSP 가치 사슬 (ONE ONDAM 428억 · P1\til{}P5)}
\keymsg{편익이 프로젝트 경계를 넘어 발생한다는 사실 자체가 프로그램의 존재 이유다}
\note{§1.1 "온담의 사례로 읽기". 총 428억, 다섯 프로젝트 — P1 플랫폼 168억 · P2 물류 자동화 194억 · P3 변화관리 22억 · P4 규제 대응 31억 · P5 운영 이관 13억. 이사회가 승인한 편익은 M+24 연간 절감 214억이고, 어느 한 프로젝트가 이것을 만들지 않는다. 그림의 가치 사슬(outputs\ra{}capabilities\ra{}outcomes\ra{}benefits)을 손으로 짚으며 다음 슬라이드의 폐기율 43억 사슬로 넘어간다. PMBOK 8판 첫 원칙 "총체적 관점"이 바로 이 관계를 보라는 요구다. 수강생 회사의 "TF"가 정의상 프로그램인지 묻는다(생각해 볼 질문 1). 예상 소요 4분.}
\end{frame}

\begin{frame}{P1의 편익은 P1 혼자 만들지 못한다}
\begin{itemize}
\item 폐기율 절감 \textbf{43억} = P1 실시간 재고(산출물) \ra{} 매장·물류의 발주 행동 변화(역량\ra{}성과) \ra{} 편익
\item 물류 처리량 \textbf{96억} = P2 설비 + P1 재고 원장이 \textbf{함께} 작동해야 나온다
\item P1의 산출물은 명확하다 — POS, 주문 허브, 통합 재고, 가맹 포털, 통합 계층·데이터 전환
\item 운영과의 경계: P5 종료(2027-05-29) 이후는 서비스운영팀(박준영, 11명)의 것
\item 연 운영비 27억은 428억 \textbf{밖}
\end{itemize}
\keymsg{프로젝트가 끝나는 지점과 운영이 시작되는 지점을 착수 시점에 정의하지 않으면 종료 때 반드시 다툰다}
\note{§1.1 "온담의 사례로 읽기" 후반. 편익 214억을 항목별로 뜯으면 어느 것도 P1 산출물만으로 나오지 않는다 — 폐기율은 매장의 발주 행동이 바뀌어야(P3 교육·영업본부 발주 정책), 물류 처리량은 P2 설비가 있어야 나온다. 이것이 3교시 ITO·BDN의 예고다. 운영과의 경계도 착수 시점에 그어야 한다 — P5 종료 이후는 서비스운영팀의 것이고 연 운영비 27억은 프로그램 예산 밖이다. "프로젝트가 끝나는 지점과 운영이 시작되는 지점을 착수 시점에 정의하지 않으면 종료 시점에 반드시 다툼이 일어난다"를 말하고 Day4(종료·이관)를 예고한다. 예상 소요 4분.}
\end{frame}

\begin{frame}{정의를 틀리는 두 방향 — 프로그램을 프로젝트로, 운영을 프로젝트로}
\begin{columns}[T]
\begin{column}{0.49\textwidth}
\subhead{프로그램을 프로젝트 도구로}
\begin{itemize}
\item "전사 혁신 과제", "DX TF"
\item 편익 소유자 소멸
\item 프로젝트 간 의존 미관리
\item 종료 조건이 흐려진다
\end{itemize}
\end{column}
\begin{column}{0.49\textwidth}
\subhead{운영을 프로젝트로 포장}
\begin{itemize}
\item 연례 업그레이드에 헌장·킥오프·종료 보고서
\item 반복되고 끝이 없는 일은 \textbf{운영}
\item 그것은 ITIL의 영역이다
\end{itemize}
\end{column}
\end{columns}
\keymsg{여러분 회사에서 "프로젝트"라 불리지만 정의상 프로그램이거나 운영인 것은?}
\note{§1.1 "실무에서 어떻게 틀리는가". 첫째 방향 — 프로그램을 프로젝트 도구로 관리하면 편익 소유자가 사라지고 프로젝트 간 의존이 관리되지 않으며 종료 조건이 흐려진다. 한국 기업의 "전사 혁신 과제", "DX TF"가 대개 이 상태다. 둘째 방향 — 운영을 프로젝트로 포장하면 연례 업그레이드마다 헌장·킥오프·종료 보고서를 쓰는 낭비가 생긴다. 반복되고 끝이 없는 일은 운영이고 ITIL의 영역이다. 쉬는 시간 전 과제로 "여러분 회사에서 프로젝트라 불리지만 정의상 프로그램이거나 운영인 것"을 하나씩 떠올려 오게 한다. 예상 소요 3분.}
\end{frame}

\begin{frame}{프로젝트는 왜 실패하는가 — "성공률 30\%"를 믿으면 안 되는 이유}
\begin{columns}[T]
\begin{column}{0.58\textwidth}
\fig[0.74]{fig_1_2_chaos_bias}
\figcap{그림 1-2. CHAOS 성공 정의의 결함 — 6\% / 94\% / 70\%}
\end{column}
\begin{column}{0.4\textwidth}
\begin{itemize}
\item Standish CHAOS(1994\til{}): 성공 16\% · challenged 53\% — 30년간 업계 상식
\item \textbf{Eveleens \& Verhoef (2010)}: "성공" = 최초 추정 안에 완료 \ra{} 추정 편향의 \textbf{방향}을 잴 뿐
\item 추정을 부풀리는 조직 70\% 성공 / 정확하되 약간 과소 6\% / 편향만 뒤집으면 94\%
\end{itemize}
\end{column}
\end{columns}
\keymsg{결론은 "프로젝트는 실패하지 않는다"가 아니다 — 이 숫자로는 아무것도 알 수 없다}
\note{§1.2 "흔히 인용되는 숫자"와 "Eveleens \& Verhoef (2010, IEEE Software)". Standish의 "성공"은 최초 추정한 시간·비용·기능 안에 완료한 것이므로, 추정을 10\til{}100배 부풀리는 조직(Organization X)은 약 70\% 성공, 추정이 정확하지만 약간 과소인 조직은 6\%, 편향만 뒤집으면 94\%가 된다. 이 지표는 관리 품질이 아니라 추정이 어느 방향으로 편향되어 있는가를 측정한다. 네 가지 문제 — 오도적 / 일방적 / 추정 관행을 왜곡 / 의미 없는 숫자 — 중 세 번째가 실무에서 가장 중요하다: 이 지표를 KPI로 쓰는 조직에서 합리적 행동은 추정을 부풀리는 것이다. "여러분 회사가 일정 준수율을 PM 평가 지표로 쓰고 있다면 같은 메커니즘이 이미 작동하고 있다"를 던진다(생각해 볼 질문 2). Jørgensen \& Moløkken-Østvold(2006)의 "189\% 초과" 재검토도 한 줄. 예상 소요 8분.}
\end{frame}

\begin{frame}{Flyvbjerg — 다른 데이터, 다른 결론}
\fig[0.63]{fig_1_2_flyvbjerg_multiply}
\figcap{그림 1-3. Flyvbjerg의 두 관찰 — iron law와 세 차원의 곱셈 구조 (1/1,000 · 8/1,000)}
\keymsg{조직은 대개 한 차원만 관리하면서 성공이라 부른다 — 2019 스마트오더: 예산 안, 8개월 만에 중단}
\note{§1.2 "Flyvbjerg". 응답 설문이 아니라 실제 인프라 프로젝트의 최초 승인 예산 vs 최종 실적을 수십 년 수집한 데이터다. iron law("over budget, over time, over and over again" — 열 중 아홉이 원가 초과)와 곱셈 구조(원가·일정·편익 각 1/10이면 셋 동시 만족은 약 1/1,000, 관대하게 2/10씩이면 8/1,000)는 자주 뒤섞여 인용되므로 분리해서 말한다. 조직은 대개 한 차원(일정)만 관리하면서 성공이라 부른다 — 온담 2019 스마트오더는 34억 예산 안에서 8개월 만에 중단됐다(원가 성공, 편익 실패). 영국 재무부 Green Book·DfT가 RCF를 공식 지침에 탑재했다는 사실로 "학술 논쟁이 아니라 제도"임을 보인다. 예상 소요 4분.}
\end{frame}

\begin{frame}{처방 — Reference Class Forecasting: 외부 분포 먼저, 내부 추정 나중}
\fig[0.5]{fig_1_2_rcf}
\figcap{그림 1-4. RCF 3단계 — ① 유사 과거 20\til{}30건 \ra{} ② 실적÷최초 예산의 분포 \ra{} ③ uplift·상대 위험 판단}
\begin{itemize}
\item 순서를 뒤집으면 \textbf{앵커링} / reference class를 너무 좁게 잡으면 \textbf{uniqueness bias}
\end{itemize}
\keymsg{IT 원가 초과는 멱함수 분포 — 평균에 맞춘 컨틴전시가 틀리는 이유 (Day2·Day3)}
\note{§1.2 "Flyvbjerg" 후반. RCF의 세 단계 — 유사 과거 프로젝트 20\til{}30건 확보, 관심 변수(실적÷최초 예산)의 분포, 그 분포에서 uplift 또는 상대 위험 판단 — 에서 핵심은 순서다: 외부 분포를 먼저 보고 내부 추정은 나중에. 반대로 하면 자기 추정에 앵커링된다. reference class를 너무 좁게 잡으면 "우리는 특별하다"는 uniqueness bias가 들어온다. IT 원가 초과가 멱함수 분포(2022 JMIS, 5,392건)라는 점 — 평균에 맞춘 컨틴전시가 틀린 이유 — 은 Day2 원가·Day3 리스크 예고로 한 줄만(그림 1-5 powerlaw는 그때 쓴다). 예상 소요 4분.}
\end{frame}

\begin{frame}{그러면 Flyvbjerg의 숫자는 믿어도 되는가 — 자기 조직의 base rate}
\begin{columns}[T]
\begin{column}{0.5\textwidth}
\begin{itemize}
\item Flyvbjerg 표본은 교통·에너지·올림픽 중심. 168억 SI에 그대로 적용은 \textbf{class를 너무 넓게} 잡는 반대쪽 오류
\item 배울 것은 숫자가 아니라 태도 — \textbf{남의 통계 대신 자기 조직의 base rate}
\item 네 건은 class로는 적다. 그러나 "정보화 4건 중 2건 중단·축소"는 P1 리스크 등록부에 들어갈 base rate
\end{itemize}
\end{column}
\begin{column}{0.48\textwidth}
\begin{tbl}[\footnotesize]{lL{0.19\textwidth}rl}
\textbf{연도} & \textbf{사업} & \textbf{규모} & \textbf{결과} \\ \midrule
2014 & SAP 도입 & 74억 & (실적 미기재) \\
2019 & 스마트오더 & 34억 & \textbf{중단} (8개월) \\
2022 & WMS 고도화 & 12억 & \textbf{범위 63\% 축소} \\
2025 & 검사기록 디지털화 & 34억 & 반려 \\
\end{tbl}
\end{column}
\end{columns}
\keymsg{오늘의 첫 실무 처방 — 제안서·임원 보고서에서 CHAOS 수치를 빼라}
\note{§1.2 "그러면 Flyvbjerg의 숫자는 믿어도 되는가"와 "실무에서 어떻게 틀리는가". 2022년 IT 5,392건도 특정 DB의 표본이다. 온담의 과거 4건 중 2014 SAP는 계획 대비 실적이 미기재라 워크북에서 가정한다. 통계 인용의 네 가지 오류 — 임원 보고서에 CHAOS를 "업계 표준 실패율"로 인용(그 자리에서 신뢰를 잃는다) / "그러니 PM은 소용없다"는 결론 / iron law를 냉소로("best practice는 예외, average practice가 재앙"이 실제 결론) / 자기 base rate를 만들지 않음. 이사회 승인 조건 5항(내부감사 지적 재발 방지)이 같은 인식이라는 점을 연결한다. 워크북 §1 심화 과제(과거 4건을 reference class 출발점으로)와 6교시 프리모템의 세 번째 원천(base rate)을 예고. 예상 소요 5분.}
\end{frame}

\begin{frame}{성공의 두 정의 — Cooke-Davies (2002)}
\fig[0.63]{fig_1_2_cooke_davies}
\figcap{그림 1-6. 성공 세 정의와 2×2 조합 — 관리 성공은 G4(2027-05-29), 프로젝트 성공은 G5(2028-03-31)에 판정}
\keymsg{세 답은 서로 독립이다 — 온타임·온버짓인데 아무도 안 쓰는 시스템 / 지연·초과인데 회사를 살린 시스템}
\note{§1.2 "성공의 두 정의". 온담 P1에서 관리 성공은 156억 집행 한도·다섯 산출물·SAC 충족, 프로젝트 성공은 폐기율 2.2\%·결품률 1.5\%·리드타임 12h, 지속적 성공은 온담의 과거 4건이 답한다. 이 구분이 없으면 PRINCE2 비즈니스 케이스, PMBOK 8판의 가치, 편익 관리 논의가 전부 뒤섞인다. 과정 전 제출 과제("최근 종료 프로젝트의 온타임·온버짓 O/X와 편익 달성 O/X")를 여기서 꺼낸다 — 두 답이 다른 프로젝트가 있는 수강생에게 한 문장씩 말하게 한다(5분 토론). 이 구분은 헌장 항목 3(성공 기준 두 층)에서 그대로 쓰인다. 예상 소요 7분.}
\end{frame}

\begin{frame}{프로젝트를 임시조직으로 본다는 것 — Lundin \& Söderholm (1995)}
\fig[0.66]{fig_1_3_temporary_org}
\figcap{그림 1-7. 임시조직 이론 — 4T · 4 행위 개념 · 실무 현상}
\keymsg{라이프사이클 그림이 설명 못 하는 것 — 왜 프로젝트 안의 사람들이 그렇게 행동하는가}
\note{§1.3 개념. 네 가지 구조 개념(4T): 시간 — 영구조직에서 시간은 흐르지만 임시조직에서는 소모된다 / 과업 — 과업이 끝나면 존재 이유가 사라진다 / 팀 — 소속감은 과업에 묶여 있고 과업 후 흩어진다 / 전환 — 이전 상태에서 이후 상태로 바꾸기 위해 존재한다. 네 가지 행위 개념: 행위 기반 기업가정신(착수는 분석이 아니라 행위의 결과) / 몰입 형성을 위한 단편화(WBS는 계획 도구이기 전에 몰입 도구) / 계획된 고립(별도 사무실·예산·보고선) / 제도화된 종료(종료 보고서, 해산, 복귀). 흔한 오류: "임시조직"을 "한시적 조직"이라는 사전적 의미로만 쓰는 것, 4T를 헌장 체크리스트로 만드는 것. 예상 소요 6분.}
\end{frame}

\begin{frame}{이 렌즈가 설명하는 네 가지 현상}
\begin{tbl}[\footnotesize]{L{0.3\textwidth}L{0.3\textwidth}L{0.32\textwidth}}
\textbf{실무자가 매일 겪는 것} & \textbf{설명하는 개념} & \textbf{온담에서} \\ \midrule
교훈(lessons learned)은 왜 항상 실패하는가 & 제도화된 종료 — 배울 조직이 종료와 함께 해체 & 스마트오더 종료 보고서의 중단 사유 세 줄 \\
종료 직전에 왜 인력이 빠지는가 & 팀 — 끝이 보이면 다음 과업을 찾는 것이 정상 & 안정화 기간의 핵심 개발자 이탈 \\
프로젝트는 왜 모조직과 싸우는가 & 계획된 고립 — 프로젝트의 고립은 운영의 부담 & IT본부 47명 중 14명 투입, 33명이 6종·47본 운영 \\
착수 결정은 왜 분석적이지 않은가 & 행위 기반 기업가정신 & "615개 매장을 하나의 채널처럼"이 먼저 \\
\end{tbl}
\keymsg{종료 이후 경로(P5, 다음 배치)를 착수 시점에 설계해야 인력 이탈을 막는다}
\note{§1.3 "왜 유용한가". 표의 네 행은 이름 없던 현상에 이름을 붙이는 것이다 — 교훈이 실패하는 것은 담당자의 게으름이 아니라 배울 조직이 종료와 함께 해체되기 때문이고, 종료 직전 인력 이탈은 배신이 아니라 정상 작동이다. IT본부 47명 중 14명이 투입되면 남은 33명이 시스템 6종·인터페이스 47본을 운영해야 한다 — 프로젝트의 고립은 운영의 부담이다. 김선호 대표의 "매장 615개를 하나의 채널처럼"이 먼저 있었고 비즈니스 케이스는 나중이었다. 처방 한 줄: 종료 이후 경로(P5, 다음 배치)를 착수 시점에 설계해야 인력 이탈을 막는다. 예상 소요 4분.}
\end{frame}

\begin{frame}{어떤 프로젝트도 섬이 아니다 — Engwall (2003)}
\begin{itemize}
\item 일정·예산·조직 구성·핵심 갈등은 착수 전에 \textbf{조직의 역사와 정치}로 이미 결정되어 있다
\item 컷오버 2027-02 = 2019 사모펀드 계약의 IPO 조항
\item 인터페이스 우선 = 2011 OD-POS를 박세연이 직접 만듦
\item 가맹 경계 = 2019 스마트오더 중단 경험
\item 처방: 새 프로젝트 첫 주에 \textbf{직전 프로젝트 두세 개의 종료 보고서}를 읽어라
\end{itemize}
\keymsg{킥오프 시점 PM의 진짜 일은 이미 결정된 것들을 목록화하는 것}
\note{§1.3 "어떤 프로젝트도 섬이 아니다". Engwall은 "맥락이 중요하다"는 상투구가 아니라 인과 주장임을 강조한다 — 직전 프로젝트의 정치적 결과가 이번 프로젝트의 특징을 설명한다. 온담의 세 예시(컷오버 날짜, 인터페이스 우선순위, 가맹 경계)는 모두 P1 착수 전에 정해졌고 PM이 바꿀 수 없다. 그래서 킥오프 시점 PM의 진짜 일은 새로 정하는 것이 아니라 이미 결정된 것들을 목록화하는 것이다 — 3교시 헌장의 가정·제약 항목과 마무리 §4.2 Morris로 이어진다. 생각해 볼 질문 4(착수 전 결정된 비율). 예상 소요 5분.}
\end{frame}

\begin{frame}{프로젝트 유형은 하나가 아니다 — NTCP 다이아몬드 (Shenhar \& Dvir, 2007)}
\begin{tbl}[\footnotesize]{L{0.13\textwidth}L{0.24\textwidth}L{0.27\textwidth}L{0.28\textwidth}}
\textbf{축} & \textbf{질문} & \textbf{단계} & \textbf{이 축이 결정하는 것} \\ \midrule
\textbf{Novelty} & 시장·사용자에게 얼마나 새로운가 & derivative / platform / breakthrough (\textbf{3}) & 요구사항 확정 시점, 시장 조사 신뢰도 \\
\textbf{Technology} & 기술이 얼마나 성숙한가 & low / medium / high / super-high (\textbf{4}) & 설계 동결 시점, 설계 반복 횟수 \\
\textbf{Complexity} & 범위·계층이 얼마나 넓은가 & assembly / system / array (\textbf{3}) & 조직·절차의 형식성, 하도급·인터페이스 관리 \\
\textbf{Pace} & 시간이 얼마나 결정적인가 & regular / fast-competitive / time-critical / blitz (\textbf{4}) & 의사결정 위임 수준, 상위 관리자 개입 \\
\end{tbl}
\keymsg{네 축은 직교한다 — 합산해 "난이도 3.2등급"을 만드는 순간 모델은 무의미해진다}
\note{§1.4 개념. 단계 수 3/4/3/4를 정확히 — 교재 오류가 자주 나는 지점이다. 다이아몬드가 바깥쪽일수록 리스크가 크고 관리는 더 유연해야 하며, 축마다 다른 처방이 붙는다. PMBOK 7·8이 테일러링을 강조하면서 판단 기준을 주지 않는 공백을 이 모델이 메운다고 위치시킨다. 예상 소요 5분.}
\end{frame}

\begin{frame}{온담 P1을 NTCP로 진단하면 — "표준 SI 프로젝트"가 아니다}
\begin{columns}[T]
\begin{column}{0.55\textwidth}
\fig[0.72]{fig_1_4_ntcp}
\figcap{그림 1-8. NTCP 다이아몬드 — 온담 P1의 위치}
\end{column}
\begin{column}{0.43\textwidth}
\begin{itemize}
\item \textbf{N = platform} — 사용자 3,000명대의 업무 방식이 바뀐다
\item \textbf{T = medium} — 기술은 성숙, \textbf{기존 기술의 미지}(문서 없는 13본, 7.4TB)가 크다
\item \textbf{C = array에 가까운 system} · \textbf{P = time-critical} — IPO·규제 역산
\end{itemize}
\end{column}
\end{columns}
\keymsg{기술 리스크보다 통합·일정 리스크에 자원 집중 — 같은 프로그램의 P2는 전혀 다른 프로파일}
\note{§1.4 "온담 P1을 이 모델로 진단하면". Novelty = platform: 옴니채널·클라우드 POS는 시장에 흔하지만 온담 사용자(매장 1,510명·가맹 331명·B2B 260개사)에게는 업무 방식이 바뀌는 플랫폼급이고, 로컬 DB+야간 동기화에서 중앙 실시간으로 가므로 derivative가 아니다 \ra{} 요구사항은 초기에 대부분 확정 가능하되 사용자 검증 여러 번(스프린트 36회와 정합). Technology = medium: 클라우드·API GW·MDM·SAP 애드온은 성숙하지만 문서 없는 인터페이스 13본, 1,842 테이블·7.4TB, 3PL WMS라는 기존 기술의 미지가 크다 \ra{} 설계 동결 1\til{}2회(설계 스프린트 6회, 전환 리허설 3회). Complexity: P1 자체는 system, 프로그램 관점(P2·P4·P5 인터페이스, 벤더 11개사, 별도 법인, 3PL)에서는 array \ra{} 형식적 계약·문서·인터페이스 관리, 별도 통합 조직(프로그램 PMO 4 + P1 PMO 3). Pace = time-critical: 컷오버 2027-02-26\til{}28은 IPO 역산, 규제 2026-09-11·2027-03-31은 협상 불가 \ra{} 일정 최우선, 권한 위임, 상위 관리자 밀착. 대표("더 앞당기라")와 CFO("148억으로")가 이 축에서 정면 충돌하므로 헌장 우선순위 항목이 적어야 한다. 이 판단은 예시이며 워크북에서 수강생이 다시 판단한다. 오류 셋: 합산, 착수 때 한 번만 그림(Pace는 중간에 승격), 그림만 그리고 처방을 안 정함 — 처방 규칙표가 2교시 TDR이다. 생각해 볼 질문 5(system vs array). 예상 소요 10분.}
\end{frame}

\begin{frame}{분야별 대조 — 같은 단어, 다른 리듬, 다른 되돌림 비용}
\fig[0.63]{fig_1_5_reversal_cost}
\figcap{그림 1-9. 분야별 되돌림 비용 시간축 대조 (SI · EPC · NPD · 제약)}
\keymsg{되돌릴 수 있으면 실험하고 반복한다. 되돌릴 수 없으면 결정 전에 최대한 알아야 하므로 앞단에 시간을 쓴다}
\note{§1.5 도입. 그림의 시간축에서 각 분야가 되돌릴 수 없게 되는 지점이 어디인지 짚는다 — SI는 데이터 전환·컷오버, EPC는 기초 타설, NPD는 금형과 양산, 제약은 규제 제출. SI의 되돌릴 수 없는 두 지점(데이터 전환 정본 610테이블 재매핑 3.8억·6주 / 컷오버 34시간 615매장 펌웨어) 앞에 G2·G3 게이트가 있는 이유를 말한다. 분야를 옮긴 PM의 실수 — 건설 자원 평준화·EVM CPI를 IT에 그대로, SI의 "일단 만들고 고치자"를 설비에 그대로(선급금 29.6억). 예상 소요 4분.}
\end{frame}

\begin{frame}{분야별 대조표 — 산출물 · 되돌림 비용 · 리듬 · 범위 확정 · 주된 실패}
\begin{tbl}[\scriptsize]{L{0.1\textwidth}L{0.19\textwidth}L{0.2\textwidth}L{0.19\textwidth}L{0.19\textwidth}}
\textbf{축} & \textbf{SI} & \textbf{건설·EPC} & \textbf{제조 NPD} & \textbf{제약 임상·규제} \\ \midrule
산출물 & 소프트웨어 — 수정 가능 & 물리 구조물 & 물리 제품 + 양산 공정 & \textbf{승인} — 규제기관의 결정 \\
\textbf{되돌림 비용} & 낮음\til{}중간 (전환·컷오버는 예외) & 매우 높음 — 기초 타설 후 변경은 철거 & 높음 — 금형 후 재작업, 양산 후 리콜 & 극단적 — 제출 후 문서 잠김, 시계는 상대가 \\
리듬을 정하는 것 & 요구사항 확정, 인터페이스 상대의 준비 & 자재 리드타임, 공정 순서, 인허가 & 금형·부품 리드타임, EVT/DVT/PVT & 심사 일정, 추가자료 왕복 횟수 \\
범위 확정 시점 & 늦다 — 계약 후에도 움직임 & 이르다 — 설계 성숙도(AACE Class)가 계약 유형 결정 & 게이트마다 단계적, Gate 1\til{}2에서 kill & 프로토콜 승인 시 사실상 고정 \\
주된 실패 & 요구 팽창, 데이터 전환, 인터페이스 미확정 & 설계 미성숙 상태의 LSTK, 동시지연 클레임 & 이빨 없는 게이트(kill 0건), 금형 뒤 설계 변경 & 추가자료 0회 가정 위의 출시일 공표 \\
\end{tbl}
\keymsg{가장 중요한 열은 되돌림 비용 — 이 표가 2교시 테일러링의 토대다}
\note{§1.5 표. 다섯 축 중 되돌림 비용 열을 먼저 읽고 나머지는 그 결과로 설명한다 — 되돌림 비용이 낮으면 범위 확정이 늦어도 되고(SI), 높으면 설계 성숙도가 계약 유형을 결정하며(EPC), 게이트에서 kill이 필요하고(NPD), 제출 전에 문서를 완결해야 한다(제약). 주된 실패 행은 5교시 §3.6 분야별 착수 노하우에서 각각 다시 나온다. 예상 소요 4분.}
\end{frame}

\begin{frame}{1교시 핵심 정리}
\begin{columns}[T]
\begin{column}{0.5\textwidth}
\begin{enumerate}
\item 프로젝트는 \textbf{산출물}, 프로그램은 산출물이 역량·성과를 거쳐 \textbf{편익}이 되도록 함께 관리
\item CHAOS 성공률은 \textbf{추정 편향의 방향}을 측정한다 — KPI로 쓰면 부풀리기가 최적 행동
\item iron law와 3차원 곱셈은 별개. 처방은 RCF — \textbf{외부 분포 먼저}
\item 관리 성공(시간·비용·범위)과 프로젝트 성공(편익)은 \textbf{독립}
\end{enumerate}
\end{column}
\begin{column}{0.48\textwidth}
\begin{enumerate}\setcounter{enumi}{4}
\item 임시조직론(4T + 4행위)이 교훈의 실패·종료 직전 이탈·모조직과의 긴장을 설명. 착수의 진짜 일은 \textbf{이미 결정된 것의 목록화}
\item NTCP(3/4/3/4)는 축별로 다른 처방. 합산하지 말고 게이트마다 다시 그린다
\item 분야 차이의 핵심 변수는 \textbf{되돌림 비용}
\end{enumerate}
\end{column}
\end{columns}
\keymsg{워크북 §1 "전략적 배경"과 심화 과제(과거 4건 base rate)에 쓰인다}
\note{제1장 핵심 정리 그대로. 편익이 프로젝트 경계를 넘는다는 사실이 프로그램의 존재 이유다(1). 어떤 프로젝트도 섬이 아니다(5). 워크북 연결 — 이 장은 워크북 §1 비즈니스 케이스 요약의 "전략적 배경"과 심화 과제(과거 4건 base rate)에 쓰인다. 쉬는 시간 전에 "여러분 회사에서 프로젝트라 불리지만 정의상 프로그램이거나 운영인 것"을 하나씩 떠올려 오게 한다. 예상 소요 4분.}
\end{frame}

% =====================================================================
\section{2교시 (90분) — 표준의 지형: 무엇을 따를 것인가}
% =====================================================================

\begin{frame}{이 교시에서 배우는 것}
\begin{itemize}
\item PMBOK의 계보 — 6판, 7판, 그리고 8판의 부분 회귀 · 6 원칙과 7 성과영역 [§2.1\til{}2.2]
\item PRINCE2 7 — 프로젝트 보드, 예외 관리, 산출물 기반 계획 [§2.3]
\item ISO 21502:2020 — 계약서가 인용하는 표준 · 애자일·SAFe·하이브리드의 위치 [§2.4\til{}2.5]
\item 테일러링 — 되돌림 비용, Boehm-Turner, Cynefin, TDR [§2.6]
\item 자격 체계의 실효성 [§2.7]
\end{itemize}
\keymsg{이 교시의 실무 산출물은 재매핑 대장과 TDR이다}
\note{제2장 전체. 시간 배분: §2.1 20분, §2.2 12분, §2.3 15분, §2.4 7분, §2.5 8분, §2.6 20분, §2.7 5분, 정리 3분. PMP 보유자는 §2.1에서 자기가 공부한 판본이 현행과 어떻게 다른지 확인하게 된다고 예고한다. 이 교시의 실무 산출물은 재매핑 대장과 TDR이다. 예상 소요 2분.}
\end{frame}

\begin{frame}{왜 계보를 알아야 하는가 — PMBOK은 참고서인데}
\begin{itemize}
\item PMBOK Guide는 법도 계약 조건도 아니다. 그런데도 판본을 알아야 하는 이유 셋
\item ① 사내 PMO 표준·양식·직무기술서가 인용하는 \textbf{조항 번호}가 살아 있는가, 죽어 있는가
\item ② 계약서·RFP의 "PMBOK 준거" 문구 — 판본과 문서를 특정하지 않으면 \textbf{분쟁의 씨앗}
\item ③ 표준을 만드는 조직이 5년 사이에 자기 구조를 \textbf{두 번 뒤집었다}는 사실 자체가 이 분야의 현재 상태
\end{itemize}
\keymsg{2021 삭제 \ra{} 2022 보완 별책 \ra{} 2025 회귀 — 세 줄을 하나의 이야기로}
\note{§2.1 "왜 계보를 알아야 하는가". 세 이유 중 첫째는 오늘 오후의 재매핑 대장으로, 둘째는 §2.4 ISO 21502(계약의 언어)로, 셋째는 다음 연표로 이어진다. 다음 슬라이드의 연표를 읽어 주되 2021\ra{}2022\ra{}2025의 세 줄(프로세스 삭제 \ra{} PMI 스스로 Process Groups 별책으로 보완 \ra{} 8판에서 40 프로세스 회귀)을 하나의 이야기로 묶는다. 예상 소요 3분.}
\end{frame}

\begin{frame}{PMBOK 계보와 주변 표준 연표 1996\til{}2026}
\fig[0.7]{fig_2_1_pmbok_timeline}
\figcap{그림 2-1. 1996 9 KA \ra{} 2017 6판 49 프로세스 \ra{} 2021 7판 12 원칙 \ra{} 2025-11 8판 40 프로세스}
\keymsg{APM BoK 7판(2019)이 PMI보다 먼저였고, 뿌리는 2006 Rethinking PM까지 올라간다}
\note{§2.1 연표. 1996 1판 계열 9 Knowledge Areas / 2017 6판 — 10 KA · 5 Process Groups · 49 processes · ITTO, Agile Practice Guide(Agile Alliance 공동) / 2019 APM BoK 7판 — PMI보다 먼저 원칙·맥락 중심 재편 / 2020 ISO 21502 / 2021 7판 — 12 principles · 8 performance domains, 프로세스·ITTO 전면 삭제 \ra{} PMIstandards+ / 2022 PMI 「Process Groups: A Practice Guide」 — 7판의 공백을 PMI 스스로 보완 / 2023 PRINCE2 7 — 7 principles · 7 practices · 7 processes / 2025-11 8판(디지털) — 6 principles · 7 domains · 5 Focus Areas · 40 processes, 2026-01 인쇄본 / 2026-07-09 PMP 신시험 시행(근거는 ECO, PMBOK은 참고 문헌). APM BoK 7판이 먼저였고 뿌리가 2006 Rethinking PM(마무리 §4.1)까지 올라간다는 점을 예고. 예상 소요 3분.}
\end{frame}

\begin{frame}{6판 (2017) — 아직 살아 있는 사내 표준의 원본}
\begin{itemize}
\item 10 지식영역 × 5 프로세스 그룹 = \textbf{49 프로세스}, 각각 ITTO(Inputs / Tools \& Techniques / Outputs)
\item 한국 기업에 뿌리내린 이유 ① \textbf{감사 가능성} — 감사·감리·검수는 "이 문서가 있는가"를 판정할 문장을 요구한다
\item 이유 ② \textbf{양식이 자산} — 수년간 축적된 산출물 서식. 원칙 12개로는 서식이 나오지 않는다
\item 이 교재는 6판을 "폐기된 판"으로 다루지 않는다 — 한국 기업 문서체계의 \textbf{원본(source document)}
\end{itemize}
\keymsg{8판이 40 프로세스를 되살린 것은 "프로세스 없는 표준은 조직이 쓸 수 없다"는 실무 피드백의 결과}
\note{§2.1 "6판". 예: 4.1 Develop Project Charter — 입력: 비즈니스 문서(비즈니스 케이스·편익 관리 계획), 협약, EEF, OPA / 출력: 헌장, 가정 기록부. 49개 프로세스의 출력물 목록이 정확히 감사가 요구하는 판정 문장이다. 2026년 현재도 한국 대기업·공공 PMO 사내 표준은 대부분 6판 ITTO가 뼈대라는 현실을 정상 상태로 말한다. 8판의 회귀는 6판이 옳았다는 뜻이 아니다. 수강생에게 "여러분 회사 표준 문서가 인용하는 판본은?"을 묻는다. 예상 소요 5분.}
\end{frame}

\begin{frame}{7판 (2021) 실험 \ra{} 8판 (2025) 부분 회귀}
\begin{columns}[T]
\begin{column}{0.48\textwidth}
\subhead{7판 — 12 원칙 / 8 성과영역, 프로세스 삭제}
\begin{itemize}
\item 두 문서가 한 권 \ra{} "7판 준거"는 어느 문서인지 불명확
\item 원칙 12개는 \textbf{감사할 수 없는 문장}
\item 실무 콘텐츠는 PMIstandards+(구독형)로
\end{itemize}
\end{column}
\begin{column}{0.5\textwidth}
\subhead{8판 — 6 원칙 + 7 성과영역 + 5 Focus Areas·40 프로세스}
\begin{itemize}
\item Focus Areas = Initiating / Planning / Executing / M\&C / Closing — \textbf{5 Process Groups가 이름을 바꿔 돌아왔다}
\item 실질 변경: \textbf{Integration \ra{} Governance}, \textbf{Cost \ra{} Finance}
\item Quality는 QA\ra{}Governance / QC\ra{}Scope로 분산, Communications는 Stakeholders에 흡수
\end{itemize}
\end{column}
\end{columns}
\keymsg{7판은 학술 논쟁의 20년 늦은 반영, 8판은 그것이 실무와 충돌한 뒤의 재타협}
\note{§2.1 "7판"과 "8판". 7판은 ANSI 표준 『The Standard for PM』과 『PMBOK Guide』가 한 권이어서 "7판 준거"가 어느 문서인지 불명확했고, "성실하고 존중하며 배려하는 관리자가 되라"를 참/거짓으로 판정할 수 없으므로 ISO 9001·CMMI·감리·위임전결 어디와도 연결되지 않았다. 비회원 직원은 PMIstandards+에 접근할 수 없었다. 8판의 40 프로세스는 non-prescriptive이며 예측형·애자일·하이브리드 모두에 적용된다. Procurement도 독립 도메인이 아니다. 온담 CFO의 "1차연도 148억, 자산 지출 일부를 리스·운영비로" 요구가 Cost가 아니라 Finance의 질문이라는 예시. Quality·Communications·Procurement의 정확한 배치는 인쇄본으로 확인하라고 단서. 예상 소요 6분.}
\end{frame}

\begin{frame}{PMBOK 8판 구조 — 6 원칙 / 7 도메인 / 5 Focus Areas · 40 프로세스}
\fig[0.64]{fig_2_2_pmbok8_structure}
\figcap{그림 2-2. Governance · Scope · Schedule · Finance · Stakeholders · Resources · Risk}
\keymsg{Governance가 첫 번째 도메인 — 실패의 가장 흔한 원인은 기술이 아니라 "누가 결정하는가"의 부재}
\note{§2.2 "7 performance domains"와 "도메인 배치의 함의". 도메인별 출처와 온담 P1 대응 — Governance(6판 Integration 승격·교체: 결정 권한 구조, 변경 통제, QA; 온담 3층 변경 통제와 게이트 G0\til{}G5) / Scope(그대로: 요구 1,180건 = StRS 214 · SyRS 396 · SRS 570, 제외 목록) / Schedule(7판 Planning·Measurement에서 재집결: 스프린트 36회, G1\til{}G4, 컷오버 창) / Finance(6판 Cost 개칭·확장 — 원가 + 자금 조달·현금흐름: 168 / 156 / 145.8억 세 층, 1차연도 148억, 5,000만 원 CFO 결재) / Stakeholders(6판 Communications 흡수: 3교시 §3.3) / Resources(7판 Team에서 이름 복귀: 피크 82명, IT본부 47명 중 14명 투입의 운영 부담) / Risk(7판 Uncertainty에서 이름 복귀: 5교시 §3.5). 이름이 6판 지식영역풍으로 돌아왔으므로 재매핑이 생각보다 쉽다 — 단 Quality·Communications·Procurement의 소멸과 Governance·Finance의 등장은 조항 번호 교체로 끝나지 않는다. "여러분 회사 품질보증팀은 QA(Governance)인가 QC(Scope)인가"(생각해 볼 질문 5). 예상 소요 5분.}
\end{frame}

\begin{frame}{한국 사내 표준은 6판에 머물러 있다 — 함의와 처방}
\begin{itemize}
\item 여러분 다수는 \textbf{두 판본 뒤처진 사내 표준}으로 일하며 \textbf{한 판본 뒤처진 자격 지식}을 갖고 있다 — 정상 상태
\item 함의 ①: "6판 4.6 Perform Integrated Change Control" 같은 조항 번호는 현행에 없다 \ra{} "현행 표준 준거" 요구에 방어 불가
\item 처방은 재작성이 아니라 \textbf{재매핑 대장} 네 열
\end{itemize}
\begin{tbl}[\footnotesize]{L{0.2\textwidth}L{0.1\textwidth}L{0.24\textwidth}L{0.34\textwidth}}
\textbf{사내 문서·조항} & \textbf{6판 근거} & \textbf{8판 대응 위치} & \textbf{조치} \\ \midrule
통합변경통제 절차서 & 4.6 & Governance 도메인 & 조항 번호만 교체, 절차 본문 유지 \\
원가관리계획서 & 7.1 & Finance 도메인 & 명칭 변경, 자금 조달·현금흐름 항목 신설 검토 \\
품질보증 절차서 & 8.x & QA\ra{}Governance / QC\ra{}Scope & 조직마다 결정 필요 \\
\end{tbl}
\keymsg{함의 ② PMP 승진 요건의 형평 문제 · 함의 ③ 판본이 바뀌어도 살아남는 산출물}
\note{§2.1 "한국 기업 사내 표준이 6판에 머물러 있다는 현실"과 "실무에서 어떻게 틀리는가". 함의 ②: PMP를 승진 요건으로 건 회사는 시험 개정 전후 형평 문제가 생긴다(인사의 문제이지만 PM이 먼저 알아야). 함의 ③(가장 중요): 판본이 바뀌어도 살아남는 산출물 — 헌장, 범위·일정·원가 기준선, 리스크 등록부, 변경 기록, 이해관계자 등록부. 이 소수를 중심으로 사내 표준을 다시 보는 것이 Day4 과제. 오류 사례: 7판이 나왔다고 프로세스·산출물 목록을 지운 조직이 6개월 뒤 감사에서 "판정 기준 없음" 지적으로 원복. "8판이 테일러링 원칙을 버렸다"도 오류 — 7판 \#7이 8판에서는 각 도메인 안에 분산 배치. 생각해 볼 질문 1. 예상 소요 5분.}
\end{frame}

\begin{frame}{6판 10 지식영역 \ra{} 8판 7 도메인 재매핑}
\fig[0.66]{fig_2_2_domain_mapping}
\figcap{그림 2-3. 실선 = 이름만 바뀜 / 파선 = 분산·흡수 (Quality · Communications · Procurement)}
\keymsg{워크북 §2 재매핑 대장의 원본 — 조항 번호 교체로 끝나는 행과 끝나지 않는 행을 구분하라}
\note{§2.2 재매핑. 그림에서 실선(Scope · Schedule · Stakeholders · Resources · Risk · Integration \ra{} Governance · Cost \ra{} Finance)은 조항 번호 교체로 대응되고, 파선(Quality가 QA/QC로 갈라지고 Communications가 Stakeholders에, Procurement가 독립 도메인에서 사라짐)은 조직이 결정해야 하는 행이다. 재매핑 대장은 이 그림을 자기 회사 문서 목록에 대고 한 행씩 채우는 작업이다. 예상 소요 3분.}
\end{frame}

\begin{frame}{PMBOK 8판의 6 core principles}
\begin{tbl}[\footnotesize\renewcommand{\arraystretch}{1.0}]{L{0.23\textwidth}L{0.22\textwidth}L{0.47\textwidth}}
\textbf{원칙} & \textbf{7판에서} & \textbf{온담에서 보이는 모습} \\ \midrule
1 \textbf{총체적 관점} (holistic view) & \#5 시스템 상호작용 + \#9 복잡성 & "P2가 늦으면 P1 컷오버가 밀린다"를 리스크 등록부에 \\
2 \textbf{가치 집중} (focus on value) & \#4 유지 + \#3·\#11·\#12 흡수 & 편익 214억의 근거를 항목별로 묻는 CFO = 이 원칙의 실행 \\
3 \textbf{품질 내재화} (embed quality) & \#8 "Build" \ra{} "Embed" & 사후 검사가 아니라 프로세스·산출물 안에. Quality 도메인 소멸과 짝 \\
4 \textbf{책임 있는 리더십} (lead accountably) & \#1 stewardship + \#6 병합 & CCB "반대의견 실명 기재" vs 2022 WMS 회의록의 빈 서명란 \\
5 \textbf{지속가능성 통합} & \textbf{신설} — 대응물 없음 & P2 물류 인력 246명 문제 = 사회적 지속가능성이 리스크로 \\
6 \textbf{권한 위임 문화} (empowered teams) & \#2 협력적 팀 환경의 진화 & 한국에서 이상과 현실의 거리가 가장 먼 원칙 \ra{} 5교시 \\
\end{tbl}
\keymsg{원칙 6(권한 위임)은 PRINCE2의 Manage by Exception·허용범위로 제도화된다}
\note{§2.2 "6 core principles". 원칙의 영어 문구는 출처에 따라 축약형·완전형이 혼재하므로 인쇄본 기준이라고 단서. 원칙 2 — 편익 214억 산출 근거를 항목별로 요구하는 CFO의 질문은 PM의 적이 아니라 이 원칙의 실행이다. 원칙 4 — CCB 규정의 "반대의견 실명 기재"와 2022 WMS 회의록의 비어 있는 물류본부장 서명란을 대비시킨다. 원칙 6(권한 위임)이 PRINCE2의 Manage by Exception·허용범위로 제도화된다는 연결을 미리 걸어 둔다. 예상 소요 6분.}
\end{frame}

\begin{frame}{PRINCE2 7 (2023) — 구조와 PMBOK과의 결정적 차이}
\begin{columns}[T]
\begin{column}{0.46\textwidth}
\fig[0.6]{fig_2_3_prince2}
\figcap{그림 2-4. PRINCE2 7 삼중 구조 + Project Board}
\end{column}
\begin{column}{0.52\textwidth}
\begin{itemize}
\item 7 principles 중 핵심: continued business justification · \textbf{manage by exception} · focus on products
\item 7판 변경: 통합 요소 4\ra{}5(People), 성과목표축 6\ra{}7(sustainability)
\item \textbf{PMBOK = "무엇을 알아야 하는가"의 지식체계 / PRINCE2 = "누가 무엇을 결정하는가"의 방법론}
\end{itemize}
\end{column}
\end{columns}
\keymsg{네 장치 — Project Board / Manage by Exception / Product-Based Planning / 살아 있는 Business Case}
\note{§2.3 "구조". 7 principles — continued business justification / learn from experience / defined roles, responsibilities and relationships / manage by stages / manage by exception / focus on products / tailor to suit the project. 7 practices(구 themes) — business case / organizing / plans / quality / risk / issues / progress. 7 processes — starting up / directing / initiating / controlling a stage / managing product delivery / managing a stage boundary / closing. 관리 접근법은 9개. 소유자가 PeopleCert(구 AXELOS)라는 점, PRINCE2 자격의 한국 시장 가치는 낮지만 네 장치는 자격과 무관하게 이식 가능하다는 점을 먼저 말해 자격 이야기로 새지 않게 한다. 예상 소요 4분.}
\end{frame}

\begin{frame}{장치 ① Project Board 3역할 · ② Manage by Exception}
\begin{tbl}[\footnotesize]{L{0.17\textwidth}L{0.52\textwidth}L{0.22\textwidth}}
\textbf{역할} & \textbf{책임} & \textbf{온담 P1} \\ \midrule
\textbf{Executive} & 비즈니스 케이스 — "돈을 쓸 가치가 있는가" & 정미란 CFO \\
\textbf{Senior User} & 산출물이 \textbf{사용되어 편익이 나오는 것}. 요구사항 대표, 편익 실현 책무 & 오정근 영업본부장 \\
\textbf{Senior Supplier} & 산출물을 만드는 쪽의 자원·기술적 실현 가능성 & 박세연 IT본부장 \\
\end{tbl}
\begin{itemize}
\item 이해가 상충하는 당사자를 \textbf{강제로 한 의사결정체에} 넣는다
\item MbE: 7 성과목표축 × 3계층의 \textbf{허용범위 격자} — 벗어날 것으로 \textbf{예측되는 순간} Exception Report
\end{itemize}
\keymsg{Exception Report는 문제 보고서가 아니라 권한 반납 문서 — PM의 자율성을 규정으로 방어하는 유일한 수단}
\note{§2.3 앞 두 장치. 한국 프로젝트가 표류하는 이유는 스폰서가 없어서가 아니라 사용자 대표와 공급자 대표가 같은 테이블에 앉지 않아서다 — 온담 G0 착수 조건에 "3역할 실명 지명"이 들어 있는 이유. Manage by Exception은 허용범위 안에서는 보고하지 않을 권리다: 7 성과목표축(benefits, cost, time, quality, scope, sustainability, risk) × 3계층(project / stage / work package). 한국 조직에 주는 실익은 보고 규율 강화가 아니라 그 반대 — 주간·수시보고 문화의 비용을 줄이는 것. 허용범위 초안(원가 +6\% = 10.08억 / 일정 +15영업일 / 범위 ±35건 / EMV 8억 / 편익 $-$5\%p)은 5교시 §3.4에서 표로 다시 본다. "운영위 명단에 3역할을 배정해 보라 — 배정 안 되는 인물(대개 절반)과 비어 있는 역할(대개 Senior Supplier)이 거버넌스 결함"(생각해 볼 질문 2). 예상 소요 7분.}
\end{frame}

\begin{frame}{장치 ③ Product-Based Planning · ④ 살아 있는 Business Case}
\subhead{Product-Based Planning — 활동이 아니라 산출물부터}
\begin{itemize}
\item Project Product Description \ra{} PBS \ra{} Product Descriptions(\textbf{품질 기준·검수 방법·검수 책임자} 강제) \ra{} Product Flow Diagram
\item 활동 중심 WBS("설계하기")는 검수 기준을 계획에서 누락시킨다 — 검수 분쟁의 원인 대부분
\end{itemize}
\subhead{Business Case — 승인용 문서가 아니라 단계 경계마다 갱신}
\begin{itemize}
\item 각 stage 경계에서 "계속할 이유가 아직 있는가"를 Board가 되묻는다
\item 편익 검토는 종료 \textbf{이후}에 예약 — 온담 이사회 조건 3항과 G5(2028-03-31)
\end{itemize}
\keymsg{PBS 노드는 명사여야 하고, 명사가 아니면 검수할 수 없다 \ra{} Day2 WBS}
\note{§2.3 뒤 두 장치, "한국에서 PRINCE2가 실무에 주는 것", "실무에서 어떻게 틀리는가". 편익 실현 책무는 Senior User, 정당성 유지는 Executive. 실무에서 틀리는 넷 — 허용범위 없이 Exception Report 양식만 도입 / 허용범위를 전 프로젝트 공통 ±10\%로(편익 민감도에서 역산해야) / PBS를 동사로 채움 / PMBOK CCB(권한 회수)와 PRINCE2 Change Authority(권한 위임)를 섞음 — 온담 3층 변경 통제가 의도적으로 섞어 놓았고 Day3에서 그 비용을 확인한다. 온담의 "본부장 위임전결 5억"과 "원가 허용범위 10.08억"이 서로 맞지 않는다는 불일치를 찾는 것이 격자의 첫 용도. 해외 발주처 대응에서 Highlight / End Stage / Exception Report, PID라는 이름을 아는 것의 값. 예상 소요 6분.}
\end{frame}

\begin{frame}{ISO 21502:2020 — 계약서가 인용하는 표준}
\begin{itemize}
\item 『Guidance on project management』, ISO/TC 258. \textbf{자격시험이 없다} \ra{} 교육시장에서 무명
\item 글로벌 EPC·발주처 계약서와 RFP에서는 "PMBOK 준거"보다 자주 — \textbf{중립성}
\item 구조: §4 맥락 \ra{} §5 전제 조건 \ra{} §6 통합 실무(\textbf{§6.2 pre-project} … \textbf{§6.9 post-project}) \ra{} §7 관리 실무 17종
\item \textbf{ISO 21500의 함정}: 2021년판 제목은 "Context and concepts" \ra{} "ISO 21500에 따라 관리한다"는 \textbf{무의미한 문장}
\item 지침은 \textbf{21502}, 용어는 \textbf{21506:2024}, 거버넌스는 \textbf{21505:2017}
\end{itemize}
\keymsg{계약에는 "ISO 21502:2020 §7 관리 실무에 따라"처럼 문서·연도·조항을 특정하라}
\note{§2.4 전체. 발주처 PMI, 시공사 PRINCE2, 본사 자체 방법론일 때 유일하게 인용 가능한 표준이다. 프로젝트의 앞(§6.2)과 뒤(§6.9)를 표준 안에 넣었다 — Morris의 front-end 주장과 편익 관리를 표준 수준에서 수용. 2012년판 ISO 21500의 제목은 "Guidance on project management"였으나 2021년판은 "Context and concepts"로 바뀌었다. 계열: 21503(프로그램, 2022) / 21504(포트폴리오, 2022) / 21505:2017 거버넌스 / 21511:2018 WBS / 21508:2018 EVM(2판 2026 예정) / 21512:2024 EVM 구현. 21505 거버넌스는 이사회·오너 계층 책무를 다루는 유일한 ISO 문서 — PRINCE2 Project Board를 "표준 근거 있는 것"으로 정당화할 때 인용하면 경영진 설득력이 달라진다. 21502는 계약과 감리의 언어이지 팀의 언어가 아니다. 예상 소요 7분.}
\end{frame}

\begin{frame}{표준 지형도 — 지식체계 / 방법론 × 거버넌스 층 / 인도 층}
\fig[0.62]{fig_2_4_standards_map}
\figcap{그림 2-5. PMBOK · PRINCE2 · ISO 21502/21505 · 애자일 프레임워크가 각각 어느 층에 있는가}
\keymsg{같은 층에 있는 것끼리만 대체된다 — 애자일은 인도 층, PRINCE2는 거버넌스 층}
\note{§2.4\til{}2.5의 지도(개념도 — 위치는 상대적 성격). 가로축은 지식체계·지침("무엇을 알아야 하는가": ISO 21502·21505, PMBOK 8판, Agile Practice Guide, Disciplined Agile)과 방법론·프레임워크("누가 무엇을 결정하는가": PRINCE2 7, PRINCE2 Agile, SAFe, Scrum·Kanban·XP), 세로축은 계약·거버넌스 층(계약서·이사회·오너·게이트)과 인도 팀 층(팀이 일하는 리듬). PMBOK 8판은 두 층에 걸쳐 있고 예측형·애자일·하이브리드 모두에 적용된다. Scrum·Kanban·SAFe는 인도 층에만 있으므로 PMBOK·PRINCE2·ISO의 대체물이 아니라 그 아래에서 조합되는 개발 접근이다 — 다음 슬라이드의 논지를 그림 하나로 미리 보여 준다. 그림 오른쪽의 메모("PRINCE2 stage 경계 = 승인 게이트, SAFe PI 경계 = 정렬 리듬, 1:1 매핑 금지")가 다음 슬라이드의 핵심 한 줄이다. 한국 조직이 "표준을 바꾼다"고 할 때 실제로는 어느 층을 바꾸려는 것인지 묻는 데 쓴다. 예상 소요 2분.}
\end{frame}

\begin{frame}{애자일·SAFe·하이브리드의 위치 — 표준의 반대편이 아니다}
\begin{itemize}
\item 애자일(Scrum, Kanban / SAFe, LeSS, DA)은 표준이 \textbf{개발 접근(development approach)}이라 부르는 것의 한 종류
\item \textbf{Agile Practice Guide (2017)} — Agile Alliance 공동 저작. 실무 자산은 Suitability Filter
\item \textbf{PRINCE2 Agile} — 게이트 거버넌스 유지, \textbf{인도 계층만 애자일화}. 시간·비용 고정, 범위는 MoSCoW
\item \textbf{SAFe} — 한국 실패 원인은 Personnel·Culture 축 무시. \textbf{PI ≠ Stage} — 1:1 매핑하면 중단 결정 지점이 사라진다
\item \textbf{온담 P1 하이브리드} — 스프린트 36회, 요구사항 월 1회 동결, FP 12,400점 고정가, 게이트 G0\til{}G4
\end{itemize}
\keymsg{하이브리드 = 어느 결정을 언제까지 되돌릴 수 있게 남기고 언제 동결하는가의 설계}
\note{§2.5 전체. Agile Practice Guide의 Suitability Filter는 문화·팀·프로젝트 3범주 방사형 그림. PRINCE2 Agile은 SAFe와 반대 방향 — 한국 대기업이 게이트를 버릴 수 없을 때 유일한 표준급 절충. Disciplined Agile은 PMI가 2019 인수, process goal + decision points + options로 테일러링을 선택지 카탈로그로 조작화. SAFe에서 stage 경계는 "계속할 것인가"의 승인 게이트, PI 경계는 팀 간 정렬 리듬 — 1:1로 매핑해 End Stage Report·갱신된 Business Case를 없애면 중단 결정 지점이 조직에서 사라진다. 온담 P1은 사업 정당성 재확인을 게이트(G1·G2·G3)에, 팀 간 정렬을 스프린트 경계와 월 동결에 둔다. 오류 둘: "애자일이 waterfall을 대체했다"는 교체 서사(Lenfle \& Loch 2010 — 맨해튼·폴라리스는 원래 병렬 실험이었고 표준화 과정에서 통제만 남았다; 애자일은 신문물이 아니라 잃어버린 뿌리) / 하이브리드를 "스프린트 돌리면서 간트도 그리는 것"으로 이해. 예상 소요 8분.}
\end{frame}

\begin{frame}{테일러링의 기준은 하나 — 되돌림 비용}
\begin{columns}[T]
\begin{column}{0.46\textwidth}
\fig[0.62]{fig_2_6_reversal_axis}
\figcap{그림 2-6. 되돌림 비용 축 위의 P1 vs P2}
\end{column}
\begin{column}{0.52\textwidth}
\begin{itemize}
\item \textbf{P1}: 스프린트 1회 ≈ 1.41억 — 잘못된 결정 하나 = 1.4\til{}2.8억 \ra{} \textbf{실험 가능}
\item \textbf{P2}: 설비 발주 후 선급금 \textbf{29.6억} 잠김, 리드타임 34주 \ra{} \textbf{예측형}, 앞단에 시간
\item P1의 예외 — G2 전환 정본(3.8억·6주), G3 컷오버(34시간) 앞은 예측형
\end{itemize}
\end{column}
\end{columns}
\keymsg{프로그램 하나에 방법론 하나를 강제하면 안 된다 — 다섯 프로젝트, 다섯 리듬}
\note{§2.6 "개념"과 "되돌림 비용이 방법을 결정한다". 표준 셋(PMBOK 7·8, PRINCE2 7, ISO 21502) 모두 테일러링을 요구하지만 판단 도구를 주지 않는다. 가장 단순하고 강력한 기준은 결정을 되돌리는 데 얼마가 드는가 — 적응형과 예측형은 이념이 아니라 되돌림 비용에 대한 두 개의 합리적 반응이다. P1 스프린트 1회 ≈ 1.41억(44명 × 0.5개월 × 640만 원). P2는 리드타임 34주(41주 가능)를 압축할 수 없고 설비 벤더는 2주마다 인도하지 않으므로 스프린트가 무의미하며, 처리량 4,200 검증·사양 동결 같은 앞단에 시간을 써야 한다. P1의 되돌릴 수 없는 두 지점 앞에는 point of no return을 명시한다. P4는 기한 고정, P5는 SAC. 프로그램 운영위의 일은 이 다섯 리듬을 통역하는 것. 1교시 §1.5 표의 "되돌림 비용" 열이 여기서 테일러링 규칙으로 승격된다. 생각해 볼 질문 3. 예상 소요 7분.}
\end{frame}

\begin{frame}{Boehm-Turner 5축 (2003) — 축이 갈리면 시스템을 쪼개라}
\begin{columns}[T]
\begin{column}{0.46\textwidth}
\fig[0.6]{fig_2_6_boehm_turner}
\figcap{그림 2-7. 5축 극좌표 — P1 전체와 SRS/SyRS 분할}
\end{column}
\begin{column}{0.52\textwidth}
\begin{itemize}
\item P1: Size 82명 \ra{} 바깥 / Criticality 615개 매장 결제 / Dynamism 월 동결 \ra{} 중간 / Culture 질서 쪽
\item 처방은 "한 방법을 골라라"가 아니라 \textbf{축이 갈리면 쪼개서 부분마다 다른 방법}
\item 매장 화면·가맹 포털(SRS 570) \ra{} \textbf{더 애자일} / 인터페이스·데이터·결제(SyRS 396) \ra{} \textbf{더 규율}
\end{itemize}
\end{column}
\end{columns}
\keymsg{한국에서 결정적인 축은 Personnel과 Culture — 프로세스만 이식하면 실패한다}
\note{§2.6 "Boehm-Turner 5축". 다섯 축: Size(소규모\til{}대규모) / Criticality(편의 손실\til{}생명·막대한 자금 손실) / Dynamism(월 50\%\til{}1\% 변화) / Personnel(고숙련\til{}저숙련 비율) / Culture(자율로 번성\til{}질서로 번성). P1의 Personnel은 수행사 61명 분포 미상이라 워크북에서 가정한다. Culture는 기관투자자 시계, 5,000만 원 CFO 결재로 질서 쪽. 분할 처방은 영업본부 vs IT본부 우선순위 긴장을 방법론 수준에서 해소하는 한 답이다. 2003년 문헌인데 2026년 SAFe 실패 분석에 가장 잘 맞는다. 5축 결과를 "애자일 점수"로 환산해 조직 전체에 하나를 강제하는 것은 원저의 결론과 정반대. 워크북 §3 심화 과제(SRS/SyRS 분할 근거)가 이 슬라이드다. 생각해 볼 질문 4. 예상 소요 6분.}
\end{frame}

\begin{frame}{Cynefin — 방법론 선택 도구가 아니다 (Snowden \& Boone, 2007 HBR)}
\begin{columns}[T]
\begin{column}{0.44\textwidth}
\fig[0.6]{fig_2_6_cynefin}
\figcap{그림 2-8. 네 도메인과 의사결정 순서}
\end{column}
\begin{column}{0.54\textwidth}
\begin{itemize}
\item 오용 "complex하니까 애자일" — 처방은 \textbf{의사결정 양식의 전환}이고, \textbf{도메인 판정은 사후적}
\item 프로젝트 전체가 아니라 \textbf{결정 하나하나}에 — SAP 애드온 = complicated / 가맹점주 반응 = complex \ra{} 파일럿 12개점 / OD-POS 보안 사고 = chaotic
\item 시중의 Stacey matrix 2×2는 \textbf{Stacey의 원 도식이 아니라 후대 각색본} — ``만든 사람이 버린 그림''이라고 말하지 말 것
\end{itemize}
\end{column}
\end{columns}
\keymsg{하나의 프로젝트 안에 네 도메인이 공존한다 — 결정마다 도메인을 감각하라}
\note{§2.6 "Cynefin". 네 도메인 — Clear(인과 명백, sense\ra{}categorize\ra{}respond, best practice) / Complicated(분석 필요, sense\ra{}analyze\ra{}respond, 전문가 분석) / Complex(사후에만 알 수 있음, probe\ra{}sense\ra{}respond, 안전한 실험과 창발) / Chaotic(인과 없음, act\ra{}sense\ra{}respond, 일단 행동해 안정화). 사전에 "이 프로젝트는 complex"라고 라벨 붙이는 것은 프레임워크의 논지에 어긋난다. Stacey matrix를 쓸 거라면 Ralph Stacey 본인이 후속 저작에서 폐기했다는 사실을 먼저 말해야 한다 — 경력 있는 청중 앞에서 이 한마디가 신뢰를 만든다. 예상 소요 5분.}
\end{frame}

\begin{frame}{테일러링을 문서로 남기기 — TDR 네 칸}
\begin{itemize}
\item 세 도구는 \textbf{다른 질문}에 답한다: NTCP = 유형 / Boehm-Turner = 준비도 / Cynefin = 개별 결정 — 충돌 지점이 설계할 지점
\item 근거 기록이 없으면 모든 생략이 감사에서 "표준 미준수" \ra{} 방어용 문서 양산 — \textbf{한국 조직에서 문서가 줄지 않는 진짜 메커니즘}
\end{itemize}
\begin{tbl}[\footnotesize]{L{0.26\textwidth}L{0.24\textwidth}L{0.22\textwidth}L{0.12\textwidth}}
\textbf{무엇을} & \textbf{왜} & \textbf{어떤 리스크를 수용하는가} & \textbf{누가 승인} \\ \midrule
스프린트 단위 상세 일정표 \textbf{생략} & 월 동결·게이트로 통제 & 스프린트 내 지연 가시성 저하 & 스폰서 \\
데이터 전환 리허설 3회 \textbf{추가} & 2023 감사 지적 "형상관리 미비", G2 비가역성 & 시험 기간 압박 & 스폰서 \\
\end{tbl}
\keymsg{생략 항목별로 행이 있어야 감사 조항에 대응된다 — 테일러링은 줄이는 것만이 아니다}
\note{§2.6 "테일러링을 문서로 남기기"와 "실무에서 어떻게 틀리는가". Tailoring Decision Record는 생략·조정한 프로세스·산출물마다 한 행이다. 서술 문단으로 처리하면 안 된다. 테일러링은 줄이는 것만이 아니다 — P1의 리허설 3회, 외부 감리 4.2억, 발주사 PMO 3명은 표준이 아니라 온담의 과거 실패가 요구하는 추가다. 워크북 §3에서 P1의 TDR을 작성한다고 되어 있으나 이 워크북 판본의 §3은 이해관계자 등록부이므로, TDR은 Day2 워크북 또는 심화 과제로 안내한다. 예상 소요 5분.}
\end{frame}

\begin{frame}{자격 체계의 실효성 — 문을 여는 열쇠이지 판단을 대신하지 않는다}
\begin{tbl}[\footnotesize]{L{0.17\textwidth}L{0.24\textwidth}L{0.52\textwidth}}
\textbf{자격} & \textbf{기반 문서} & \textbf{실효} \\ \midrule
\textbf{PMP} (PMI) & \textbf{ECO} — PMBOK은 참고 문헌. 2026-07-09 신버전 & 한국 최고 인지도. 증명하는 것은 "지식과 경력 시간"이지 "상황 속 판단"이 아님 \\
PMI-ACP & 애자일 실무 & \\
PRINCE2 7 F/P (PeopleCert) & PRINCE2 7 매뉴얼 & 한국 자격 가치 낮음. 해외 발주처·글로벌 본사 보고에서 "문서 이름 맞추는 능력" \\
SAFe SA/SSM/SPC & SAFe & 전환 컨설팅 수임에 유리, 연 1회 갱신 \\
IPMA A\til{}D & \textbf{ICB4} 29 역량 요소 & 자격 실효는 낮으나 ICB4는 사내 PM 등급제·직무기술서의 원천 \\
\end{tbl}
\keymsg{"PMP를 8판으로 다시 봐야 하나요?" \ra{} 아니오. 시험 근거는 ECO. 판본 논쟁은 사내 문서 체계의 문제}
\note{§2.7 전체. PMP는 공공 정보화 가점, 대기업 PMO 채용 요건. SAFe는 개인 급여 프리미엄이 PMP보다 낮은 경향. 사내 교육을 PMBOK 목차에 맞추면 합격률은 오히려 역효과. 자격을 "역량 증명"으로 승진에 쓰면 자격은 있으나 판단이 없는 PM을 승진시킨다. Crawford, Morris, Thomas, Winter(2006) — 자격 기반 역량 모델은 trained technician을 측정하며 성과는 reflective practitioner의 판단에서 나온다(마무리 §4.1 방향 5). 반대로 자격 무용론도 오류 — 원 논문의 주장은 대체가 아니라 발달 경로이고, 자격은 공통 어휘를 얻는 가장 빠른 길. Business Environment 배점 상승은 2차 출처이므로 숫자 단정 금지. 예상 소요 5분.}
\end{frame}

\begin{frame}{2교시 핵심 정리}
\begin{columns}[T]
\begin{column}{0.5\textwidth}
\begin{enumerate}
\item PMBOK 6판(49 프로세스) \ra{} 7판(12원칙, 삭제) \ra{} 8판(6원칙/7도메인/40 프로세스, \textbf{부분 회귀})
\item 사내 표준은 대부분 6판 — 처방은 \textbf{4열 재매핑 대장}. Integration\ra{}Governance, Cost\ra{}Finance
\item PRINCE2의 자산 = \textbf{3역할 / 허용범위 격자와 MbE / PBP / 갱신되는 BC}
\item 계약서는 ISO 21502 — "21500에 따라"는 틀린 문장. 용어는 21506
\end{enumerate}
\end{column}
\begin{column}{0.48\textwidth}
\begin{enumerate}\setcounter{enumi}{4}
\item 애자일은 개발 접근의 한 종류. PRINCE2 Agile은 게이트 유지. Stage ≠ PI
\item 테일러링의 기준은 \textbf{되돌림 비용}. NTCP / Boehm-Turner / Cynefin은 다른 질문. 결과는 \textbf{TDR 네 칸}
\item PMP의 근거는 ECO. 자격은 문을 여는 열쇠
\end{enumerate}
\end{column}
\end{columns}
\keymsg{점심시간에 워크북 §0.3 온담 한 페이지 요약과 §1.5·§2.5 완성 예시본을 훑어 올 것}
\note{제2장 핵심 정리 그대로. 회귀는 프로세스 없는 표준을 조직이 쓸 수 없다는 피드백의 결과(1). Quality·Communications·Procurement 소멸이 실질 변경(2). Boehm-Turner는 갈리면 분할(6). 점심 전 마지막 슬라이드 — 오후 4교시 실습을 위해 워크북 §0.3(온담 한 페이지 요약)과 §1.5·§2.5 완성 예시본을 점심시간에 훑어 오도록 안내한다. 예상 소요 3분.}
\end{frame}

% =====================================================================
\section{3교시 (90분) — 착수 이론: 비즈니스 케이스 · 편익의 소유 · 헌장 · 이해관계자}
% =====================================================================

\begin{frame}{이 교시에서 배우는 것}
\begin{itemize}
\item 착수 이전 — 비즈니스 케이스와 편익의 소유권 [§3.1]: 이미 결정된 것들 / BC의 항목 / Zwikael \& Smyrk의 ITO / Ward \& Daniel의 BDN
\item 프로젝트 헌장 — 항목별 의미와 한국 기업에서의 실질화 [§3.2]
\item 이해관계자 — 식별·분석·참여 계획 [§3.3]: Power/Interest의 한계 / Mitchell-Agle-Wood 현저성 모델 / 한국 특유의 이해관계자
\item 4교시 실습 ①(비즈니스 케이스 요약 + 헌장)의 이론 부분
\end{itemize}
\keymsg{오전보다 온담 숫자가 훨씬 많이 나온다 — 워크북 §0.3 한 페이지 요약을 펼쳐 놓고 듣는다}
\note{제3장 §3.1\til{}3.3. 시간 배분: §3.1 30분, §3.2 30분, §3.3 25분, 정리 5분. 오전보다 온담 숫자가 훨씬 많이 나온다 — 워크북 §0.3 한 페이지 요약을 펼쳐 놓고 듣게 한다. 예상 소요 2분.}
\end{frame}

\begin{frame}{착수의 첫 번째 일은 검증이다 — 표준이 "착수 전"을 수용하는 방식}
\begin{itemize}
\item ISO 21502 \textbf{§6.2 pre-project activities} / PRINCE2 \textbf{Starting Up a Project}가 Initiating 앞에
\item PMBOK 8판 Initiating은 프로세스 둘뿐(헌장 개발, 이해관계자 식별) — 입력인 BC·편익 관리 계획은 \textbf{프로젝트 바깥}에서 온다
\item P1 PM이 받는 것: 이사회 부의안(2025-12) + 편익 표(연 214억) + 프로그램 관리 규정 + 승인 조건 5개항
\item 착수의 첫 번째 일은 새로 만드는 것이 아니라 \textbf{이미 만들어진 BC의 숫자와 가정을 검증하는 것}
\end{itemize}
\keymsg{읽지 않고 헌장을 쓰는 것은 남이 정한 목표를 자기가 정한 것처럼 서명하는 일}
\note{§3.1 "착수 전에 이미 결정된 것들". 1교시 Engwall("이미 결정된 것들의 목록화")이 표준에서 어떻게 자리 잡았는지로 연결한다 — 세 표준 모두 프로젝트 이전 단계를 명시적으로 둔다. P1 PM이 착수 시점에 손에 쥐는 네 문서를 읽지 않고 헌장을 쓰는 것은 남이 정한 목표를 자기가 정한 것처럼 서명하는 일이다. 다음 슬라이드에서 BC 항목을 온담 P1로 채우며 검증한다. 예상 소요 4분.}
\end{frame}

\begin{frame}{비즈니스 케이스를 온담 P1로 채우면}
\begin{tbl}[\scriptsize]{L{0.11\textwidth}L{0.3\textwidth}L{0.5\textwidth}}
\textbf{항목} & \textbf{무엇을 적는가} & \textbf{온담 P1에서} \\ \midrule
\textbf{문제 / 기회} & 무엇이 잘못되어 얼마의 비용을 만드는가. 숫자로 & OD-POS OS 지원 종료(2024-07), 문서 부재 인터페이스 13본. 폐기율 3.9\% = 연 \textbf{98.7억}(2,530억 × 3.9\%), 결품률 4.7\%, 가맹 발주 리드타임 38시간 \\
\textbf{대안} & do-nothing / do-minimum / 제안안 / 그 밖 & (a) 현행 유지 — SAP ECC 2027-12 종료 시 이중 위기 (b) 부분 개선 (c) 통합 플랫폼(P1) (d) 상용 패키지 — \textbf{검토 기록 없음} \\
\textbf{편익} & 기준값·목표값·측정 방법·시점·\textbf{측정 책임자} & B-01\til{}B-08, M+24 연 214억. 측정 책임자 "프로그램 착수 후 지정 예정" — \textbf{결함 ①} \\
\textbf{비용} & 프로젝트 + 운영 = TCO & P1 168억 + 나머지 260억. 연 운영비 27억(3년 81억)이 편익에서 \textbf{미차감} — \textbf{결함 ②} \\
\textbf{리스크} & 편익이 안 나올 주요 이유 & 과거 2건 중단·축소, P2 리드타임, 3PL WMS 개방, 가맹 부속서 388점 서면 동의, 규제 기한 \\
\textbf{권고} & 어느 대안을 왜. 승인 조건 & 이사회 원칙 승인 + 5개 조건 \\
\end{tbl}
\keymsg{표를 채우면서 알게 되는 것 — 온담 BC는 존재하지만 완결되어 있지 않다}
\note{§3.1 "비즈니스 케이스에는 무엇이 들어가는가". 문제/기회 행은 재고 실시간 가시성 부재까지 포함한다. 대안 (a) 현행 유지는 패치 없는 결제 시스템 + SAP ECC 2027-12 종료 시 이중 위기다. 편익 행의 측정 책임자 "추후 지정"과 비용 행의 운영비 미차감이 두 결함이며, 다음 슬라이드에서 결함을 목록화한다. 예상 소요 5분.}
\end{frame}

\begin{frame}{온담 비즈니스 케이스의 결함 — PM이 고칠 것이 아니라 목록으로 만들 것}
\begin{itemize}
\item B-03 결품률: "기회손실, 미산정" / B-02 재고 정확도 91.3\%: 분모가 \textbf{순환실사 대상 SKU 69\%}뿐
\item B-06 물류 처리량 96억: 이사회 자료에 \textbf{분모가 없고}, 노조는 "인원 절반이라는 뜻이냐", 한동석은 "4,200은 현장 숫자가 아니다"
\item 편익 합계 214억에서 연 운영비 27억 미차감 \ra{} 정직한 숫자는 \textbf{187억}
\item 측정 책임자 "추후 지정" — 그것은 지정하지 않겠다는 뜻
\item CFO: "근거가 없는 항목은 승인하지 않겠다" / 내부감사: "편익 산정 근거를 문서로 제출하라"
\end{itemize}
\keymsg{PM의 착수 시점 일은 결함을 고치는 것(권한 밖)이 아니라 목록으로 만들어 스폰서에게 결정을 요구하는 것}
\note{§3.1의 결함 문단과 "실무에서 어떻게 틀리는가". 오류 넷: BC를 승인용으로 한 번 쓰고 서랍에(좀비 프로젝트) / 편익을 종료와 함께 폐기(측정 불가) / Gate 1에서 5년 매출·전체 NPV 요구(허구 \ra{} survival of the unfittest; 게이트는 "다음 단계로 갈 권리를 사는 것") / 운영비 미차감. 4교시 워크북 §1 심화 과제가 이 결함 목록 + 결정 요청 문안이다. 예상 소요 6분.}
\end{frame}

\begin{frame}{편익은 누가 소유하는가 — Zwikael \& Smyrk의 ITO 모델 (2011)}
\begin{columns}[T]
\begin{column}{0.55\textwidth}
\fig[0.66]{fig_3_1_ito}
\figcap{그림 3-1. Input \ra{} Output(PM) \ra{} Outcome(사업 오너) — 책임의 분리선}
\end{column}
\begin{column}{0.43\textwidth}
\begin{itemize}
\item \textbf{산출물} = PM의 책임 / \textbf{활용·성과} = 사업 오너의 책임
\item PM에게 매출·절감 목표를 주면 \textbf{범위 방어와 책임 회피가 합리적 행동}이 된다
\item 온담: 폐기율 43억의 오너는 \textbf{오정근}(Senior User), 물류 96억은 \textbf{한동석}, 인건비율 38억은 영업본부
\end{itemize}
\end{column}
\end{columns}
\keymsg{편익 측정 시점(G5 2028-03-31)에 프로젝트 조직은 없다 \ra{} 측정 책임자는 상설 조직의 실명}
\note{§3.1 "편익은 누가 소유하는가". 한국 기업에서 가장 자주 잘못 설계되는 지점은 benefit owner를 PM이나 PMO장으로 지정하는 것이다. PRINCE2가 편익 책무를 Senior User에게 두는 것도 같은 논리. 부서명을 쓰면 그 편익은 6개월 안에 유령이 된다. 1교시 임시조직론의 "제도화된 종료"가 왜 실명이어야 하는지의 근거다. "편익 오너의 이름이 없는 편익은 착수 시점에 이미 실현되지 않기로 결정된 편익이다." 예상 소요 6분.}
\end{frame}

\begin{frame}{편익 실현은 왜 대부분 실패하는가 — Ward \& Daniel의 BDN}
\fig[0.68]{fig_3_1_bdn}
\figcap{그림 3-2. 온담 편익의 Benefits Dependency Network — IT 조력자 \ra{} 조력 변화 \ra{} 사업 변화 \ra{} 편익}
\keymsg{P1이 만드는 것은 맨 왼쪽뿐 — "P1이 성공해도 편익이 안 나올 경로"가 착수 시점에 보여야 한다}
\note{§3.1 "편익 실현은 왜 대부분 실패하는가". 편익은 기술이 아니라 조직 변화에서 나오는데, 대부분의 프로젝트는 기술을 인도하고 끝난다. BDN은 투자 목적 \ra{} 편익 \ra{} 사업 변화 \ra{} 조력 변화 \ra{} IT 조력자를 역방향으로 연결한다. 폐기율 3.9\%\ra{}2.2\%(연 43억)는 실시간 재고 화면(P1) \ra{} 매장 직원이 화면을 쓸 줄 안다(P3 교육 3,900명) + 발주 정책 개정(영업본부 규정) \ra{} 매장 직원이 재고를 보고 발주량을 바꾼다는 사슬을 거친다. 나머지는 P3와 영업본부의 일이다. 이것이 헌장의 가정·제약 항목과 리스크 등록부의 재료다. Bradley 『Benefit Realisation Management』(2판, 2010) — benefit profile(편익마다 기준값·목표값·측정 방법·오너·시점 한 장), benefit owner 지정 절차. 워크북 §1 편익 표의 "분모·기여 구분·측정 책임자" 열이 BDN과 benefit profile에서 왔다. 예상 소요 5분.}
\end{frame}

\begin{frame}{프로젝트 헌장은 권한 문서다}
\begin{columns}[T]
\begin{column}{0.46\textwidth}
\fig[0.6]{fig_3_2_doc_flow}
\figcap{그림 3-3. 착수 문서 흐름 — BC \ra{} 헌장 \ra{} 계획서}
\end{column}
\begin{column}{0.52\textwidth}
\begin{itemize}
\item 헌장 = 프로젝트의 존재를 공식화하고 \textbf{PM에게 조직 자원을 쓸 권한을 부여}하는 문서
\item 답하는 질문은 "무엇을 어떻게"가 아니라 \textbf{"왜 하고, 누가 결정하며, PM은 무엇을 할 수 있는가"}
\item PMBOK 헌장 ≈ PRINCE2 \textbf{Project Brief}(\ra{} PID) ≈ ISO 21502 §6.4 "charter or equivalent"
\end{itemize}
\end{column}
\end{columns}
\keymsg{워크북 헌장 = Project Brief에 PID의 핵심(권한·허용범위·완료 기준)을 얹은 15개 항목}
\note{§3.2 "왜 쓰는가". 표준 위치: PMBOK 8판 Initiating 첫 프로세스 Develop Project Charter의 출력(스폰서 발행, PM 작성 지원) / PRINCE2 7 Project Brief(Starting Up) \ra{} PID(Initiating) / ISO 21502 §6.4. 워크북 §2.3에 15개 항목 전부의 "무엇을/어떻게/흔한 오류/온담에서는" 가이드가 있다고 안내한다. 이어지는 네 슬라이드는 그중 헌장의 성격을 결정하는 항목만 다룬다 — 상위 요구사항(5\til{}10줄, 상세가 아님), 산출물(명사로 — PBP의 출발), 마일스톤(되돌릴 수 없는 지점 G2·G3와 외부 고정 기한을 구분)은 구두로 짧게. 예상 소요 4분.}
\end{frame}

\begin{frame}{헌장 항목 ① 목적 / 정당성 · 측정 가능한 성공 기준}
\subhead{목적 — 3\til{}5문장, 반드시 숫자}
\begin{itemize}
\item 흔한 오류: \textbf{솔루션을 목적으로} — "차세대 플랫폼을 구축한다"는 목적이 아니다
\item 온담: "패치가 중단된 시스템으로 615개 매장의 결제를 처리하고, 연 98.7억어치가 폐기되며, 2027 IPO 전에 벗어나야 한다" — \textbf{프로젝트가 없어져도 남는 문장}
\end{itemize}
\subhead{성공 기준 — 두 층으로 분리}
\begin{itemize}
\item 프로젝트 관리 성공(PM 책임): 2027-05-29까지 156억 한도 내 다섯 산출물 인도, SAC 충족
\item 프로젝트 성공(편익 오너 책임): M+24 폐기율 2.2\%, 결품률 1.5\%, 리드타임 12h — \textbf{측정 책임자 실명}
\end{itemize}
\keymsg{흔한 오류 — 편익 목표를 PM의 성공 기준에 넣는 것}
\note{§3.2 "항목별 의미" 앞부분. 목적 항목은 사업 배경·해결할 문제·기대 효과를 3\til{}5문장으로, 반드시 숫자를 넣어 쓴다. 온담 원문: "보안 패치가 중단된 시스템으로 615개 매장의 결제를 처리하고 있고, 재고를 실시간으로 볼 수 없어 연 약 98.7억어치 식자재가 폐기되며, 2027 IPO 예비심사 전에 이 상태를 벗어나야 한다." 성공 기준의 두 층은 1교시 Cooke-Davies 구분 그대로이며, 편익 층에는 측정 책임자 실명이 붙어야 한다. 예상 소요 5분.}
\end{frame}

\begin{frame}{헌장 항목 ② 범위 / 제외 범위 — 무엇을 만들지 않는가}
\begin{itemize}
\item 무엇을 만드는가보다 \textbf{무엇을 만들지 않는가}가 더 중요 — 제외 범위 없는 헌장은 범위 팽창의 초대장
\item P1 제외: SAP ECC 교체 / 이천 설비·WMS(P2) / 교육·가맹 계약 개정(P3) / 규제 등록 자체(P4) / 운영 조직(P5)
\item 각 제외 항목에 \textbf{"그러면 누가 하는가"}를 붙인다
\item 흔한 오류: "기타 협의된 사항"처럼 모호하게
\end{itemize}
\keymsg{제외 범위는 포함 범위와 같은 무게, 같은 분량으로}
\note{§3.2 "항목별 의미" — 범위/제외 범위. P1 제외 목록의 각 항목은 프로그램의 다른 프로젝트가 맡는다는 점에서 "누가 하는가"의 답이 이미 있다. 워크북 §2.3에서는 여기에 가맹점 단말 하드웨어 교체를 굵게 추가한다 — 헌장에서 가장 논쟁적인 줄이고 그래서 가장 필요한 줄(4교시). 예상 소요 3분.}
\end{frame}

\begin{frame}{헌장 항목 ③ 예산 층 · 가정과 제약}
\subhead{예산 요약 — 하나의 숫자가 아니라 층}
\begin{tbl}[\small]{ccccc}
\textbf{승인 예산} & \textbf{집행 한도} & \textbf{계획 원가} & \textbf{PM 통제 컨틴전시} & \textbf{관리예비비 (스폰서)} \\ \midrule
168.0억 & 156.0억 & 145.8억 & 10.2억 & 12.0억 \\
\end{tbl}
\subhead{가정 = 참이라고 믿고 계획하지만 확인되지 않은 것 / 제약 = 바꿀 수 없는 것}
\begin{itemize}
\item 제약: 컷오버 창 2027-02-26\til{}28, 규제 기한 2개, 집행 한도, FP 12,400점 계약
\item 가정: P2 인터페이스 준비 / 대성로지스 준실시간 API 합의 / 가맹 부속서 2026-08-14 / 박세연의 13본 분석 시간
\end{itemize}
\keymsg{가정은 리스크 등록부의 첫 번째 원천 — 가정이 틀리면 그것이 리스크}
\note{§3.2 "항목별 의미" 뒷부분. 흔한 오류: 168억 하나만 적어 "우리 예산이 얼마인가"에 답이 세 개라는 사실을 숨기는 것. 가정 네 개는 모두 P1 PM의 통제 밖에 있다 — P2가 통합시험 전 인터페이스를 준비한다, 대성로지스가 준실시간 API에 합의한다, 가맹 부속서 개정이 2026-08-14 완료된다, 박세연이 인터페이스 13본 분석에 시간을 낸다. 5교시 §3.5에서 착수 리스크의 첫 원천이 바로 이 가정 목록임을 다시 본다. 예상 소요 4분.}
\end{frame}

\begin{frame}{헌장 항목 ④ 거버넌스와 PM 권한 — 헌장의 심장}
\begin{itemize}
\item 스폰서(정미란), 운영위 구성·주기(월 1회), 에스컬레이션 경로
\item 변경 통제 경로: CCB 월 셋째 목요일 / 변경협의체 격주 수요일
\item \textbf{PM이 스스로 결정할 수 있는 범위}: 원가 +6\%(10.08억), 일정 +15영업일, 범위 ±35건, 컨틴전시 10.2억
\item 흔한 오류: "PM은 프로젝트를 총괄한다" 한 줄
\end{itemize}
\keymsg{리스크(5\til{}10개, 프리모템 결과) · 이해관계자 요약 · 승인 — 서명 없는 헌장은 헌장이 아니다}
\note{§3.2 "항목별 의미" — 거버넌스와 PM 권한. 이 항목이 헌장의 심장인 이유는 다른 모든 항목이 "무엇"을 말할 때 이 항목만 "누가"를 말하기 때문이다. 나머지 항목은 한 줄씩: 리스크(5\til{}10개, 프리모템 결과), 이해관계자 요약(§3.3), 승인(서명 없는 헌장은 헌장이 아니다). 4교시에 이 칸을 쓰면서 막히는 지점이 5교시 "권한 없는 PM"의 출발점이다. 예상 소요 4분.}
\end{frame}

\begin{frame}{한국 기업에서 헌장이 형식화되는 네 가지 이유 \ra{} 실질화하는 다섯 가지 방법}
\begin{columns}[T]
\begin{column}{0.44\textwidth}
\subhead{왜 형식이 되는가}
\begin{enumerate}
\item 결정이 헌장 \textbf{밖에서} 이미 내려졌다 — 부의안·품의서의 사본
\item PM 권한 항목을 채울 권한이 PM에게 없다 — 전결 5억과 안 맞는다
\item 스폰서가 헌장을 읽지 않는다 — 대리 서명, 전자결재 자동 승인
\item 헌장이 계약서와 충돌한다 — SI의 진짜 범위 문서는 계약서
\end{enumerate}
\end{column}
\begin{column}{0.54\textwidth}
\subhead{어떻게 살리는가 — 스폰서에게 결정을 요구}
\begin{enumerate}
\item 헌장을 \textbf{결정 요청서}로 — 우선순위·허용범위·편익 오너·미확정 가정을 빈칸으로
\item \textbf{가정·제약 목록을 가장 긴 항목으로}
\item PM 권한을 \textbf{위임전결규정과 대조}
\item \textbf{과거 실패에 대한 답}을 넣는다 — 승인 조건 5항
\item \textbf{게이트마다 다시 연다} — G1·G2·G3
\end{enumerate}
\end{column}
\end{columns}
\keymsg{여러분의 마지막 헌장에는 여덟 가지 오류 중 몇 개가 해당하는가}
\note{§3.2 "한국 기업에서 헌장이 형식화되는 이유"와 "실질화하는 방법". 방법 1의 문형: "이 빈칸을 채우면 헌장이 완성됩니다" — 대표(일정)·CFO(비용)·수행사(범위)의 세 시계를 그대로 적어 스폰서가 하나를 고르게 하고, 고르지 않으면 그 사실을 기록한다. 방법 2는 "내가 통제하지 못하는 것이 내 일정을 결정한다"를 문서로 남기는 것. 방법 3의 예: "컨틴전시 10.2억은 전결 5억 초과 \ra{} 5억 초과 건은 스폰서 사전 결재" — 어느 쪽이든 문서에. 방법 4: 승인 조건 5항(2020·2023 감사) \ra{} 편익 프로필 8종 오너 실명, 중단 시 손실 확정 절차, 위탁 산출물 형상관리·소스 인수. 방법 5: 게이트에서 목적·성공 기준·가정 재확인. 마지막에 "실무에서 어떻게 틀리는가"의 여덟 오류(킥오프 자료로 대체 / 목적에 솔루션 / 제외 범위 비움 / 예산 한 숫자 / 권한 "총괄" / 편익을 PM 기준에 / 가정 없음 / 읽지 않은 헌장에 서명)를 읽고 "여러분 마지막 헌장에 몇 개가 해당하는가"를 묻는다. 예상 소요 8분.}
\end{frame}

\begin{frame}{이해관계자 — Mitchell-Agle-Wood 현저성 모델 (1997, AMR): 3속성 7유형}
\fig[0.7]{fig_3_3_salience}
\figcap{그림 3-4. 권력 · 정당성 · 긴급성의 벤 다이어그램 7유형 + 온담 인물 (G0 예시 판단)}
\keymsg{Power/Interest 격자의 세 한계 — 긴급성 축 없음 / 권력자와 약자를 같은 칸에 / 정적}
\note{§3.3 "개념", "Power/Interest Grid", "Mitchell-Agle-Wood". 이해관계자 = 프로젝트에 영향을 주거나 받는 개인·집단·조직. 정의는 넓어야 한다 — 좁게 잡은 프로젝트는 "그런 사람이 있는 줄 몰랐다"를 반복한다. 세 단계: 식별 \ra{} 분석 \ra{} 참여 계획. 표준: PMBOK 8판 Initiating의 둘째 프로세스, Stakeholders 도메인이 Communications 흡수 / PRINCE2 organizing practice + 3역할 / ISO 21502 §7.12. Power/Interest 2×2(밀착 관리 / 만족 유지 / 정보 제공 / 모니터링)는 10분이면 그리지만 그림 오른쪽 상자의 세 한계가 있다. 그래서 세 속성 — 권력(자기 의지를 관철할 수단), 정당성(요구가 조직·사회적으로 정당한가), 긴급성(즉각 대응을 요구하는가) — 으로 7유형: Dormant(P), Discretionary(L), Demanding(U), Dominant(P+L), Dangerous(P+U), Dependent(L+U), Definitive(P+L+U). 워크북 §3은 2×2 대신 이 모델을 쓴다. 7유형 이름을 한국어·영어로 함께 익히게 한다 — 등록부 열 이름이다. 예상 소요 8분.}
\end{frame}

\begin{frame}{온담의 이해관계자를 이 모델로 읽기 — 가장 위험한 사람은 권력자가 아니다}
\begin{columns}[T]
\begin{column}{0.5\textwidth}
\subhead{권력이 있는 쪽}
\begin{itemize}
\item \textbf{정미란 CFO} — P+L+U, Definitive
\item \textbf{김선호 대표} — Definitive, 일상은 Dormant처럼 \ra{} 게이트에서 깨어난다
\item \textbf{오정근} Dominant \ra{} 가맹 반발 시 Definitive
\item \textbf{박세연} Dominant — 권력의 원천은 직위가 아니라 \textbf{지식}(유일한 해독자)
\end{itemize}
\end{column}
\begin{column}{0.48\textwidth}
\subhead{권력이 없는 쪽 — 연합할 동기}
\begin{itemize}
\item \textbf{서정필 협의회장} — L+U, Dependent. 388점 서면 동의 거부라는 집합적 권력 \ra{} 오정근과 연합 시 Definitive
\item \textbf{김미르 노조 지부장} — Dependent \ra{} \textbf{Dangerous} 가능
\item \textbf{최영주 CPO} — Dependent, 박세연이 움직여야
\item \textbf{강윤지 지역매니저} — Discretionary \ra{} 슈퍼유저 124명 대표가 되면 활용률 좌우
\end{itemize}
\end{column}
\end{columns}
\keymsg{속성은 고정되어 있지 않다 — 2019년의 실패는 정당성·긴급성만 있는 이들이 연합했을 때 일어났다}
\note{§3.3 "온담의 이해관계자를 이 모델로 읽기". 실무 가치 ①: 속성은 고정되어 있지 않다 — Dormant가 긴급성을 얻으면 Dangerous, Dependent가 Dominant와 연합하면 Definitive. 실무 가치 ②: 유형마다 대응이 다르다 — Dangerous에게는 정당성을 부여(테이블에 앉힌다)하거나 권력을 차단, Dependent에게는 그를 대변할 Dominant를 찾아 준다. 한동석 물류본부장은 Dominant(P1 관심 낮음)이나 P2·WMS 연동에서 Definitive가 된다. 판단은 예시이며 워크북 §3에서 25행으로 다시 한다. 윤태호 감사팀장(Dominant), 이재현 수행사 PM(Dominant, 61명·변경대가 청구권), 박준영(Dependent \ra{} 종료 시 Definitive)은 구두로. 생각해 볼 질문 4(서정필과 김미르가 연합하면 — 그 연합을 막는 것이 PM의 일인가, 정당한 것으로 테이블에 올리는 것이 PM의 일인가). 예상 소요 7분.}
\end{frame}

\begin{frame}{한국 조직 특유의 이해관계자 넷 + 참여 계획}
\begin{columns}[T]
\begin{column}{0.5\textwidth}
\subhead{서구 문헌에 잘 안 나오지만 한국에서 결정적인}
\begin{itemize}
\item \textbf{노동조합} — 인력 효율이 편익로 명시된 프로그램에서 착수부터 Definitive에 가깝다
\item \textbf{가맹점주협의회} — 데이터와 단말은 가맹점주 소유. "우리 시스템"의 절반은 본사 것이 아니다
\item \textbf{감사실 / 준법} — "어떤 문서가 있으면 지적하지 않을 것인가"를 착수에 묻는다
\item \textbf{그룹 지주사 / 별도 법인} — 계약 주체 둘, 개인정보 처리자 둘, 보안 정책 둘
\end{itemize}
\end{column}
\begin{column}{0.48\textwidth}
\subhead{참여 계획 — 분석은 대응이 있어야 의미가 있다}
\begin{itemize}
\item 현재 참여 수준 \ra{} 목표 수준 \ra{} 간극을 메울 \textbf{행동} (unaware / resistant / neutral / supportive / leading)
\item 서정필: "발주 포털은 환영, 데이터·단말은 저항" \ra{} 한 사람을 두 행으로. 행동은 \textbf{오정근을 통해}, 2026-08-14 이전
\end{itemize}
\end{column}
\end{columns}
\keymsg{행동에 이름과 날짜가 없는 참여 계획은 계획이 아니다}
\note{§3.3 "한국 조직 특유의 이해관계자", "참여 계획", "실무에서 어떻게 틀리는가". 노조: 시스템 도입은 노동 조건 변경이며 단체협약상 사전 협의 대상일 수 있다. 매장 인건비율 21.4\ra{}20.1\%, 물류 처리량이 편익로 명시되어 있으므로 "P1은 인력 프로젝트가 아니다"는 맞지만 노조가 그렇게 보지 않는다는 사실이 리스크다. 가맹점주협의회: 가맹사업법상 협의 의무, 부속서 개정은 개별 서면 동의. 감사실: 과거 지적사항의 작성자이자 종료 후 감사의 수행자 — 착수 시점에 물어 두면 종료 비용이 준다. ㈜온담푸드시스템은 100\% 종속이지만 별도 법인. 서정필의 목표는 "데이터 연동에 neutral 이상", 행동은 배포 범위·연동 수준을 동결 전에 협의회에 제시. 오류 넷: 착수 때 한 번 만들고 갱신 안 함(현저성은 게이트마다 재배치) / 직위만 적고 사람을 안 적음("오정근, 2005년 입사, 가맹사업팀장 출신, 재계약률 91\% KPI, 2019 스마트오더 중단을 겪음") / 부정적 이해관계자를 등록부에서 뺌 / PM 혼자 관리 — 서정필은 오정근이, 김미르는 한동석이, 윤태호는 정미란이. 예상 소요 7분.}
\end{frame}

\begin{frame}{3교시 핵심 정리}
\begin{enumerate}
\item 착수의 첫 번째 일은 \textbf{이미 만들어진 BC의 숫자와 가정을 검증}하고 결함을 목록으로 만들어 스폰서에게 결정을 요구하는 것
\item \textbf{산출물은 PM의 책임, 활용·성과는 사업 오너의 책임}(ITO). 편익 오너는 \textbf{상설 조직의 실명}
\item 편익은 기술이 아니라 조직 변화에서 나온다(BDN) — P1이 성공해도 편익이 안 나올 경로가 착수 시점에 보여야 한다
\item 헌장은 계획서가 아니라 \textbf{권한 문서} — 제외 범위·예산 층·가정·PM 권한을 채우고, 스폰서가 결정할 빈칸을 남긴다
\item 3속성·7유형으로 보면 P1의 가장 위험한 이해관계자는 권력자가 아니라 \textbf{연합할 동기가 있는 Dependent들}
\end{enumerate}
\keymsg{다음 교시는 바로 실습 ① — 워크북 §1·§2와 온담 정본을 준비한다}
\note{제3장 핵심 정리 1\til{}5. 편익 오너를 PM으로 지정하면 범위 방어와 책임 회피가 합리적 행동이 된다(2). 목적에 솔루션을 적지 말고 우선순위 빈칸을 남겨 결정을 요구한다(4). Power/Interest는 긴급성과 정당성을 못 잡는다; 한국 특유의 넷 — 노조·가맹점주협의회·감사실·별도 법인(5). 다음 교시가 바로 실습 ①(워크북 §1 비즈니스 케이스 요약, §2 헌장)이므로 워크북과 온담 정본을 준비시킨다. 예상 소요 3분.}
\end{frame}

% =====================================================================
\section{4교시 (90분) — 실습 ①: 온담 P1 비즈니스 케이스 요약 + 프로젝트 헌장}
% =====================================================================

\begin{frame}{이 교시에서 하는 것}
\begin{itemize}
\item \textbf{산출물 ① 비즈니스 케이스 요약} — 워크북 §1, 약 35분. P1 몫만 발라내고 각 숫자의 \textbf{분모와 책임자}를 확인한 두 쪽
\item \textbf{산출물 ② 프로젝트 헌장} — 워크북 §2, 약 50분. 15개 항목 중 \textbf{목적·성공 기준·범위/제외·예산 층·가정/제약·PM 권한·우선순위·완료 기준}
\item 진행: 과제 지시 \ra{} 작성 요령 \ra{} 흔한 오류 점검 \ra{} 완성 예시본과 대조
\item 준비물: 워크북 §0.3 온담 한 페이지 요약, §1.4·§2.4 빈 템플릿, §1.5·§2.5 완성 예시본
\item 3\til{}7년차: 각 절의 \textbf{심화 과제}(자기 회사 프로젝트로 같은 문서)를 병행
\end{itemize}
\keymsg{예시본은 "정답"이 아니라 "판단의 흔적이 보이는 문서"다}
\note{워크북 §0.1 사용법 — X.3 항목별 가이드가 가장 중요하고 X.5 완성 예시본이 그 다음. 숫자에는 값·단위·출처(C-01, O-08 같은 정본 ID)를 붙이는 습관을 강조한다. 빈 템플릿 인쇄본: 템플릿/PM\_Day1/01\_비즈니스케이스요약\_빈템플릿.pdf, 02\_프로젝트헌장\_빈템플릿.pdf(통합본 00\_Day1\_템플릿팩\_빈양식.pdf). 예상 소요 3분.}
\end{frame}

\begin{frame}{산출물 ① 비즈니스 케이스 요약 — 과제 지시}
\begin{itemize}
\item \textbf{과제}: 워크북 §1.4 빈 템플릿 9개 항목 중 \textbf{항목 2(전략적 배경) · 4(편익) · 6(재무 평가) · 9(편익 책임자)}를 온담 P1로 채우라
\item 나머지 항목은 §1.5 예시본을 읽고 확인
\item 왜 PM이 다시 쓰는가: 승인 문서의 숫자를 PM이 검산하지 않으면 그 숫자는 나중에 \textbf{PM의 약속}이 된다 — 214억 안의 "물류 96억"은 P2 설비의 숫자
\item BC는 게이트마다 갱신되는 살아 있는 문서 — 갱신 기준선이 없으면 중단 결정을 못 내리는 좀비
\end{itemize}
\keymsg{시간이 부족하면 항목 4(편익)와 9(책임자)만 — 이 둘이 헌장 성공 기준의 입력이다}
\note{워크북 §1.1\til{}1.2. 분야별 무게(워크북 §1.1): SI는 수행사 PM이 발주사 BC를 본 적조차 없는 경우가 많다 / EPC는 FEL \ra{} FID에서 AACE Class 3 이상과 함께 고정 / NPD는 게이트 문서 자체가 BC이며 매 게이트 갱신 — 한 번 쓰고 끝내는 분야는 IT뿐. 표준 위치: PMBOK 8판 Governance·Initiating의 입력 문서 / 6판 1.2.6·4.1 입력 / PRINCE2 Outline BC \ra{} BC \ra{} 단계마다 갱신 / ISO 21502 §6.2·§7.3 / (참고) Green Book Five Case Model. 예상 소요 5분.}
\end{frame}

\begin{frame}{작성 요령 ① — 전략적 배경 · 대안 (워크북 §1.3 항목 2\til{}3)}
\subhead{항목 2 전략적 배경 — "왜 지금인가"}
\begin{itemize}
\item 3\til{}5문단, 문단마다 숫자 하나 이상. 현재 상태를 측정값으로(폐기율 3.9\%), \textbf{분모}를 밝히고(취급원가 2,530억)
\item 온담의 세 압력 — 기술 부채 한계 / IPO 시계 / 규제 시행일 + 두 번의 실패 이력. \textbf{do-nothing을 한 문단 이상}
\end{itemize}
\subhead{항목 3 대안 검토 — 최소 셋 + 기각된 대안 하나}
\begin{itemize}
\item ① 현행 유지 ② do-minimum(OS 마이그레이션 + 인터페이스 문서화, 약 18억) ③ 상용 리테일 패키지 — 기각 ④ 제안안 168억
\item 기각 사유 한 줄이라도 반드시 — "가맹·B2B·직영 세 모델을 한 패키지가 못 다루고 SAP 병행 요건 미충족"
\end{itemize}
\keymsg{do-nothing = "OS 지원 종료 상태로 615개 매장 결제를 계속 처리한다"}
\note{워크북 §1.3 항목 2\til{}3. 세 압력의 숫자: OD-POS OS 종료 2024-07, SAP ECC 2027-12 / 2027 3분기 예비심사 / 2026-09-11, 2027-03-31. do-minimum 18억은 [추가 설정]이다. 대안은 do-nothing / do-minimum / 제안안이 최소이고 기각된 대안이 하나 이상 있어야 검토이지 사후 정당화가 아니다. 예상 소요 4분.}
\end{frame}

\begin{frame}{작성 요령 ② — 편익 · 비용 · 재무 평가 · 편익 책임자 (항목 4\til{}6, 9)}
\begin{itemize}
\item \textbf{항목 4 편익} — 기준값·목표값·시점·금액·\textbf{산출 분모}·\textbf{측정 책임자}. P1 직접(43+37+38 = \textbf{118억})과 P2 의존(물류 96억)을 \textbf{분리}
\item \textbf{항목 5 비용} — 구축 8항목과 운영(연 27억 × 3년 = 81억) 구분, 3층(168.0 / 156.0 / 145.8) 정의, 5년 TCO 한 줄
\item \textbf{항목 6 재무 평가} — 회수기간 약 3.4년 / NPV(8\%, 5년) 약 +134억 / 민감도 70\%: +16, 50\%: $-$63 / \textbf{손익분기 편익 실현률 약 66\%}
\item \textbf{항목 9 편익 책임자} — 직책과 실명: 폐기율·결품률·인건비율 오정근 / 재고 정확도 한동석 / 검토 시점은 종료 후(M+12, G5)
\end{itemize}
\keymsg{PM은 한 명도 없다 — G5에 PM은 그 자리에 없다. 이 표의 서명이 헌장 서명보다 어렵고, 그래서 더 중요하다}
\note{워크북 §1.3 항목 4\til{}6, 9와 §1.6 해설. 항목 4에서 "P1 편익을 118억으로 한정"하는 결정이 이 문서의 핵심 판단 — "편익을 정직하게 나누는 것은 겸손이 아니라 자기 방어이며 프로그램 전체의 편익 관리를 가능하게 하는 유일한 방법". 분모 열과 책임자 열이 없는 편익 표는 편익 표가 아니다. 운영비 27억 미차감 명시, 재고 정확도 91.3\%의 분모(순환실사 SKU 69\%)를 주석으로. 항목 5의 흔한 오류: "SI 개발비 68.2억"만 적어 예산이 실제의 40\%로 보이는 것 / 운영비 누락 / 예비비를 빼고 "예산 안에 맞췄다". 항목 6: 편익과 비용을 같은 시간축에, 할인율과 실현 기간 명시. "편익이 몇 \% 실현되어야 손익분기인가"가 NPV보다 의사결정자에게 유용하다(118억의 66\% ≈ 78억이 M+24부터 매년). 이사회의 관심은 IRR이 아니라 "부채비율 187\%인 회사가 1차연도에 얼마를 쓰는가" \ra{} 연차별 현금 유출을 함께. 66\%가 헌장 항목 14 조기 종료 조건으로 흘러간다는 점을 미리 말한다. 운영비 27억 전액을 P1에 귀속한 것은 보수적 검산. 항목 9의 온라인·배달수수료 책임자 오정근은 [추가 설정], 재고 정확도 분모 확장은 나영선과 협의. 예상 소요 8분.}
\end{frame}

\begin{frame}{산출물 ① 흔한 오류 점검}
\begin{itemize}
\item 원본 승인 문서(2025-12 이사회 부의안)를 적지 않아 요약본이 독립 문서처럼 돌아다닌다
\item 솔루션을 배경으로("클라우드 POS가 필요하다") / 형용사만 있고 숫자가 없다("노후화가 심각하다")
\item 제안안 하나만 적고 "대안 없음" / 214억을 그대로 옮기고 분모를 확인하지 않았다 / 측정 책임자 "추후 지정"
\item 편익 100\% 실현 전제로 IRR 40\% / 할인율이 없다 / 리스크 칸에 "일정 지연" 같은 일반론
\item 이사회 승인 조건 5개를 빠뜨렸다 / 다음 갱신 시점(G1 2026-04-30)이 없다
\end{itemize}
\keymsg{점검 질문 — 편익 표 모든 행에 분모와 실명이 있는가 / 손익분기 실현률을 한 문장으로 말할 수 있는가}
\note{워크북 §1.3 각 항목의 "흔한 오류" 칸. 15분 작성 후 이 목록으로 자기 초안을 점검하게 한다. 추가 오류: 대안 간 비교 기준이 다르다 / 다른 프로젝트 편익을 자기 것으로 / 편익 실현 시점을 종료 직후로. 점검 질문(§1.7) 셋째: do-nothing 비용이 0으로 적혀 있지 않은가. 예상 소요 3분 (설명) + 작성 시간.}
\end{frame}

\begin{frame}{산출물 ① 예시본 참조 — 워크북 §1.5\til{}1.6}
\begin{itemize}
\item 왜 P1 편익을 \textbf{118억}으로 한정했는가
\item 왜 do-minimum을 G1까지 \textbf{폴백}으로 남겼는가 — 예산 삭감이 곧 중단이 된 2019년의 결말
\item 왜 손익분기 \textbf{66\%}를 굳이 계산했는가
\item 왜 재고 정확도 분모 문제를 별도 항목으로 뺐는가
\end{itemize}
\keymsg{완성 예시 PDF: 템플릿/PM\_Day1/01\_비즈니스케이스요약\_완성예시.pdf}
\note{5분 예시본 대조. 예시본은 정답이 아니라 판단의 흔적이 보이는 문서이며, 조의 판단이 예시본과 다르다면 그 이유를 말하게 한다. §1.7 기본 과제(CFO 148억 확정, P1 2026년 집행 102\ra{}86억 시 항목 5·6 재작성 — 이월 vs 총예산 152억 축소, 후자면 do-minimum이 되살아나는가)는 시간이 남는 조에만. 예상 소요 5분.}
\end{frame}

\begin{frame}{산출물 ② 프로젝트 헌장 — 과제 지시}
\begin{itemize}
\item \textbf{과제}: 워크북 §2.4 빈 템플릿 15개 항목 중 \textbf{2(목적) · 3(성공 기준) · 5(범위/제외) · 8(예산) · 9(가정/제약) · 12(거버넌스·PM 권한) · 13(우선순위) · 14(완료 기준)}. 나머지 7개는 §2.5 예시본으로 확인
\item 헌장이 답하는 질문은 여섯 개뿐: \textbf{왜 / 무엇이 성공 / 어디까지 / 누가 결정 / PM은 무엇을 혼자 결정 / 언제 끝}
\item 없을 때 터지는 것: \textbf{범위 경계 분쟁}(가맹 단말 교체 비용) / \textbf{PM 권한의 공백}(전결 5억 vs 10.08억) / \textbf{성공 기준 부재}(2019 "중단"이 실패인지 손절인지 판정 불가)
\end{itemize}
\keymsg{가장 중요한 칸은 "무엇을 만드는가"가 아니라 "무엇을 만들지 않는가"와 "누가 무엇을 결정하는가"}
\note{워크북 §2.1\til{}2.2. 범위 경계 분쟁은 동결 직전 주에 388점 반발과 함께 돌아오고, 2019 스마트오더는 감사 지적 "손실 확정 절차 부재"로 남았다. 분야별 실체: SI는 발주사 헌장 vs 수행사 착수계·사업수행계획서의 간극이 변경대가 분쟁의 씨앗 / EPC는 FID 문서 + PEP, "제약" 칸이 LD·성능 보증·준공 기한으로 훨씬 무겁다 / NPD는 게이트 통과 문서, 시장 진입 시점과 target cost가 성공 기준. PM이 초안을 쓰고 스폰서가 서명하는 순서가 헌장을 왜곡한다(PM은 자기 권한을 넓게 쓰기 어렵고 스폰서는 자기 책임을 좁게 읽는다) \ra{} 가장 협상이 필요한 칸이 항목 12. 예상 소요 5분.}
\end{frame}

\begin{frame}{작성 요령 ③ — 목적 · 성공 기준 · 범위/제외 (워크북 §2.3 항목 2·3·5)}
\begin{itemize}
\item \textbf{항목 2 목적} — 3\til{}5문장, 숫자 필수, 편익 문서 ID 인용. 오류: "플랫폼을 구축한다"(수단) / 프로그램 전체의 목적을 그대로(P1 ≠ ONE ONDAM)
\item \textbf{항목 3 성공 기준} — 표: 지표 / 목표값 / \textbf{판정 시점} / \textbf{판정자}. 관리 성공(G4): +15영업일·+6\% 이내, RTM 100\%, 결함밀도 0.4건/FP, SAC 0건 / 프로젝트 성공(G5): §1 편익 표 인용
\item \textbf{항목 5 범위} — 제외를 포함과 \textbf{같은 무게}로, 각 제외에 "그러면 누가 하는가". 제외에 \textbf{가맹점 단말 하드웨어 교체} 추가
\item 경계 조건: 가맹 단말은 기존 HW에 SW 배포까지 / 3PL WMS 연동 명세는 작성하되 계약 변경 협상은 물류본부
\end{itemize}
\keymsg{오류 — "기타 부대 업무 일체": 수행사는 무한 범위로, 발주사는 무상 서비스로 읽는다}
\note{워크북 §2.3 항목 2·3·5. 목적 온담 문안: "615개 매장·물류센터·공장의 주문과 재고를 하나의 원장으로 실시간 관리할 수 있게 하여 폐기율 3.9\%\ra{}2.2\%, 결품률 4.7\%\ra{}1.5\%, 가맹 발주 리드타임 38h\ra{}12h를 달성하고(연 직접 편익 118억), 지원 종료된 OD-POS를 규제 기한(2026-09-11, 2027-03-31) 안에 대체한다". 성공 기준 오류: "안정적으로 가동한다" / 판정자·시점 없음 / 편익 목표를 G4에 판정. 항목 4(상위 요구사항 5\til{}10줄, 요구 제공자 괄호), 6(산출물 — 소스·형상 인수, 운영 문서 포함; 2023 감사 지적), 7(마일스톤 — 외부 고정 날짜 굵게)은 구두로 한 줄씩. §2.6 해설 "왜 제외 범위에 가맹점 단말 HW를 굵게 넣었는가" — 헌장에서 가장 논쟁적인 줄이고 그래서 가장 필요한 줄. 예상 소요 8분.}
\end{frame}

\begin{frame}{작성 요령 ④ — 예산 · 가정/제약 · PM 권한 · 우선순위 · 완료 기준 (항목 8·9·12·13·14)}
\begin{itemize}
\item \textbf{8 예산} — 3층 표 + 컨틴전시(PM 10.2억, 단 5,000만 원 이상 단일 건은 CFO 결재)와 관리예비비(스폰서 12.0억) 구분
\item \textbf{9 가정/제약} — 가정에 \textbf{확인 시점·책임자}(3PL 2026-06 한동석 / 부속서 2026-08-14 오정근 / 겸임 투입률 50\% 박세연), 제약에 \textbf{출처}
\item \textbf{12 PM 권한} — 허용범위 표 + 에스컬레이션 3단계 + \textbf{"PM은 이 범위 안에서 스폰서 보고 없이 결정한다"} + 전결 충돌 처리
\item \textbf{13 우선순위} — 고정 / 조정 가능 / 수용. \textbf{스폰서가 결정}. 예시본: 컷오버 창 고정, 원가 한도 내 고정, \textbf{범위(옴니채널 고도화)를 조정 변수로}
\item \textbf{14 완료 기준} — 정상 종료(G4 SAC 0건, 판정 박준영·승인 정미란) + \textbf{조기 종료 조건}(손익분기 66\% 미달 전망 / 컷오버 창 2회 이탈)
\end{itemize}
\keymsg{세 개를 다 고정하면 조정 변수는 품질뿐 = 결함을 운영으로 넘긴다는 뜻}
\note{워크북 §2.3 항목 8·9·12·13·14. 항목 9의 제약: 성수기 변경 금지 창 / 규제 기한 2개 / 컷오버 창 / 본부장 전결 5억 / OD-POS·SAP 병행 / 자사앱 소스 벤더 보유. 오류: 가정이 희망사항("충분한 인력이 확보된다") / 제약에 자기 팀이 정한 것. 항목 12의 허용범위: 원가 +6\% / 일정 +15영업일 / 범위 ±3\% / 품질 0.4건/FP / 리스크 EMV 8억 / 편익 $-$5\%p, 전결 충돌은 "허용범위 내라도 5억 초과 건은 CFO 결재 병행". 항목 13: 대표는 일정, CFO는 원가, 수행사는 범위를 고정하려 한다; 범위 중 규제·가용성은 고정, 고도화는 이연 가능. 항목 14: 조기 종료 시 스폰서가 이사회에 중단 부의, 손실 확정은 재경팀장. §2.6 해설의 두 문장을 반드시 — "초과가 예상되는 시점"이라는 단어 하나가 거버넌스가 작동하는지 장식인지를 가른다 / 중단은 실패가 아니다, 손익분기에 못 미칠 것이 확실한 프로젝트를 계속하는 것이 실패다(이사회 조건 5항에 대한 직접 답). 예상 소요 10분.}
\end{frame}

\begin{frame}{산출물 ② 흔한 오류 점검 — 헌장의 여덟 가지 오류 + 넷}
\begin{itemize}
\item 킥오프 발표 자료로 대체 / 목적 칸에 솔루션 / 제외 범위 비움 / 예산을 한 숫자로
\item PM 권한을 "총괄"로 / 편익 목표를 PM 성공 기준에 / 가정을 안 적음 / 스폰서가 읽지 않은 헌장에 서명
\item 승인자 칸에 "경영진"·"운영위원회" — 위원회는 서명하지 않는다
\item 우선순위를 전부 "고정"으로 — 희망이지 우선순위가 아님 / 종료 조건이 없어 무한 하자보수
\item Senior User 서명 없이 나중에 "우리는 동의한 적 없다"
\end{itemize}
\keymsg{점검 — PM이 스폰서에게 묻지 않고 결정할 수 있는 금액과 기간을 한 문장으로 말할 수 있는가}
\note{교재 §3.2의 여덟 오류와 워크북 §2.3의 추가 넷. 25분 작성 후 이 목록으로 점검한다. 점검 질문(§2.7): 목적 문단에 "구축한다", "도입한다"로 끝나는 문장이 있는가 \ra{} 뒤에 "그래서 무엇이 얼마나 좋아지는가"를 붙이라 / 제외 범위 각 항목 옆에 "그러면 누가 하는가"가 있는가 / PM 권한 금액과 기간을 한 문장으로 말할 수 없다면 아직 권한 부여 문서가 아니다. 예상 소요 3분 (설명) + 작성 시간.}
\end{frame}

\begin{frame}{산출물 ② 예시본 참조 — 워크북 §2.5\til{}2.6}
\begin{itemize}
\item 왜 제외 범위에 \textbf{가맹점 단말 HW}를 굵게 / 왜 성공 기준을 \textbf{두 층}으로
\item 왜 우선순위에서 \textbf{범위}를 조정 변수로 / 왜 \textbf{"초과가 예상되는 시점"}이라고 썼는가
\item 왜 허용범위 10.08억과 전결 5억의 \textbf{충돌을 헌장에} 적었는가
\item 왜 조기 종료 조건에 \textbf{손실 확정 절차}까지 넣었는가
\end{itemize}
\keymsg{완성 예시 PDF: 템플릿/PM\_Day1/02\_프로젝트헌장\_완성예시.pdf}
\note{10분 예시본 대조. 조별로 항목 13(우선순위) 결정을 발표시키면 대표·CFO·수행사의 세 시계가 조마다 다르게 해소되는 것이 보인다. §2.7 기본 과제(대표의 "컷오버를 2026-11로" 요구 시 항목 13·9·10 재작성과 스폰서 제시 대안 세 개)는 시간이 남는 조에만. 예상 소요 5분.}
\end{frame}

\begin{frame}{4교시 정리 — 두 문서에서 반복된 규칙}
\begin{enumerate}
\item \textbf{숫자에는 값·단위·출처·분모} — 214억이 아니라 "P1 직접 118억(43+37+38), 분모 취급원가 2,530억, 운영비 27억 미차감"
\item \textbf{사람은 부서명이 아니라 직책·실명}, 편익 책임자는 PM이 아니다 — G5에 PM은 그 자리에 없다
\item 헌장은 채워 넣는 문서가 아니라 \textbf{스폰서에게 결정을 요구하는 문서} — 우선순위·허용범위·편익 오너·미확정 가정
\item \textbf{PM이 통제하지 못하는 것을 문서로 남기는 것}(가정·제약·제외 범위)이 PM의 권한을 지키는 방법
\item BC의 66\%와 헌장의 조기 종료 조건은 한 줄로 이어진다 — 오늘 쓴 것은 G1(2026-04-30)에 v1.0\ra{}v1.1로 확정되는 \textbf{초안}
\end{enumerate}
\keymsg{5교시가 설명하는 것 — "방금 쓴 PM 권한 칸이 왜 그렇게 쓰기 어려웠는가"}
\note{워크북 §6 "공통 규칙" 세 가지를 앞당겨 말한다. 5교시는 헌장 항목 12(거버넌스·PM 권한)와 10(리스크)의 이론이다 — "방금 쓴 PM 권한 칸이 왜 그렇게 쓰기 어려웠는지"를 5교시가 설명한다고 연결. 예상 소요 3분.}
\end{frame}

% =====================================================================
\section{5교시 (75분) — 거버넌스 · 권한 없는 PM · 프리모템 · 분야별 착수 노하우}
% =====================================================================

\begin{frame}{이 교시에서 배우는 것}
\begin{itemize}
\item 거버넌스 — 스폰서·운영위원회·PMO·PM의 권한 구조 [§3.4]: Manage by Exception과 허용범위 격자 / "권한 없는 PM"이 한국에서 왜 일반적이며 어떻게 일하는가
\item 초기 리스크와 프리모템 [§3.5]: 착수 시점 리스크의 세 원천 / 계획 오류와 외부 시각 / Gary Klein의 프리모템
\item 분야별 착수 노하우 — SI, 건설·EPC, 제조 NPD, 제약 [§3.6]
\item 6교시 실습 ②(이해관계자 등록부·RACI·프리모템)의 이론 부분
\end{itemize}
\keymsg{4교시에서 헌장 항목 12를 쓰며 막혔던 지점 — "PM 권한을 무엇으로 채우나" — 이 §3.4의 출발점}
\note{제3장 §3.4\til{}3.6. 시간 배분: §3.4 25분, §3.5 22분, §3.6 20분, 정리 5분. 4교시에서 수강생이 헌장 항목 12를 쓰며 막혔던 지점("PM 권한을 무엇으로 채우나")이 §3.4의 출발점이다. 예상 소요 2분.}
\end{frame}

\begin{frame}{거버넌스 — 조직도가 아니라 결정의 흐름}
\fig[0.7]{fig_3_4_governance}
\figcap{그림 3-5. 온담 P1 거버넌스 — 결정의 흐름(위임↓ / 예외↑)과 변경 통제 3층}
\keymsg{어떤 종류의 결정이 어디서 얼마 안에 내려지는지가 그려져야 한다 — 조직도 그리기가 아니다}
\note{§3.4 "개념". 거버넌스 = "누가 무엇을 결정하고 누구에게 책임지는가"의 구조. PMBOK 8판 첫 도메인 / ISO 21505:2017 별도 표준 / PRINCE2 첫 프로세스 Directing a Project — 세 표준이 모두 맨 앞에 두는 이유는 실패의 가장 흔한 원인이 기술이 아니라 결정의 부재이기 때문이다. 그림에서 위로 올라가는 화살표(예외 보고)와 아래로 내려오는 화살표(위임)를 짚고, 오른쪽의 변경 통제 3층(사내 CCB / 계약상 변경협의체 / 기준선 재확인)이 Day3의 주제임을 예고한다. 가장 근본적인 오류는 거버넌스를 "조직도 그리기"로 이해하는 것이다. 예상 소요 4분.}
\end{frame}

\begin{frame}{거버넌스 구성 요소 — 원칙과 온담 P1}
\begin{tbl}[\scriptsize]{L{0.16\textwidth}L{0.4\textwidth}L{0.36\textwidth}}
\textbf{구성 요소} & \textbf{원칙} & \textbf{온담 P1} \\ \midrule
\textbf{스폰서} (Executive) & 자원을 대고 BC에 책임. \textbf{한 명}. 위원회는 스폰서가 될 수 없다 & 정미란 CFO — 프로그램 스폰서 겸임(P1 별도 스폰서 없음) \\
\textbf{운영위원회} & 스폰서가 의장인 결정 기관. \textbf{"결정 요청 사항"이 없는 운영위는 열지 않는 편이 낫다} & 프로그램 운영위(월 1회) + 프로그램 이사회(분기 + 게이트, 의장 김선호) \\
\textbf{PMO} & 표준·도구·보고·형상관리·QA 지원. \textbf{PMO는 결정하지 않는다 — 결정을 준비하고 기록한다} & 프로그램 PMO 4 + P1 PMO 3(상주 24개월, 12.6억) + P2 PMO 2 \\
\textbf{PM} & 산출물 인도, 허용범위 안에서 결정, 벗어나면 에스컬레이션 & \\
\textbf{변경 통제 기구} & 권한 회수(CCB)와 권한 위임(Change Authority)을 구분 & \textbf{3층} — 사내 CCB(7인, 정족수 5, 반대의견 실명) / 변경협의체(3+3) / 기준선 재확인(분기 또는 FP ±3\%) \\
\end{tbl}
\keymsg{PMO가 결정하기 시작하면 스폰서가 결정하지 않게 된다}
\note{§3.4 "개념"과 "실무에서 어떻게 틀리는가". 사내 CCB는 월 셋째 목요일, 변경협의체는 격주 수요일. P1 스폰서가 프로그램 스폰서와 같은 사람이라는 점은 워크북 §4의 논점이다 — 생각해 볼 질문 2(P1 스폰서 = 프로그램 스폰서는 장점인가 단점인가). 예상 소요 4분.}
\end{frame}

\begin{frame}{Manage by Exception — 허용범위 격자 (온담 초안, 사례 §5.6)}
\begin{columns}[T]
\begin{column}{0.55\textwidth}
\fig[0.66]{fig_3_4_tolerance_grid}
\figcap{그림 3-6. 허용범위 격자 7축 × 3계층 (온담 초안 수치)}
\end{column}
\begin{column}{0.43\textwidth}
\begin{itemize}
\item P1 PM은 원가 10.08억 · 일정 15영업일 · 범위 ±35건 안에서 \textbf{스폰서에게 묻지 않고 결정}한다
\item 검증 ① \textbf{사내 규정} — 원가 10.08억은 본부장 위임전결 5억을 초과
\item 검증 ② \textbf{BC 편익 민감도} — 편익 $-$5\%p = 약 10.7억이 사라져도 프로젝트 계층에서 처리한다는 뜻
\end{itemize}
\end{column}
\end{columns}
\keymsg{벗어날 것이 예측되는 순간(실제로 벗어난 후가 아니라) Exception Report}
\note{§3.4 "Manage by Exception과 허용범위". 격자 수치: 원가 +2\% / +6\%(10.08억) / +12\% — 일정 +5 / +15 / +30영업일 — 범위 ±1\% / ±3\%(±35건) / ±6\% — 품질 결함밀도 0.13 / 0.4 / 0.8건/FP — 리스크 EMV 2.7억 / 8억 / 16억 — 편익 $-$1.7\%p / $-$5\%p / $-$10\%p — 안전·지속가능성 0 / 0 / 0(작업패키지 / 프로젝트 / 프로그램). 2교시 PRINCE2 슬라이드의 격자를 온담 숫자로 다시 본다. CFO가 "근거 없는 편익은 승인 안 함"이라 한 상황에서 편익 폭이 적절한지는 스폰서가 결정한다. 6교시 워크북 §4 항목 4(허용범위 위임 요약)에서 이 표를 3층으로 옮기고 전결 충돌 칸에 주석을 단다. Day2에서 확정. 예상 소요 6분.}
\end{frame}

\begin{frame}{"권한 없는 PM"이 한국에서 왜 일반적인가}
\begin{itemize}
\item 예산 집행권은 재경과 본부장에게, 인사권은 부서장에게, 계약 변경권은 구매·법무에게, 요구사항 확정권은 현업 임원에게 — PM에게 남는 것은 \textbf{일정표와 회의 소집권}
\item 개인의 문제가 아니라 \textbf{구조의 문제} — 이유 셋
\item ① \textbf{매트릭스 조직에서 기능 부서의 권한이 프로젝트의 권한보다 세다} — 임시조직에 권한을 주면 영구조직이 잃는다
\item ② \textbf{위임전결규정이 직위 기준이지 역할 기준이 아니다} — "본부장 5억"은 있어도 "PM 10억"은 없다
\item ③ \textbf{스폰서가 권한을 위임하지 않고 보고를 요구한다} — 허용범위가 없으니 모든 결정이 보고 대상
\end{itemize}
\keymsg{보고가 많으니 PM은 결정하는 사람이 아니라 보고하는 사람이 된다}
\note{§3.4 "권한 없는 PM이 한국에서 왜 일반적이며". 이유 ①은 1교시 임시조직론의 귀결이고, ②는 PM이 직위가 아니라 역할이므로 규정에 자리가 없다는 뜻이며, ③은 PMBOK 8판 원칙 6 "권한 위임 문화"가 한국에서 이상과 현실의 거리가 가장 먼 원칙인 이유다. 4교시에 헌장 항목 12를 쓰면서 겪은 어려움이 바로 이 세 이유라고 연결한다. 예상 소요 5분.}
\end{frame}

\begin{frame}{권한 없는 PM이 일하는 세 가지 방법}
\begin{itemize}
\item ① \textbf{허용범위를 문서로 획득한다} — "PM은 총괄한다"가 아니라 "PM은 원가 10.08억, 일정 15영업일, 범위 ±35건 안에서 결정한다"를 \textbf{스폰서 서명}으로
\item ② \textbf{결정을 결정권자에게 가져간다 — 선택지 + 권고 + 기한을 붙여서} — 스폰서는 결정하거나 결정하지 않은 책임을 진다. 어느 쪽이든 \textbf{기록된다}
\item ③ \textbf{에스컬레이션을 BATNA로 쓰지 않는다} — 한 번 쓰면 관계 자본이 소진되고, 임원이 PM 편을 들지 않으면 공개적으로 무너진다
\item 영향력의 실체는 \textbf{상대의 BATNA} — 박세연이 인터페이스 분석에 시간을 내지 않으면 그가 잃는 것은 "우리 손으로 운영할 수 있는 시스템"이라는 그 자신의 목표
\end{itemize}
\keymsg{"권고는 (b). 결정 기한은 2026-06-30, 그 이후는 (a)만 남는다"}
\note{§3.4 "어떻게 일하는가". ①은 PM이 권한을 얻는 유일한 표준 근거 경로다. ②의 문형 예시: "P2 리드타임이 41주로 늘면 (a) 컷오버를 2027-03-12로 미룬다(G3 합류 버퍼 15영업일 소진) (b) 물류 인터페이스를 컷오버 범위에서 빼고 2차 오픈으로 (c) P2 벤더에 단축 대가 지급. 권고는 (b). 결정 기한은 2026-06-30, 그 이후는 (a)만 남는다." 이 문형은 워크북 §2.7 기본 과제(대안 세 개)와 6교시 RACI에서 "PM은 R, A는 임원"인 행에서 실제로 쓰게 된다. ③ 매트릭스 PM의 BATNA는 대개 에스컬레이션이고 그것은 나쁜 BATNA다. 오류 목록: 운영위 명단에 3역할 미배정, 운영위를 정보 공유회로, PMO가 결정, 허용범위 없이 Exception Report 양식만, 에스컬레이션을 위협으로. 예상 소요 7분.}
\end{frame}

\begin{frame}{착수 시점의 리스크는 어디서 나오는가 — 세 원천}
\begin{tbl}[\scriptsize]{L{0.12\textwidth}L{0.3\textwidth}L{0.5\textwidth}}
\textbf{원천} & \textbf{질문} & \textbf{온담 P1} \\ \midrule
\textbf{가정에서} & 헌장의 가정 각 항목 — "이 가정이 틀리면 무엇이 일어나는가" & "대성로지스가 준실시간 API에 합의한다"가 틀리면 \ra{} 센터 재고가 야간 배치로 남고 \ra{} 재고 정확도 98.5\%가 구조적으로 불가능 \ra{} B-02·B-03·B-04가 흔들린다. \textbf{2022년에 정확히 이것이 일어났다} \\
\textbf{구조에서} & 임시조직론·Engwall의 "이미 결정된 것들" & 통제 밖 의존(P2, 3PL, 규제, 노조) / 과거 실패의 반복 패턴 / 이해관계자 연합 가능성 / \textbf{핵심 인물 의존}(박세연 — OD-POS 원 개발자 3명 중 재직 1명) \\
\textbf{base rate에서} & 1교시 §1.2 reference class & 정보화 사업 4건 중 2건 중단·축소 / IT 원가 초과의 멱함수 분포 — "이 프로젝트는 다르다"는 믿음에 대한 외부 시각 \\
\end{tbl}
\keymsg{아직 WBS도 일정도 없다 — 착수 시점 리스크는 활동 단위가 아니라 가정·구조·base rate 단위}
\note{§3.5 "착수 시점에 리스크를 보는 법". 첫 번째 원천이 4교시에 쓴 헌장 항목 9(가정)이며, 그래서 가정 목록을 길게 쓰라고 한 것이다. 구조에서 나오는 리스크의 세부: 과거 실패의 반복 패턴은 편익 근거 부재·형상관리 미비, 이해관계자 연합은 가맹·노조·매장. 예상 소요 5분.}
\end{frame}

\begin{frame}{계획 오류(planning fallacy)와 외부 시각}
\begin{itemize}
\item Kahneman \& Tversky(1979): 사람들은 과제 완료 시간을 체계적으로 과소추정한다
\item Buehler, Griffin \& Ross(1994): \textbf{자기 과거 실적 정보를 주어도 예측이 개선되지 않는다}
\item \textbf{내부 시각}: 세부 계획을 보고 "이렇게 하면 된다" / \textbf{외부 시각}: 유사 프로젝트들의 한 사례로 보고 "그 부류는 대개 이렇게 되었다"
\item Lovallo \& Kahneman(2003, HBR)이 경영 언어로, Flyvbjerg의 RCF가 절차로
\item "17개월·156억이 맞는가"보다 먼저 — \textbf{"비슷한 회사가 비슷한 재구축을 했을 때 대개 얼마나 초과했는가"}
\end{itemize}
\keymsg{그 데이터가 없다면 — 대개 없다 — "데이터가 없다"는 사실 자체가 첫 번째 리스크}
\note{§3.5 "계획 오류와 외부 시각". Newby-Clark et al.(2000): 낙관 시나리오에 집중하고 비관 시나리오를 무시한다. 자신의 과거 지연을 알면서도 이번에는 다를 것이라 믿는 것이 계획 오류의 핵심이며, 그래서 outside view가 필요하다. 1교시 RCF의 심리학적 근거. 온담에는 과거 IT 투자 편익 실현률 기록이 없다는 것 자체가 감사 지적사항이라는 워크북 §1.6 해설과 연결. 예상 소요 5분.}
\end{frame}

\begin{frame}{프리모템 — Gary Klein (2007, HBR, 2페이지)}
\fig[0.62]{fig_3_5_premortem}
\figcap{그림 3-7. 프리모템 절차와 "실패했다"의 문법 — prospective hindsight (Mitchell, Russo \& Pennington, 1989)}
\keymsg{"지금은 프로젝트가 끝난 시점이고, 완전히 실패했다. 각자 2분 동안 그 실패의 이유를 적으라"}
\note{§3.5 "프리모템". 절차: 팀이 계획을 공유한 뒤 진행자가 실패를 완료된 사실로 선언하고, 각자 2분 동안 이유를 적은 뒤 돌아가며 한 가지씩 말하고 기록한다. 설계 핵심 둘 — (a) "실패했다"를 완료된 사실로 전제: "실패할 수 있다면 왜"가 아니라 "실패했다. 왜". Mitchell, Russo \& Pennington(1989)의 prospective hindsight 실험은 미래를 이미 일어난 일로 회고하게 하면 이유를 생성하는 능력이 유의하게 높아진다는 것을 보였다 — 문법의 차이가 효과를 만든다. (b) 약점을 아는 반대자가 안전하게 말할 수 있는 장: 킥오프 직후의 팀은 낙관과 몰입으로 차 있고 "이건 안 될 것 같다"의 사회적 비용이 크다. 프리모템은 그 발언을 과제로 만들어 비용을 없앤다. 게이트마다 반복한다 — G0, G1, G3 직전. "실패했다. 왜"의 답이 달라지는 것이 프로젝트가 어디로 가는지를 보여 준다. 예상 소요 4분.}
\end{frame}

\begin{frame}{P1에서 나올 법한 실패 이유 — 여러 항목이 같은 원인을 가리킨다}
\begin{itemize}
\item 가맹 388점 중 140점이 새 POS 배포 거부, 두 시스템 병행
\item G2 정본 결정이 박세연 한 사람에게 걸려 두 달 지연
\item P2 41주로 물류 인터페이스 시험이 컷오버 후로
\item CFO 148억 축소로 전환 용역 11.4억이 2027년으로, 리허설 3회\ra{}1회 / 수행사 대기 시간 변경대가 청구
\item 오픈 첫 달 615개 매장 콜이 서비스데스크 11명에게, SAC 미충족으로 인수 거부
\end{itemize}
\keymsg{공통 원인(가맹 협상 · 핵심 인물 · 프로그램 의존 · 재무 압박)이 최우선 리스크}
\note{§3.5 프리모템 예시와 "실무에서 어떻게 틀리는가". 각 항목이 리스크 등록부 한 행(원인·사건·영향·오너)이 되고, 여러 항목이 같은 원인을 가리키는 것이 보인다. 변경협의체 6개월 미합의도 예시에 포함된다. 실무 오류 다섯: 등록부를 "일반적 리스크 목록"(일정 지연·예산 초과·인력 이탈)으로 채움 / "실패했다"를 "실패할 수 있다"로 바꿔 브레인스토밍으로 / 결과를 헌장 가정·리스크 등록부에 연결 안 함 / 반대자의 발언을 분위기 해치는 것으로 / uniqueness bias("우리는 특별하다")를 리스크 식별에서 반복. 세 번째 오류의 처방이 워크북 §5 항목 6(즉시 조치와 헌장·계획 반영). 6교시 워크북 §5에서 실패 선언문 세 개(일정·편익·중단)로 실습. 예상 소요 3분.}
\end{frame}

\begin{frame}{분야별 착수 노하우 ① — SI: 계약 이후에 착수하는 프로젝트}
\begin{itemize}
\item 착수 시점에 이미 계약이 있다 — SI PM의 일은 헌장을 쓰는 것이 아니라 \textbf{계약 문서를 읽고 헌장으로 번역하는 것}
\item ① \textbf{세 문서의 우선순위}: 계약서 > 과업지시서 > 제안서. 제안서는 \textbf{"단방향 확대 문서"} — "무상 3개월 안정화", "AI 자동 분류"가 그대로 과업이 된다
\item ② \textbf{규모 기준선}: FP 12,400점 × 55만 원 = 68.2억. FP는 측정치가 아니라 \textbf{협상 인공물} — G1 "FP 기준선 확인서 서명"
\item ③ \textbf{인터페이스 상대의 소유자}: 47본 중 문서 부재 13본, 3PL WMS, 은행·PG 7본 — \textbf{"누가 상대 쪽을 준비하는가"의 이름}
\item \textbf{공공 vs 민간}: 과업심의위원회·감리 의무·분리발주는 \textbf{공공 발주에만}. 온담 같은 민간에는 없다 \ra{} 계약서로 직접 만든다
\end{itemize}
\keymsg{SI는 컷오버 주말에 죽지 않는다 — 진짜 사인은 데이터 정제 착수 지연과 인터페이스 상대 미확정}
\note{§3.6 "SI". RFP \ra{} 제안 \ra{} 평가 \ra{} 협상 \ra{} 계약이 끝나 있고 범위·금액·기간·산출물이 계약서에 있다. 제안서는 과업을 늘릴 수는 있어도 RFP를 축소할 수는 없다 — 발주사 PM은 제안서에 쓰인 것("성능 2배" 포함)을 착수 시점에 산출물 목록으로 확정해 둔다. FP는 발주사에게는 상한, 수행사에게는 변경대가 근거로 작동한다. 응찰가 198억\ra{}168억·손익률 6.5\% — 3주차에 WBS와 제안서 산출물이 1:1이 아니라고 밝혀지면 검수 분쟁 확정. 별도 법인 공장 데이터도 인터페이스 상대. 공공 제도의 근거: 소프트웨어 진흥법 §50 과업심의위원회, 전자정부법 감리 의무 — 국가기관 등 공공 발주에만 적용되며, 민간에서는 그 기능을 사내 CCB + 변경협의체 + 기준선 재확인으로 만든다. 공공 경험자가 민간에서 "과업심의위원회에 올리자"고 말하는 순간 그 제도가 여기 없다는 사실이 리스크. 발주사 PM(수강생)과 수행사 PM(이재현)은 같은 프로젝트의 다른 이해관계를 갖는다는 1교시 §1.5의 이원 구조 재확인. 예상 소요 7분.}
\end{frame}

\begin{frame}{분야별 착수 노하우 ② — 건설·EPC / 제조 NPD}
\begin{columns}[T]
\begin{column}{0.5\textwidth}
\subhead{건설·EPC — 발주 방식과 계약 유형이 리스크를 결정}
\begin{itemize}
\item \textbf{설계 성숙도}: AACE Class 5 \til{} 1. 어느 Class에서 계약하느냐가 신뢰 구간을 정한다
\item Class 4 이하에서 경쟁입찰 \textbf{LSTK}는 리스크 이전이 아니라 프리미엄·변경 청구·클레임으로 되돌아온다 — 온담 P2가 정확히 이 검토 대상
\item \textbf{기록 규율}(SCL Protocol): 지연 청구의 승패는 \textbf{사건 당시 생성된 기록}
\end{itemize}
\end{column}
\begin{column}{0.48\textwidth}
\subhead{제조 NPD — Stage-Gate의 Gate 1\til{}2}
\begin{itemize}
\item \textbf{게이트는 회의가 아니라 투자 결정} — Go / Kill / Hold / Recycle. 자원 약정 없는 게이트 = "이빨 없는 게이트"
\item \textbf{Must-meet / Should-meet} 이원 기준 — "3.8점이니 통과"는 규제 불가 항목이 점수에 가려진 것
\item \textbf{Kill이 0건인 게이트 프로세스는 고장난 것}
\end{itemize}
\end{column}
\end{columns}
\keymsg{온담에 던질 질문 — G1에서 P1을 중단할 수 있는가, 기준은, 중단 결정자는 착수 결정자와 다른가}
\note{§3.6 "건설·EPC"와 "제조 NPD". AACE 등급의 신뢰 구간은 고정 상수가 아니라 리스크 분석으로 산정한다 — "Class 3 = $-$20\%\til{}+30\%"처럼 표를 외워 적용하면 틀린다. 계약 유형의 역설(Merrow, IPA): Class 4 이하 LSTK는 (a) 프리미엄으로 얹히거나 (b) 시공자가 스펙을 동결하고 모든 해석 차이가 변경 청구가 되어 (c) 클레임·부실·부도로 되돌아온다 — "턴키가 안전하다"는 신념이 깨지는 지점. 온담 P2: "4,200은 현장 숫자가 아니다" = Class 4 사양으로 설비 발주(선급금 29.6억)를 앞두고 있다. SCL Protocol 2판 Core Principle 1 "프로그램과 기록"은 SI에도 성립 — "발주사 결정 지연 12일"을 주장하려면 당일 서면 통지. 회의록 확인 서명을 미루기 시작하는 발주사 담당자는 내부적으로 결정 권한을 잃었다는 신호. NPD: 정보가 가장 적은 시점의 결정. Go에는 다음 스테이지 인력·예산이 붙어야. 규제·안전·IP는 must-meet. 초기 게이트에서 전체 NPV를 요구하지 않는다 — 실물옵션: 각 스테이지 투자는 "다음 단계로 갈 권리를 사는 것". 중단 결정자 분리는 마무리 §4.4 Staw. EPC PM이 SI PM에게 줄 수 있는 것 — "추정치를 숫자 하나가 아니라 숫자 + 정확도 밴드 + 근거 산출물 목록으로 유통시키는 문화"(§1.5). 예상 소요 7분.}
\end{frame}

\begin{frame}{분야별 착수 노하우 ③ — 제약·규제: 규제 게이트의 비가역성}
\begin{itemize}
\item 착수는 규제 경로 결정에서 시작: IND \ra{} 임상 \ra{} NDA / 510(k) · De Novo · PMA / SW 의료기기는 IEC 62304 등급. 경로가 정해지면 게이트는 협상 불가
\item (a) \textbf{제출 순간 문서 세트가 잠긴다} (b) \textbf{심사 시계는 제출자가 통제 못 한다} (c) \textbf{추가자료 요구는 시계를 멈추고 재개 시점을 상대가 정한다} \ra{} float가 성립하지 않는다
\item 착수 오류: \textbf{추가자료 라운드 0회로 가정}한 뒤 출시일 공표. 처방: 라운드 수를 같은 경로의 과거 제출 건(reference class)에서 분포로
\item "규제가 애자일을 막는다"는 통념은 틀렸다 — GAMP 5 2판, FDA CSA. 개발 리듬과 산출물 베이스라인을 분리
\item 온담 P4: 개인정보 보호법 개정 대응(2026-09-11), 전자금융업 등록(심사 2027-03-31), ISMS-P(2027-06)
\end{itemize}
\keymsg{최영주: "시행일은 협상 대상이 아니다" — 규제 해석을 착수 시점에 서면으로 받아 두라}
\note{§3.6 "제약". De Novo 150일은 응답 대기 제외 심사일이다. 온담은 제약회사가 아니지만 P4 규제 워크스트림이 같은 구조다 — 처리항목 142개·위탁사 19개, 전자금융업 등록 신청 2027-01-29. 최영주 CPO는 "검토 의견을 제공하지 아키텍처를 결정하지 않는다"고 말한다 — 그래서 규제 해석을 착수 시점에 서면으로 받아 두어야 한다. 예상 소요 4분.}
\end{frame}

\begin{frame}{네 분야가 착수에서 공유하는 네 질문}
\begin{enumerate}
\item \textbf{무엇이 이미 결정되어 있는가} — 계약, 설계 성숙도, 규제 경로, 게이트 기준
\item \textbf{어느 결정이 되돌릴 수 없는가} — FP 기준선, 설비 발주, 금형, 규제 제출
\item \textbf{누가 그 결정을 소유하는가}
\item \textbf{그것을 문서로 남겼는가}
\end{enumerate}
\keymsg{헌장은 이 네 질문의 답을 적는 곳이다}
\note{§3.6 "네 분야가 착수에서 공유하는 것". 1교시 Engwall(이미 결정된 것), §1.5·§2.6 되돌림 비용, 3교시 ITO(소유), 헌장(문서)이 네 질문 하나씩에 대응한다. 6교시 실습의 세 문서가 이 답을 헌장에 되먹이는 문서라는 점으로 넘어간다. 예상 소요 2분.}
\end{frame}

\begin{frame}{5교시 핵심 정리}
\begin{enumerate}
\item 거버넌스는 조직도가 아니라 \textbf{결정의 흐름}. 허용범위 격자(7축 × 3계층)가 PM이 권한을 얻는 유일한 표준 근거 경로 — 전결 규정과 편익 민감도 양쪽으로 검증
\item 권한 없는 PM은 구조의 문제. 처방 셋 — \textbf{허용범위를 서명으로 획득 / 선택지·권고·기한과 함께 결정을 가져간다 / 에스컬레이션을 BATNA로 쓰지 않는다}
\item 착수 시점 리스크는 \textbf{가정·구조·base rate}에서 나온다. 자기 과거 실적을 알아도 예측은 개선되지 않는다 \ra{} outside view
\item 프리모템(Klein, 2007)은 \textbf{"실패했다. 왜"}의 문법으로 반대자의 발언을 과제로 만든다. 게이트마다 반복하고 헌장·리스크 등록부에 연결
\item SI는 \textbf{계약 이후 착수} / EPC는 \textbf{설계 성숙도 × 계약 유형} / NPD는 \textbf{kill 0건이면 고장} / 제약은 \textbf{규제 게이트의 비가역성}
\end{enumerate}
\keymsg{6교시는 바로 실습 ② — 4교시에 쓴 헌장 초안을 꺼내 놓는다. §3\til{}5는 헌장을 고치는 문서들이다}
\note{제3장 핵심 정리 6\til{}8. 권한 없는 PM의 세 구조적 이유: 매트릭스 / 직위 기준 전결 / 위임 없는 보고 요구. SI의 세부: 문서 우선순위, FP 기준선, 인터페이스 소유자 — 공공 제도는 민간에 없다. NPD: 자원 약정이 붙은 Gate 1\til{}2. 6교시는 바로 실습 ②(워크북 §3 등록부, §4 거버넌스·RACI, §5 프리모템). 예상 소요 3분.}
\end{frame}

% =====================================================================
\section{6교시 (75분) — 실습 ②: 이해관계자 등록부 · 거버넌스·RACI · 프리모템}
% =====================================================================

\begin{frame}{이 교시에서 하는 것}
\begin{itemize}
\item \textbf{산출물 ③ 이해관계자 등록부와 현저성 분석} — 워크북 §3, 약 25분
\item \textbf{산출물 ④ 거버넌스 구조와 RACI 초안} — 워크북 §4, 약 20분
\item \textbf{산출물 ⑤ 프리모템 기록과 초기 리스크 목록} — 워크북 §5, 약 25분. 조별로 실제로 실시한다
\item 세 문서는 4교시 헌장을 \textbf{고치는} 문서 — 등록부가 항목 11을, RACI가 항목 12를, 프리모템이 항목 9·10을 되먹인다
\item 준비물: 워크북 §3.4·§4.4·§5.4 빈 템플릿, §3.5·§4.5·§5.5 완성 예시본, 4교시 헌장 초안
\end{itemize}
\keymsg{75분에 세 산출물 — 각 산출물은 핵심 항목만 채우고 나머지는 예시본 대조와 사후 과제}
\note{워크북 §6의 되먹임 구조를 먼저 말한다. 75분에 세 산출물이므로 각 산출물은 핵심 항목만 채운다 — 등록부는 인물 6\til{}8명의 3속성·유형·관계 소유자, RACI는 10행, 프리모템은 실패 선언문 1개 + 사유 수집 + 리스크 3행. 빈 템플릿 인쇄본: 템플릿/PM\_Day1/03\_이해관계자등록부\_빈템플릿.pdf, 04\_거버넌스RACI\_빈템플릿.pdf, 05\_프리모템리스크\_빈템플릿.pdf. 예상 소요 3분.}
\end{frame}

\begin{frame}{산출물 ③ 이해관계자 등록부 — 과제 지시}
\begin{itemize}
\item \textbf{과제}: 워크북 §3.4 템플릿으로 8명 등록 — \textbf{정미란·오정근·박세연·한동석·서정필·김미르·최영주·박준영}
\item 각 행에 항목 3(기대·우려·KPI), 4(P/L/U 3속성과 7유형), 5(태도 현재\ra{}목표), 6(관여 전략·\textbf{관계 소유자}). 그다음 \textbf{현저성 분석 요약표}
\item 등록부의 가치는 명단이 아니라 \textbf{분석 열}에 있다 — 어떤 힘, 쓸 정당성, 지금 움직이려 하는가. 그리고 그 셋은 \textbf{바뀐다}
\item 없을 때 터지는 것: \textbf{보이지 않던 사람이 결정적 순간에 나타난다}(김미르) / \textbf{같은 직책을 같은 방식으로 대한다}(오정근과 한동석)
\end{itemize}
\keymsg{"요구할 정당성은 없지만 멈출 힘이 있는 사람"을 찾아 적는 것이 등록부의 절반}
\note{워크북 §3.1\til{}3.2. 김미르 — P1은 인력 프로젝트가 아니지만 이천센터 협조 없이 데이터 전환도 시험도 안 된다. 오정근과 한동석은 둘 다 본부장이지만 P1에 대한 긴급성·정당성 구조가 다르다 — 같은 월간 보고서를 보내면 둘 다에게 외면당한다. 분야별 무게 중심: SI는 현업-IT 긴장과 계약 관계, 최종 사용자가 가장 자주 누락 / EPC는 인허가·주민·환경단체·하도급이 절반 / NPD는 핵심 고객·인증기관·단일소스 부품사, 양산 이관받는 생산·품질이 누락. 표준 위치: PMBOK 8판 Stakeholders 도메인(Communications 흡수) / 6판 13.1·13.2(태도 5단계 표기법의 유래) / PRINCE2 organizing + communication management approach / ISO 21502 §7.12. 예상 소요 4분.}
\end{frame}

\begin{frame}{작성 요령 ⑤ — 3속성 판정 · 태도 · 관계 소유자 (워크북 §3.3 항목 3\til{}6)}
\begin{itemize}
\item \textbf{항목 3 기대·우려} — \textbf{그 사람의 말로}. 박세연의 우려: 벤더 종속, 그리고 \textbf{자신이 만든 OD-POS와 문서 없는 13본이 "실패"로 서술되는 것} — 공식 입장에 없지만 등록부에는 있어야
\item \textbf{항목 4 3속성·유형} — \textbf{판정 시점(G0)}을 적고 바뀔 시점과 방향을 비고에. 서정필 L+U 의존 \ra{} 2026-08 협상에서 거부권이 실질 권력 \ra{} 확정 / \textbf{김미르 P+U 위험}
\item \textbf{항목 5 태도} — 무인지/저항/중립/지지/주도, "현재 \ra{} 목표". \textbf{저항은 정보다} — 오정근의 저항은 2019년의 경험
\item \textbf{항목 6 관여 전략·관계 소유자} — 외부 협상 상대의 소유자는 PM이 아니라 소관 임원: 서정필 \ra{} \textbf{오정근} / 김미르 \ra{} \textbf{한동석} / 윤태호 \ra{} P1 PM
\end{itemize}
\keymsg{오류 — 직급으로 권력 판정(나영선의 P1 권력은 CCB 표 하나) / 긴급성 무시 / 한 번 판정하고 갱신 안 함}
\note{워크북 §3.3. 항목 3: 기대와 우려를 분리하고 가능하면 KPI(오정근: 가맹 재계약률 91\%). 박세연의 기대는 IT본부가 운영 가능한 아키텍처. 항목 4의 판정 기준: P(이 프로젝트의 결정·자원·진행을 바꿀 수 있는가) / L(조직·계약·법적으로 정당한가) / U(지금 당장 대응을 요구하는가). 정미란 P+L+U 확정 / 김선호 P+L 지배 \ra{} G3에서 확정 / 윤태호 L+U 의존 \ra{} 감사 착수 시 P 획득. 항목 5(Ford, Ford \& D'Amelio, 2008): 오정근 저항(조건부)\ra{}지지 / 박세연 중립(경계)\ra{}주도 / 강윤지 저항(불신)\ra{}지지(슈퍼유저) / 이재현 중립(계약 방어)\ra{}지지. 모두를 주도로 만들 필요 없다. 항목 6 전략은 정보 제공/협의/공동 결정/협상/모니터링: 서정필 — 협상, P1은 가맹 포털 프로토타입 시연으로 지원 / 김미르 — 모니터링·정보 제공, "P1은 인력 프로젝트가 아님"을 문서로 / 윤태호 — 컨틴전시 집행 기록 월 1회 사본. 항목 1(ID S-01 형식, 내부/외부·계약/비계약 구분), 항목 2(역할은 헌장 항목 12·RACI와 일치), 항목 7(갱신 시점 — G1, 2026-08 협상 후, G2, G3)은 구두로. 예상 소요 8분.}
\end{frame}

\begin{frame}{산출물 ③ 흔한 오류 점검}
\begin{itemize}
\item 부서명만 적고 사람을 안 적었다("영업본부") — 부서는 우려를 갖지 않는다 / 임원만 적고 실무 접점자(구매팀장, 재경팀장, 지역매니저)를 뺐다
\item "영향받는 자"(매장 직원, 물류 현장)를 뺐다 — 요구하지 않지만 결과를 매일 겪는다
\item PM의 추측을 그 사람의 기대로 / 우려를 안 적음 / 모두에게 "프로젝트 성공"
\item 모두 "지지" / 목표를 모두 "주도" / 저항을 개인의 결함으로
\item 전략 칸에 "정기 보고"만 / 관계 소유자를 안 적어 PM이 노조 협상까지 / 착수 시 한 번 쓰고 끝
\end{itemize}
\keymsg{점검 — 부서명만 있는 행이 있는가 / 관계 소유자가 전부 PM인가 / "우려" 칸이 "기대" 칸과 다른가}
\note{워크북 §3.3 각 항목의 흔한 오류. 12분 작성 후 이 목록으로 점검한다. 추가: 채널을 전부 회의로 / 판정 시점이 없어 어느 시점의 판단인지 모름. 점검 질문은 §3.7. 예상 소요 3분 (설명) + 작성 시간.}
\end{frame}

\begin{frame}{산출물 ③ 예시본 참조 — 워크북 §3.5\til{}3.6 (25행 + 요약표 + 속성 변화 예측)}
\begin{itemize}
\item 왜 김미르를 \textbf{"위험"}으로 판정했는가 — 정당한 경로가 없는 권력은 비공식 경로로 행사된다. 대응은 설득이 아니라 \textbf{정당성 경로를 만들어 주는 것}
\item 왜 박세연의 "우려" 칸에 공식 입장에 없는 것을 적었는가
\item 왜 관계 소유자가 PM이 아닌 행이 열 개 가까이 되는가
\item 왜 사모펀드와 규제기관처럼 접촉하지 않는 이해관계자를 넣었는가 — 등록부는 접촉 명단이 아니라 \textbf{힘의 지도}
\item 왜 정하늘 감리원의 목표 태도가 중립\ra{}중립인가 — 감리는 독립성이 가치
\end{itemize}
\keymsg{완성 예시 PDF: 템플릿/PM\_Day1/03\_이해관계자등록부\_완성예시.pdf}
\note{5분 예시본 대조. 김미르 대응의 구체: 인사본부 협의 채널, P1 편익에 인건비 없음을 문서로 — 2×2로는 이 판단을 할 수 없다. §3.7 기본 과제(2026-06 3PL 협상 결렬 + P2 41주 확정 시 한동석·김미르·대성로지스·박준영 네 행 재판정, 헌장 항목 9의 어느 가정이 깨지는가)는 사후 과제. 예상 소요 4분.}
\end{frame}

\begin{frame}{산출물 ④ 거버넌스 구조와 RACI — 과제 지시 + 작성 요령}
\begin{itemize}
\item \textbf{과제}: 워크북 §4.4 템플릿의 \textbf{항목 2(3역할과 핵심 직위)}와 \textbf{항목 5(RACI)} 10행. 항목 1·3·4는 §4.5 예시본으로 확인
\item \textbf{RACI 규칙 셋}: 행마다 A는 정확히 한 명 / R은 여럿 가능하나 조정자 필요 / C가 많으면 느려지고 I가 많으면 아무도 안 읽는다. \textbf{행은 활동이 아니라 결정과 산출물}
\item \textbf{항목 2 3역할} — Executive 정미란 / Senior User 오정근 / Senior Supplier \textbf{박세연}(수행사 이재현은 그 아래의 공급자)
\item \textbf{항목 5 RACI} — 결정 행과 산출물 행 구분. \textbf{통제 밖 결정(3PL 협상, 가맹 배포 범위, P2 규격)의 A가 P1 밖의 임원임을 명시}
\item "컨틴전시 집행" 행: \textbf{A = P1 PM, C = 재경팀장, I = 정미란·윤태호} — 2020 감사 지적에 대한 답은 사전 승인이 아니라 \textbf{사후 투명성}
\end{itemize}
\keymsg{한국 SI에서 가장 흔한 혼동 — 수행사 PM을 Senior Supplier로 적어 발주사 IT의 기술 승인 권한이 사라지는 것}
\note{워크북 §4.1\til{}4.3. 거버넌스 구조 = "누가 무엇을 결정하는가", RACI = 각 결정·산출물에 대해 R(실행) / A(최종 책임) / C(협의) / I(통보). 헌장 항목 12의 요약을 결정 단위로 펼친 것. 한 사람이 두 역할이면 명시하고 이해 충돌 여부를 적는다(오정근은 P3 겸임, 이해 충돌 없음). 3역할 오류: 스폰서가 둘이거나 없다 / Senior User를 IT 부서장이 겸함. RACI는 행 15\til{}25, 열 8\til{}12; 표 아래에 "A 없는 행 / A 둘인 행 = 0" 검산 문장. 이 문서가 없을 때 온담에서 터지는 것 셋 — 컨틴전시 승인 근거 부재(2020 감사) / 결정의 이중화(CCB vs 변경협의체가 같은 변경을 다른 각도에서) / 경계 결정의 공백(3PL·가맹·P2 규격은 PM이 R이지 A가 아닌데 A가 적혀 있지 않으면 아무도 안 한다). 항목 1의 관계 문장 "변경협의체 합의 \ra{} CCB 승인 \ra{} 계약 변경"과 항목 3의 시한(수행사 PM \ra{} P1 PM 3영업일 \ra{} 스폰서 5영업일 \ra{} 이사회 임시 소집 10영업일, 발주사 응답 시한 5영업일 초과 시 대기 비용 산정 개시)은 구두로. 예상 소요 8분.}
\end{frame}

\begin{frame}{산출물 ④ 흔한 오류 점검}
\begin{itemize}
\item 회의체 이름만 적고 권한을 안 적음 / 정족수·반대의견 기록 방식이 없어 "전원 찬성"만 남음 / 사내 CCB와 변경협의체의 선후 관계를 안 적음
\item 경로는 있는데 시한이 없다 / "필요 시 보고"(필요는 항상 사후에 판정된다) / 수행사의 에스컬레이션 경로를 뺐다
\item 모든 프로젝트에 같은 ±10\% — 허용범위는 BC 편익 민감도(손익분기 66\%)에서 역산 / 작업패키지 계층을 비워 수행사에 재량이 없다
\item \textbf{A가 둘 이상인 행(공동 책임 = 무책임)} / PM이 모든 행의 A / 활동("주간 회의")을 행으로
\item 위원회와 개인이 섞인 열에서 위원회 안의 누가 A인지 불명확
\end{itemize}
\keymsg{점검 — 모든 행에 A가 정확히 한 명인가 / 확인하지 않은 A는 A가 아니다}
\note{워크북 §4.3. 10분 작성 후 점검. 점검 질문(§4.7): 모든 행에 A가 정확히 한 명인가 / PM이 A가 아닌 행의 A들이 자기가 A라는 사실을 서명으로 확인했는가 / 사내 회의체와 계약상 회의체가 같은 변경을 어떤 순서로 다루는지 한 문장으로. 예상 소요 3분 (설명) + 작성 시간.}
\end{frame}

\begin{frame}{산출물 ④ 예시본 참조 — 워크북 §4.5\til{}4.6 (26행 RACI)}
\begin{itemize}
\item \textbf{왜 P1 PM이 A인 행이 26개 중 8개뿐인가} — 착수·기획 결정 26개 중 18개의 최종 책임은 PM 밖에 있다(3PL 한동석, 가맹 배포 오정근, 아키텍처 박세연, 규제 해석 최영주, 기준선 이사회)
\item 왜 컨틴전시 집행의 A를 PM으로 두고 CCB를 I로
\item 왜 가맹 SW 배포 범위(A 오정근)에 CCB·정미란을 C로 — C를 빼면 배포 범위가 조용히 줄고 P1은 G2에서 발견
\item 왜 CCB와 변경협의체 순서를 문장으로 못박았는가
\item 왜 박세연의 이해 충돌(교체 대상 시스템의 설계자가 아키텍처 승인자)을 표에 기록했는가 — 정본 판정에서 그는 R이지 A가 아니고 감리 확인 병행
\end{itemize}
\keymsg{이 표를 스폰서에게 보여 주는 이유 — A인 사람들이 자기가 A라는 것을 알게 하기 위해서}
\note{5분 예시본 대조. PM 책임을 줄이려는 것이 아니다 — 2022 WMS 회의록의 비어 있는 서명란은 그가 A였다는 것을 아무도 문서로 말하지 않았기 때문이다. 완성 예시 PDF: 템플릿/PM\_Day1/04\_거버넌스RACI\_완성예시.pdf. 심화 과제(§4.7): 자기 회사 운영위 명단에 3역할 배정 \ra{} 배정 안 되는 인물 수와 빈 역할, 최근 결정 10개로 RACI \ra{} A가 둘이거나 없는 행 수. 예상 소요 4분.}
\end{frame}

\begin{frame}{산출물 ⑤ 프리모템 — 과제 지시 (조별로 실제로 실시한다)}
\begin{itemize}
\item \textbf{과제}: 조별로 프리모템을 실시하고 §5.4 템플릿의 항목 2(실패 선언문) · 3(원문) · 4(군집화) · 5(초기 리스크 3행) · 6(즉시 조치)
\item ① 진행자가 \textbf{실패 선언문}을 읽는다 — 과거형, 날짜·숫자 포함. F1 "2027-08, 컷오버 창을 놓쳐 IPO 예비심사 서류에 '구축 중'으로 기재되었다"
\item ② 각자 \textbf{2분간 개별 작성} — "무엇이 일어났다"의 과거형. "\til{}할 것 같다"는 다시 쓰게 한다
\item ③ 돌아가며 한 가지씩, \textbf{편집 없이} 번호를 붙여 기록. 중복도 그대로(중복 횟수가 우선순위 정보)
\item ④ 원인별로 군집화, 언급 수 + 1인 3표로 순위. 군집명은 범주("일정")가 아니라 \textbf{원인}("가맹 부속서 협상 지연")
\end{itemize}
\keymsg{브레인스토밍은 "요구사항 변경"을, 프리모템은 "협의회가 2026-08 부속서 서명을 거부했고 우리는 이미 가맹 SW를 개발 중이었다"를 낸다}
\note{워크북 §5.1\til{}5.3(Klein 2007). 실패 선언문 세 개: F1(일정) / F2 "2028-03-31 G5에서 편익 실현률 41\%로 판정되었다" / F3 "2026-11 G2에서 프로그램 이사회가 P1 중단을 의결했다". 왜 착수 리스크 식별의 첫 도구인가(§5.1): 착수 시점 편향(낙관, 계획 오류, 고유성)이 가장 강하고 그 절반은 정치적 편향이다. 분야별: SI는 발주사 내부·계약 경계가 대부분, 기술 리스크는 의외로 적음. 수행사 포함 여부는 논쟁적 — 발주사 세션과 합동 세션 분리가 절충 / EPC는 HAZOP·Constructability Review가 제도화 / NPD는 FMEA가 부품·공정 단위 프리모템. 조당 5\til{}6명, 진행자는 강사가 아니라 조원 중 한 명. 세션 정보(항목 1)는 참가자를 역할로 적어도 된다 — 익명성이 기법의 전제. 실패 선언문 세 개를 다 쓰면 좋지만 시간상 F1 또는 F2 하나로. 예상 소요 5분 (설명) + 실시 12분.}
\end{frame}

\begin{frame}{작성 요령 ⑥ — 리스크 서술 3부 · 소유자 · 즉시 조치 (워크북 §5.3 항목 5\til{}6)}
\begin{itemize}
\item \textbf{항목 5} — 군집을 리스크 서술로: \textbf{원인 \ra{} 사건 \ra{} 영향}. 원인은 \textbf{현재 사실}, 사건은 \textbf{불확실한 미래}, 영향은 \textbf{목표에의 결과}. \textbf{기회 1\til{}2건 포함}
\item R-01: "P2 벤더가 리드타임 34\ra{}41주 가능성을 통지한 상태다 \ra{} MC가 2027-02-12를 넘긴다 \ra{} 물류 통합시험을 G3 전에 못 하고 컷오버를 분리하거나 창을 놓친다"
\item 원인 \ra{} \textbf{트리거}, 사건 \ra{} \textbf{근접성}, 영향 \ra{} \textbf{대응 방향}. Day2 정량화는 확률을 사건에, 금액을 영향에
\item 소유자는 원인에 가장 가까이 있고 대응 자원을 가진 사람 — \textbf{PM이 아닌 경우가 많다}. 예시본 13개 위협 중 PM 소유는 3개
\item \textbf{항목 6 즉시 조치} — 리스크 대응 계획을 기다리지 않고 \textbf{지금 바꿔야 하는 것}: 헌장 가정·제외 범위 수정, 등록부 추가, RACI 행 추가. 각 줄에 "어느 문서의 어느 항목을" + 담당·기한
\end{itemize}
\keymsg{오류 — 원인과 사건을 섞음("P2 지연 리스크") / 이미 발생한 이슈를 리스크로 / 기회를 안 적음}
\note{워크북 §5.3 항목 5·6과 §5.6 해설. 항목 5의 열: ID / 원인\ra{}사건\ra{}영향 / 범주 / 출처(군집 번호, 헌장 가정 번호) / 정성 확률·영향(상중하) / 근접성 / 소유자(실명) / 초기 대응 방향 / 트리거 신호 / 다음 검토. PM 소유 3개는 R-07 요구사항 집중, R-13 컷오버 단말, R-04 대응안. 오류: 영향 없음 / 소유자 전부 PM. 항목 6 예시본: 헌장 가정에 "자사앱 벤더 연동 2026-12" 추가(원문 8번이 드러냄) / 등록부 S-19 자사앱 벤더를 위험 유형으로 / RACI 12행(앱 소스 인수) 추가 / BC 가정 A3(분모)에 R-12 연결, 편익 책임자 서명을 G1 통과 조건에 / 가맹 단말 기종 전수 조사(R-13) 착수 요청 — 오정근, 2026-04. 오류 3번(결과가 리스크 목록에만 남고 헌장은 그대로)의 처방이 항목 6이다 — "다섯 문서는 서로를 고친다". 예시본의 기회 O-02(가맹 포털을 2026-Q4 먼저 릴리스하면 부속서 협상의 지렛대 \ra{} R-02 확률 감소)가 프리모템에서 나온 가장 실용적 발견. 예상 소요 7분.}
\end{frame}

\begin{frame}{산출물 ⑤ 흔한 오류 점검}
\begin{itemize}
\item 실패 선언문을 조건문으로("만약 \til{}한다면") — 그러면 브레인스토밍 / 실패의 모습이 추상적("잘 안 됐다") / 한 종류의 실패만 제시
\item PM 혼자 쓴다(PM의 걱정 목록) / 임원이 참석해 발언 위축 / 운영·현업이 없어 기술 리스크만
\item 수집 단계에서 토론이 시작되어 사유가 걸러진다 / PM이 "그건 이미 대응 중"이라며 삭제 / 원문을 곧바로 범주로 뭉갬
\item 군집을 너무 크게("외부 요인") / 득표만 보고 언급 수 적은 치명적 사유(원 개발자 1명 이탈)를 탈락시킴
\item 리스크 목록이 일반론 / 소유자 전부 PM / 이미 발생한 이슈를 리스크로 / 결과를 헌장에 연결 안 함
\end{itemize}
\keymsg{점검 — 소유자 열에서 PM이 절반을 넘는가: 넘는다면 통제 밖 리스크를 PM이 떠안고 있거나 아직 못 찾은 것}
\note{워크북 §5.3, 교재 §3.5. 조별 리스크 3행을 발표시키고 "원인이 현재 사실인가"만 검사한다. 점검 질문(§5.7): 실패 선언문이 과거형인가 / 리스크 "원인"이 지금 확인 가능한 사실인가 / 소유자 열에서 PM이 절반을 넘는가. 예상 소요 3분.}
\end{frame}

\begin{frame}{산출물 ⑤ 예시본 참조 — 워크북 §5.5\til{}5.6}
\begin{itemize}
\item 예시본 구성: 세션 정보, 실패 선언문 F1\til{}F3, 원문 15개, 군집 12개, 리스크 R-01\til{}R-13 + O-01·O-02, 즉시 조치 7건
\item 왜 실패 선언문을 \textbf{세 개}로 — F1만 주면 사유가 전부 일정으로 쏠린다. F2(제때 끝나도 편익 실패), F3(통제 밖의 것이 결합해 중단) = Flyvbjerg의 세 성공은 별개
\item 왜 리스크 서술을 세 부분으로 강제했는가 / 왜 소유자에 PM이 세 번밖에 없는가(RACI 8/26과 같은 그림)
\item 왜 기회를 넣었는가 / 왜 "즉시 조치"에서 헌장을 고쳤는가
\end{itemize}
\keymsg{완성 예시 PDF: 템플릿/PM\_Day1/05\_프리모템리스크\_완성예시.pdf}
\note{4분 예시본 대조. §5.7 기본 과제(2026-02-16 내부통제 지침 — 5,000만 원 이상 CFO 직접 결재를 컨틴전시에도 적용 \ra{} R-14 서술, R-04·R-05와의 결합)는 사후 과제. 예상 소요 4분.}
\end{frame}

\begin{frame}{다섯 문서의 상호 관계 — 순서대로 쓰지만 순서대로 완성되지 않는다}
\begin{center}\usebeamercolor[fg]{normal text}
\begin{tikzpicture}[
  doc/.style={draw=mDarkTeal, rounded corners=2pt, fill=mDarkTeal!6, align=center,
              font=\scriptsize\bfseries, inner sep=3pt, minimum height=8mm, text width=36mm},
  lab/.style={font=\tiny, text=black!75, align=center, inner sep=1pt, fill=white, fill opacity=0.85, text opacity=1},
  arr/.style={-{Stealth[length=2mm]}, thick, draw=mDarkTeal}, x=1mm, y=1mm]
\node[doc] (bc)   at (0,0)    {비즈니스 케이스 요약};
\node[doc] (ch)   at (92,0)   {헌장 §2 목적 / §3 성공 기준 / §14 중단 조건};
\node[doc] (pm)   at (0,-21)  {프리모템·초기 리스크};
\node[doc] (sr)   at (92,-21) {이해관계자 등록부};
\node[doc] (ch2)  at (0,-42)  {헌장 §9·§10 수정 (v1.1)};
\node[doc] (raci) at (92,-42) {거버넌스·RACI};
\draw[arr] (bc) -- node[lab, above] {편익 118억 · 손익분기 66\% · 승인 조건 5} (ch);
\draw[arr] (ch) -- node[lab, left, text width=26mm] {범위 경계·제외 범위·가정·제약} (sr);
\draw[arr] (sr) -- node[lab, above, pos=0.55] {누가 경계 안팎에 있는가 / 가정 1$\sim$6 $\rightarrow$ 리스크 원인} (pm);
\draw[arr] (pm) -- node[lab, right, text width=24mm] {분모 가정 A3 $\leftarrow$ R-12} (bc);
\draw[arr] (pm) -- node[lab, right, text width=24mm] {소유자 실명, 즉시 조치} (ch2);
\draw[arr] (sr) -- node[lab, left, text width=26mm] {현저성 유형 $\rightarrow$ 의사결정체 배정, 관계 소유자} (raci);
\draw[arr] (raci) -- node[lab, below] {A가 누구인가 $\rightarrow$ 리스크 소유자, 에스컬레이션 경로} (ch2);
\end{tikzpicture}
\end{center}
\keymsg{따로 쓰면 각 문서는 다른 문서가 이미 의심한 것을 계속 참이라고 말한다}
\note{워크북 §6 그대로. BC \ra{} 헌장: 목적은 118억·규제 기한을 인용, 성공 기준의 편익 층은 편익 표를 그대로, 중단 조건 66\%는 재무 평가에서. 헌장 \ra{} 등록부: 가맹 단말 HW를 제외 범위에 넣는 순간 서정필의 우려가 등록부의 한 칸이 되고, 3PL 경계 조건을 적는 순간 대성로지스가 휴면 유형으로. 등록부 \ra{} RACI: 확정 유형 4명 = 3역할 + 수행사 PM = 기준선 동결의 필수 합의자. 위험 유형의 관계 소유자(한동석·박세연) = RACI 해당 행의 A. 헌장·등록부·RACI \ra{} 프리모템: 가정 6개가 R-01·02·03·05·08의 원인, RACI의 A가 리스크 소유자. 프리모템 \ra{} 앞 문서들의 수정: 원문 8번(앱 벤더) \ra{} 헌장 v1.1 가정 추가, 등록부 S-19, RACI 12행. 이것이 다섯 문서를 하루에 함께 쓰는 이유. 오늘 쓴 것은 G1(2026-04-30)에 v1.0\ra{}v1.1로 확정되는 초안이라고 다시 말한다. 예상 소요 4분.}
\end{frame}

\begin{frame}{6교시 정리 — 세 문서에서 반복된 규칙}
\begin{enumerate}
\item \textbf{사람은 부서명이 아니라 직책·실명}이며, 책임(관계 소유자, RACI의 A, 리스크 소유자)은 대부분 \textbf{PM 밖}에 있다 — 관계 소유자 10행, A 8/26, 리스크 소유자 3/13
\item \textbf{PM이 통제하지 못하는 것을 문서로 남기는 것}이 PM의 권한을 지키는 방법
\item 가장 위험한 이해관계자는 권력자가 아니라 \textbf{정당성 경로가 없는 권력(위험)}과 \textbf{연합할 동기가 있는 의존} — 대응은 정당성 경로 제공과 후원자 연결
\item 프리모템은 \textbf{"실패했다. 왜"}의 문법이고, 완료 조건은 리스크 목록이 아니라 \textbf{헌장·등록부·RACI의 수정}
\item Day2로: 헌장 §5·§6 \ra{} WBS / §4·§13 \ra{} RTM·MoSCoW / §7·§9 \ra{} 일정 기준선 / BC §5·헌장 §8 \ra{} 원가 기준선 / 리스크 13+2건·허용범위·66\% \ra{} 정량화
\end{enumerate}
\keymsg{"내가 통제하지 못하는 것이 내 일정을 결정한다는 사실을 문서로 남기지 않은 PM은 지연이 발생했을 때 설명할 근거가 없다"}
\note{워크북 §6 공통 규칙과 §7 "Day2로 이어지는 것" 표의 요약. RACI 19·20행이 원가 기준선의 입력이다. Day2에서 Day1 문서로 되돌아오는 것도 있다 — WBS를 만들다 헌장 범위에 없던 산출물(파일럿 12개점용 임시 연동)이 나타나면 헌장 v1.1 범위 수정 요청, EMV 합계가 8억을 넘으면 허용범위·예비비 재협상. 예상 소요 3분.}
\end{frame}

% =====================================================================
\section{마무리 (20분) — 학술 배경 깊이 읽기 · 읽을 자료 · Day2 예고}
% =====================================================================

\begin{frame}{이 블록에서 배우는 것}
\begin{itemize}
\item 오늘 인용한 학술 문헌 네 갈래를 \textbf{왜 그런가 / 언제 틀리는가}의 관점에서 다시 본다 [교재 제4장]
\item Rethinking Project Management의 다섯 방향 / Morris의 front-end / Flyvbjerg의 착각과 기만 / Staw의 몰입 상승
\item 읽을 자료 — 필수 · 30일 · 90일 · 참고 [제5장]
\item 오늘 만든 산출물 다섯 가지와 사후 읽기 과제
\item Day2 예고 — 계획: 범위·일정·원가 기준선
\end{itemize}
\keymsg{표준이 "이렇게 하라"고만 말하는 지점에서 학술 문헌은 "왜"와 "언제 틀리는가"를 말한다 — 그것을 아는 사람만이 테일러링할 수 있다}
\note{제4\til{}5장. 실무자에게 학술 문헌이 필요한 이유는 권위를 빌리기 위해서가 아니다 — 표준이 "이렇게 하라"고만 말하는 지점에서 학술 문헌은 "왜 그런가"와 "언제 틀리는가"를 말해 주며, 그것을 아는 사람만이 테일러링할 수 있다(제4장 서문). 시간 배분: §4.1\til{}4.4 각 3\til{}4분, 읽을 자료 4분, 산출물·과제·예고 6분. 예상 소요 2분.}
\end{frame}

\begin{frame}{Rethinking Project Management (2006) — 다섯 방향이 이 교재의 뼈대다}
\fig[0.63]{fig_4_1_rpm}
\figcap{그림 4-1. 다섯 방향 FROM \ra{} TOWARDS + Day1 대응 절 (Winter, Smith, Morris \& Cicmil, IJPM 24(8))}
\keymsg{표준 변경을 뉴스가 아니라 사조로 — 7판은 이 논문의 15년 늦은 반영, 8판은 학술과 실무의 재타협}
\note{§4.1. 2003년 EPSRC 연구 네트워크 \ra{} 2006 IJPM 24(8) 특집, 대표 논문 12페이지. 다섯 방향: ① 복잡성 이론(단순한 순차 라이프사이클 \ra{} 복잡성 이론들: §1.4 NTCP, §2.6 테일러링) ② 사회적 과정(도구적 이미지 \ra{} 사람들이 상호작용하며 만들어 가는 것: §1.3 임시조직, §3.3 이해관계자) ③ 가치 창출(산출물 인도 \ra{} 편익과 가치: §3.1) ④ 개념화(사전 정의된 목표를 가진 좁은 개념 \ra{} 다중 목표·다중 이해관계자·진행 중 정의가 바뀌는 넓은 개념: §1.1, §4.2 Morris) ⑤ 실무자 발달(trained technician \ra{} reflective practitioner: 이 과정 전체가 자격증 수업이 아닌 이유). APM BoK 7판(2019)이 그 사이. 흔한 오독: "애자일이 waterfall을 대체했다"는 서사로 축약 / 다섯 방향을 순차 성숙 단계로 / "성찰"을 회고 회의로 등치(Schön의 reflection-in-action은 실시간 판단 습관) / 자격증 무용론으로 — 원 논문은 대체가 아니라 발달 경로. 같은 특집호의 Cicmil et al.(actuality of projects)과 Crawford et al.(practitioner development)을 한 줄씩. 읽을 것: Winter et al.(2006) 전체, 특히 Table 1. 그다음 Geraldi \& Söderlund(2018). 예상 소요 4분.}
\end{frame}

\begin{frame}{Morris — "project management"에서 "the management of projects"로}
\begin{itemize}
\item 전자는 실행 통제(계획·감시·변경 통제), 후자는 프로젝트가 정의되는 \textbf{앞단(front-end)}, 기술, 상업, 정치까지 — 성패는 대부분 앞단에서 결정되는데 \textbf{BOK는 실행만 다룬다}
\item \textbf{Samset \& Volden(2016)의 역설}: \textbf{정보의 가치가 가장 큰 시점에 예산·권한·데이터가 가장 적다}
\item Flyvbjerg의 break-fix model — 앞단이 부실하면 실행에서 반복적 파손과 수리, 그 수리 비용이 원가 초과의 실체
\item 한국 조직에서 PM은 품의서 결재 뒤에 지명된다 — 그래도 둘은 할 수 있다: ① 앞단의 결정을 목록화해 가정·리스크를 헌장에 ② 종료 보고서에 "앞단의 X가 실행 중 Y의 원인"을 기록
\item 오독: front-end를 "기획 문서 강화"로 — 요점은 문서량이 아니라 의사결정 권한과 \textbf{대안 유지} / "착수 지연"으로 — 처방은 think slow, \textbf{act fast}
\end{itemize}
\keymsg{"착수 시점에 잘못 적힌 것은 계획 시점에 고쳐지지 않는다" — 오늘 헌장의 가정 목록을 길게 쓰라고 한 이유}
\note{§4.2. Morris(1994, 『The Management of Projects』)가 바꾼 어법 하나. Morris(2013, 『Reconstructing Project Management』)는 표준 지식체계가 왜 front-end를 배제하게 되었는가의 계보. Samset \& Volden(2016)의 열 가지 역설 중 핵심: 착수 전 결정(무엇을, 어느 대안, 어느 규모)은 실행 중 어떤 최적화로도 되돌릴 수 없는데, 그 시점에는 정보도 자원도 없다. "실행을 잘하자"에는 한계선이 있고 그 선은 앞단에서 그어진다. ①은 Engwall, ②는 다음 프로젝트에서 앞단에 들어갈 근거를 만드는 일. 교재 맺음말 인용. 읽을 것: Morris(2013) Part 1·3, Samset \& Volden(2016) 역설 목록 먼저. 예상 소요 4분.}
\end{frame}

\begin{frame}{Flyvbjerg — 착각(delusion)과 기만(deception), 그리고 think slow, act fast}
\begin{itemize}
\item Flyvbjerg, Garbuio \& Lovallo(2009): \textbf{Delusion} — 낙관 편향·계획 오류 \ra{} 처방은 outside view·RCF / \textbf{Deception} — 전략적 왜곡 \ra{} 처방은 \textbf{인센티브·책임·리스크 자본 구조의 설계}
\item 2021 PMJ 「Top Ten Behavioral Biases」 1위는 \textbf{strategic misrepresentation} — 편향 교육으로 정치 문제를 풀려는 헛수고
\item 「Survival of the Unfittest」 — 선정되는 것은 \textbf{서류상 가장 좋아 보이는 프로젝트} = 원가 최소·편익 최대로 추정한 프로젝트
\item 온담 물류 96억의 분모 부재 — 착각인가 기만인가? \textbf{알 수 없고 알 필요도 없다. 처방이 같다} — 분모·측정 방법·오너 실명 없는 항목은 합계에서 제외
\item \textbf{Think slow, act fast}(2023) — 학습곡선은 \textbf{반복 가능한 동일 단위}에서: 매장 배포는 모듈화(파일럿 12개점\ra{}확대), 전환 7.4TB·컷오버 34시간은 bespoke \ra{} 리스크 집중
\end{itemize}
\keymsg{CFO의 "근거 없는 항목은 승인 안 함"이 정확히 Flyvbjerg의 제도 처방이다}
\note{§4.3. 2002 「Error or Lie?」 \ra{} 2009 「Survival of the Unfittest」. "behavioral bias = cognitive bias"는 오류, 절반은 political bias. 반대 오용: 모든 낙관 견적을 거짓말로 규정해 신뢰 파괴. 실행이 길수록 예상치 못한 사건의 창이 넓어진다. 반대편도 알아야 한다: Hirschman의 Hiding Hand(1967) — Flyvbjerg(2016)가 p<0.0001로 기각(방법론 비판: 11개 성공 사례만 골라 표본 — sampling on the dependent variable) \ra{} Ika(2018) 반박 \ra{} Kreiner(2020) "쟁점은 데이터가 아니라 존재론". "Flyvbjerg에 따르면"이라고 말할 때 그 반대편이 무엇인지도 말할 수 있는 실무자. 읽을 것: 입문 『프로젝트 설계자』 전체, 숫자는 2014 PMJ(arXiv 1409.0003)·2021 PMJ(arXiv 2202.00125) 무료, IT 청중 배포용은 2011 HBR 3페이지(arXiv 1304.0265). 예상 소요 5분.}
\end{frame}

\begin{frame}{Staw — 몰입의 상승(escalation of commitment)과 결정자의 분리}
\begin{columns}[T]
\begin{column}{0.55\textwidth}
\fig[0.56]{fig_4_4_staw}
\figcap{그림 4-2. 몰입의 상승 순환과 두 가지 조직 설계 처방}
\end{column}
\begin{column}{0.43\textwidth}
\begin{itemize}
\item Staw(1976): \textbf{같은 나쁜 결과에서, 그 결정을 자기가 내렸을 때 추가 투입이 유의하게 커진다} — 원인은 \textbf{자기 정당화}
\item 처방 ① \textbf{착수 결정자와 중단 결정자를 분리} — 김선호 대표가 의장·게이트 승인자·비전 소유자. 사외이사 1인이 최소 장치
\item 처방 ② \textbf{중단 기준(kill criteria)을 착수 시점에 사전 확약} — 이사회 조건 5항의 핵심
\end{itemize}
\end{column}
\end{columns}
\keymsg{escalation은 매몰비용 계산 착오가 아니라 사회적·평판적 현상 — 게이트마다 작동한다}
\note{§4.4. 실험: R\&D 예산을 두 사업부 중 하나에 배분 \ra{} 나쁜 결과 \ra{} 다시 배분. 조건은 첫 배분을 자기가 했는가 vs 남이 한 것을 이어받았는가. 원인이 자기 정당화라면 처방은 회계 교육("매몰비용을 무시하세요")이 아니라 조직 설계다. 온담: G1·G2에서 중단해야 할 상황이 오면 그 결정을 대표가 내릴 수 있겠는가. kill criteria 예: "G2에서 정본이 확정되지 않으면 컷오버를 2027 하반기로", "G1에서 편익 프로필 8종 중 5종 이상의 오너·분모가 없으면 통과시키지 않는다" — 정보가 없어 자기 정당화가 작동하지 않는 착수 시점에 정한다. 2019 스마트오더 감사 지적 "중단 시 손실 확정 절차 없음". 오독: escalation = sunk cost fallacy(매몰비용 오류는 개인의 계산 오류) / 항상 비합리로 규정 — 정보가 갱신되어 계속하는 것이 옳은 경우와 구별 / Flyvbjerg 2021 편향 목록 \#10 — 착수만이 아니라 게이트마다. 4교시 헌장 항목 14와 5교시 NPD "kill 0건이면 고장"의 학술 근거. 생각해 볼 질문 5. 읽을 것: Staw(1976) Method 절, Flyvbjerg(2021) \#10. 예상 소요 4분.}
\end{frame}

\begin{frame}{읽을 자료 — 필수 · 30일 · 90일 · 참고 (제5장)}
\subhead{필수 (과정 전, 총 90분) — 사례 회사 문서 §1\til{}3·§5 / Klein(2007) 2페이지 / 교재 §2.1 연표·§2.2 표 / Cooke-Davies(2002) 6페이지}
\begin{tbl}[\scriptsize]{L{0.42\textwidth}L{0.5\textwidth}}
\textbf{30일 안에} & \textbf{왜} \\ \midrule
Flyvbjerg \& Gardner, 『프로젝트 설계자』(2024) & "우리 회사 base rate은 어디서 구하나"로 남은 질문. 8시간 \\
Winter et al.(2006), IJPM 24(8) & 다섯 방향 — 표준 변경을 사조로 \\
Lundin \& Söderholm(1995) / Engwall(2003) & 4T·4행위 / 착수의 진짜 일이 "이미 결정된 것의 목록화"인 이유 \\
Zwikael \& Smyrk(2011) ITO 장 & 워크북 편익 프로필 "오너 실명" 칸의 이론 원전 \\
Eveleens \& Verhoef(2010) 7페이지 & 제안서에서 CHAOS 수치를 빼야 하는 이유를 데이터로 \\
Mitchell, Agle \& Wood(1997) 7유형 정의 절 & 워크북 §3 분류 근거 \\
\end{tbl}
\keymsg{읽지 않아도 되는 것 — CHAOS 원문 / Stacey matrix 2차 자료 / "PMBOK 7판 요약" 류 시중 교재(현행은 8판)}
\note{§4.5와 제5장. 90일 안에(자기 조직에 적용할 때): Shenhar \& Dvir(2007) NTCP / Boehm \& Turner(2003) 5축 / Ward \& Daniel(2012) BDN / Bradley(2010) / PRINCE2 7(Organizing·Progress·Business case practice) / PMBOK 8판(Standard 전체, Governance·Finance) / Morris(2013) Part 1·3 / Flyvbjerg(2021) 편향 10. 참고(워크북 쓰다가 손을 뻗을 것): ISO 21502:2020 §6.4·§7.12·§5.2 / ISO 21505:2017 / ISO 21506:2024 / PMBOK 6판(사내 조항 번호 원문) / Agile Practice Guide Suitability Filter / Snowden \& Boone(2007) / Mitchell, Russo \& Pennington(1989) / Staw(1976) / P3.express v2.0 / PM² Guide v3.1 / 소프트웨어 진흥법 §50. 읽지 않아도 되는 것의 대체: CHAOS 원문 \ra{} Eveleens \& Verhoef, Stacey 2차 자료 \ra{} Snowden \& Boone. 서지는 교재 §4.5·§5.1\til{}5.3과 05\_리딩리스트/00\_리딩리스트.md의 검증 표기를 따른다. arXiv 번호가 있는 것은 무료. 한국어판 표기가 없는 것은 확인되지 않은 것이며 없다는 증명이 아니다. Kahneman \& Tversky 1979 수록지 등 일부 서지는 2차 출처 확인이므로 논문 인용 시 원문 대조. 예상 소요 4분.}
\end{frame}

\begin{frame}{오늘 만든 산출물 다섯 가지}
\begin{tbl}[\scriptsize]{cL{0.25\textwidth}cL{0.28\textwidth}L{0.26\textwidth}}
\textbf{\#} & \textbf{산출물} & \textbf{워크북} & \textbf{PRINCE2 7 대응} & \textbf{상태} \\ \midrule
① & 비즈니스 케이스 요약 & §1 & Outline Business Case (Project Brief 내) & v0.9 — G1(2026-04-30)에 v1.0 확정 \\
② & 프로젝트 헌장 & §2 & Project Brief \ra{} PID(핵심부) & v1.0 — 프리모템 반영 후 v1.1 \\
③ & 이해관계자 등록부 + 현저성 분석 & §3 & Communication management approach의 분석부 & G0 판정 — G1·2026-08·G2·G3 갱신 \\
④ & 거버넌스 구조 + RACI 초안 & §4 & PM team structure, Role descriptions & 허용범위 3층은 Day2 확정 \\
⑤ & 프리모템 기록 + 초기 리스크 목록 & §5 & Risk register(초기), Risk management approach(초안) & 13+2건 정성 — Day2에서 확률·영향·EMV \\
\end{tbl}
\begin{itemize}
\item 다섯 문서 모두 \textbf{"결정되지 않은 것을 결정하는 문서"} — Day2 이후의 문서는 Day1에서 결정된 것을 분해하고 측정한다
\end{itemize}
\keymsg{헌장의 한 줄이 모호하면 WBS 열 개가 모호해진다 — 완성하지 못한 항목은 예시본과 대조해 채우고 내일 가져오라}
\note{워크북 §0.2 표의 Day1 부분. 오늘 완성하지 못한 항목은 예시본(§1.5·§2.5·§3.5·§4.5·§5.5, 통합본 템플릿/PM\_Day1/00\_Day1\_템플릿팩\_완성예시.pdf)과 대조하며 사후에 채운다. 워크북 §0.1 사용법 10번 — "Day1 산출물은 Day2 계획 산출물의 입력이다. 오늘 쓴 것을 버리지 말고 내일 가져오라." 예상 소요 3분.}
\end{frame}

\begin{frame}{사후 읽기 과제 (Day2까지)}
\begin{columns}[T]
\begin{column}{0.48\textwidth}
\subhead{읽기 (총 1시간 안)}
\begin{itemize}
\item Klein(2007) 「Performing a Project Premortem」 2페이지 — 왜 "실패했다"인지 원문에서
\item Cooke-Davies(2002) 6페이지 — 헌장 항목 3의 두 층이 왜 분리되어야 하는지
\item 교재 §3.2 "항목별 의미"와 §3.4 "권한 없는 PM" — 오늘 쓴 헌장 항목 12를 다시 읽으며
\end{itemize}
\end{column}
\begin{column}{0.5\textwidth}
\subhead{쓰기 (심화 과제, 3\til{}7년차 — 하나만)}
\begin{itemize}
\item §1.7 최근 승인 문서의 편익 표를 항목 4 형식으로 — 분모 없는 항목·부서명 책임자·남의 편익의 개수
\item §2.7 자기 회사 헌장에 \textbf{항목 12와 14}가 없다면 그 둘만 써서 스폰서에게 보여 줄 초안
\item §4.7 운영위 명단에 3역할 배정 — 배정 안 되는 인물 수와 빈 역할
\item §5.7 실패 선언문 세 개(일정·편익·중단), 각 5개 사유 — 현재 등록부에 없는 것의 개수
\end{itemize}
\end{column}
\end{columns}
\keymsg{생각해 볼 질문 — 제1장 3 / 제2장 3 / 제3장 1 중 하나씩 골라 답을 적어 오라. Day2 첫 10분에 두세 명이 말한다}
\note{읽기 과제는 총 1시간 안. 심화 과제는 네 개 중 하나만 해도 된다 — §2.7은 Day2 시작 때 한 문장씩 공유. 기본 과제 변형(시간이 남는 조): §1.7 CFO 148억 확정·P1 2026년 집행 102\ra{}86억 시 항목 5·6 재작성(이월 vs 총예산 152억 축소, 후자면 do-minimum이 되살아나는가) / §2.7 대표의 "컷오버를 2026-11로" 요구 시 항목 13·9·10 재작성과 스폰서 제시 대안 세 개. 생각해 볼 질문(교재 각 장 끝): 제1장 3 — 온담 2019 스마트오더는 예산 안에서 중단됐다. 이것은 실패인가, 어느 정의로 실패인가 / 제2장 3 — 여러분 프로젝트에서 되돌릴 수 없는 결정은 무엇이고, 그 앞에 얼마나 시간을 쓰고 있는가 / 제3장 1 — 여러분의 마지막 프로젝트에서 편익 오너는 누구였는가, 그 사람은 프로젝트가 끝난 뒤에도 그 조직에 있었는가. 답은 Day2 첫 10분에 두세 명이 말한다. 예상 소요 2분.}
\end{frame}

\begin{frame}{Day2 예고 — 계획: 범위 · 일정 · 원가 기준선}
\begin{tbl}[\scriptsize]{L{0.28\textwidth}L{0.64\textwidth}}
\textbf{Day2 산출물} & \textbf{Day1에서 받는 입력} \\ \midrule
⑥ 범위 기술서와 WBS & 헌장 §5 범위, §6 산출물 — 포함 5개 산출물이 WBS 레벨 2. \textbf{PBS 노드는 명사} \\
⑦ 요구사항 추적표(RTM) 골격 & 헌장 §4 상위 요구사항(요구 제공자), §13 우선순위 \ra{} MoSCoW \\
⑧ 일정 기준선 & 헌장 §7 마일스톤(외부 고정), §9 제약 \ra{} 컷오버·규제에서 역산한 주공정, G3 합류 버퍼 15영업일 \\
⑨ 원가 기준선과 예비비 구조 & BC §5·3층 예산, 헌장 §8, RACI 19·20행 \ra{} 145.8 / 컨틴전시 10.2 / 관리예비비 12.0. R-04 시나리오별 축소안 \\
⑩ 리스크 등록부 정량화 + 허용범위 격자 & 초기 리스크 13+2건, 허용범위 초안, 손익분기 66\% \ra{} 확률·영향·EMV, EMV 합계 vs 8억, 편익 $-$5\%p의 타당성 역산 \\
\end{tbl}
\begin{itemize}
\item 다시 만나는 Day1 개념: \textbf{RCF} / \textbf{멱함수 분포} / \textbf{되돌림 비용} / \textbf{Product-Based Planning} / \textbf{허용범위 격자}
\end{itemize}
\keymsg{Day3 실행·통제(3층 변경 통제의 비용, Kreiner의 표류 역설) · Day4 종료·이관(판본이 바뀌어도 살아남는 산출물, ITIL)}
\note{워크북 §7 표. RCF는 원가 기준선의 base rate, 멱함수 분포는 컨틴전시가 꼬리를 못 덮는 이유(처방은 modularity), 되돌림 비용은 G2·G3 앞의 예측형 요소, PBP는 활동 WBS vs 산출물 PBS, 허용범위 격자는 EMV 8억·원가 +6\%. Day2에서 Day1 문서로 되돌아오는 것(WBS에서 발견되는 헌장에 없던 산출물 \ra{} v1.1 범위 수정 요청 / EMV 합계 초과 \ra{} 허용범위·예비비 재협상)을 말하고 마친다. 교재 맺음말 — "착수 시점에 잘못 적힌 것은 계획 시점에 고쳐지지 않는다. 그것이 Morris가 말한 front-end의 무게이며, 여러분이 오늘 헌장의 가정 목록을 길게 쓰는 이유다." 예상 소요 4분.}
\end{frame}
